\paragraph{Peter--Weyl 因子化与覆盖 \texorpdfstring{$\zeta$}{zeta} 极点秩序的维数加权律}
在正规覆盖的 Peter--Weyl 块对角化下，覆盖行列式按不可约通道作 Artin 因子化。由此，任一给定特征值对覆盖 \(\zeta\) 主点的贡献都可逐通道累加，并被严格记录为“表示维数 \(\times\) 通道重数”的和。

\begin{theorem}[正则覆盖 \texorpdfstring{$\zeta$}{zeta} 的极点阶数等于不可约通道维数加权的特征值重数和]\label{thm:xi-time-part64-cover-zeta-pole-order-weighted-channel-multiplicity}
沿用命题 \ref{prop:kernel-artin-factorization} 的记号。
设 \(G\) 为有限群，\(M_{K,G}(u)\) 为群扩张邻接算子，并记
\[
\det(I-zM_{K,G}(u))
=
\prod_{\rho\in\widehat G}
\det(I-zM_{K,\rho}(u))^{d_\rho},
\qquad
d_\rho:=\dim\rho.
\]
取任意复数 \(\Lambda\neq 0\)，并令 \(m_\rho(\Lambda)\) 表示 \(\Lambda\) 作为 \(M_{K,\rho}(u)\) 特征值的代数重数。
则覆盖 \(\zeta\) 函数
\[
\zeta_{K,G}(z,u):=\det(I-zM_{K,G}(u))^{-1}
\]
在 \(z=\Lambda^{-1}\) 处的极点阶数恰为
\[
\boxed{\;
\operatorname{ord}_{z=\Lambda^{-1}}\zeta_{K,G}(z,u)
=
\sum_{\rho\in\widehat G} d_\rho\,m_\rho(\Lambda)\; }.
\]
\end{theorem}
\begin{proof}
对每个 \(\rho\in\widehat G\)，在 \(z=\Lambda^{-1}\) 的邻域内可写
\[
\det(I-zM_{K,\rho}(u))
=(1-z\Lambda)^{m_\rho(\Lambda)}h_\rho(z),
\qquad
h_\rho(\Lambda^{-1})\neq 0.
\]
将其代入命题 \ref{prop:kernel-artin-factorization} 的 Artin 因子化，得到
\[
\det(I-zM_{K,G}(u))
=
\prod_{\rho\in\widehat G}
\bigl((1-z\Lambda)^{m_\rho(\Lambda)}h_\rho(z)\bigr)^{d_\rho}
\]
\[
=(1-z\Lambda)^{\sum_{\rho}d_\rho m_\rho(\Lambda)}\,H(z),
\qquad
H(\Lambda^{-1})\neq 0.
\]
取倒数即得
\[
\zeta_{K,G}(z,u)
=(1-z\Lambda)^{-\sum_{\rho}d_\rho m_\rho(\Lambda)}H(z)^{-1},
\]
故其在 \(z=\Lambda^{-1}\) 处的极点阶数恰为
\(
\sum_{\rho}d_\rho m_\rho(\Lambda)
\)。
证毕。
\end{proof}

\begin{corollary}[主极点的覆盖稳定性判据]\label{cor:xi-time-part64-cover-zeta-leading-pole-stability}
\leanverified{paper\_xi\_time\_part64\_cover\_zeta\_leading\_pole\_stability}
沿用定理 \ref{thm:xi-time-part64-cover-zeta-pole-order-weighted-channel-multiplicity} 的记号，并令
\[
\lambda_1(u)
\]
表示平凡通道 \(M_{K,\mathbf 1}(u)\) 的 Perron 根。
若满足：
\begin{enumerate}
  \item \(\lambda_1(u)\) 在平凡通道中是简单特征值；
  \item 对所有非平凡 \(\rho\neq \mathbf 1\)，都有
  \[
  \operatorname{spr}\!\bigl(M_{K,\rho}(u)\bigr)<\lambda_1(u),
  \]
  其中 \(\operatorname{spr}\) 表示谱半径，
\end{enumerate}
则覆盖 \(\zeta\) 在主点 \(z=\lambda_1(u)^{-1}\) 处满足
\[
\boxed{\;
\operatorname{ord}_{z=\lambda_1(u)^{-1}}\zeta_{K,G}(z,u)=1\; }.
\]
亦即有限正则覆盖在主谱层面不改变极点阶数，只可能修改次主通道。

更一般地，若恰有一族非平凡通道
\[
\mathcal R\subseteq \widehat G\setminus\{\mathbf 1\}
\]
满足
\(
\lambda_1(u)\in\Spec(M_{K,\rho}(u))
\)
对每个 \(\rho\in\mathcal R\) 成立，而对其余非平凡通道 \(\eta\notin\mathcal R\) 有
\(
\lambda_1(u)\notin\Spec(M_{K,\eta}(u))
\)，
则
\[
\operatorname{ord}_{z=\lambda_1(u)^{-1}}\zeta_{K,G}(z,u)
=
1+\sum_{\rho\in\mathcal R}d_\rho\,m_\rho(\lambda_1(u)).
\]
\end{corollary}
\begin{proof}
在平凡通道 \(\mathbf 1\) 上，\(\lambda_1(u)\) 的代数重数为 \(1\)。
若对所有非平凡 \(\rho\neq\mathbf 1\) 都有
\(
\operatorname{spr}(M_{K,\rho}(u))<\lambda_1(u)
\)，
则这些通道根本不可能含有特征值 \(\lambda_1(u)\)，故
\[
m_\rho(\lambda_1(u))=0
\qquad(\rho\neq\mathbf 1).
\]
由定理 \ref{thm:xi-time-part64-cover-zeta-pole-order-weighted-channel-multiplicity}，
\[
\operatorname{ord}_{z=\lambda_1(u)^{-1}}\zeta_{K,G}(z,u)
=
d_{\mathbf 1}\,m_{\mathbf 1}(\lambda_1(u))
=1.
\]

一般情形下，仍由同一定理，
\[
\operatorname{ord}_{z=\lambda_1(u)^{-1}}\zeta_{K,G}(z,u)
=
d_{\mathbf 1}\,m_{\mathbf 1}(\lambda_1(u))
+\sum_{\rho\in\mathcal R}d_\rho\,m_\rho(\lambda_1(u)),
\]
而 \(d_{\mathbf 1}=1\)、\(m_{\mathbf 1}(\lambda_1(u))=1\)，遂得所示公式。证毕。
\end{proof}

\noindent
于是 Peter--Weyl/Artin 因子化不仅把覆盖 \(\zeta\) 分解为有限多个不可约通道；它还把主极点的升阶机制完全刚性化：覆盖只能按表示维数与通道共振重数的乘积累加极点阶数，而在非平凡通道不触及主谱时，主极点保持简单。

\paragraph{非阿贝尔通道能量账本与覆盖 Perron 主相的平凡通道刚性}
在非阿贝尔 Chebotarev--Plancherel 能量恒等式、群扩张邻接算子的 Peter--Weyl 块对角化以及覆盖 \(\zeta\) 的 Artin 因子化已经建立之后，还可进一步抽取出下列两条严格后果：商映射在类函数层面损失的 \(L^2\) 能量可逐不可约通道精确记账，而正规覆盖算子的 Perron 根则始终由平凡通道完全控制。

\begin{theorem}[非阿贝尔商映射的精确能量损失公式]\label{thm:xi-time-part64-nonabelian-quotient-energy-loss-exact}
\leanverified{paper\_xi\_time\_part64\_nonabelian\_quotient\_energy\_loss\_exact}
沿用命题 \ref{prop:kernel-chebotarev-plancherel-nonabelian} 的记号。设
\[
\varphi:G\twoheadrightarrow H
\]
为有限群满同态，并记
\[
K:=\ker\varphi.
\]
令 \(\varepsilon\in \mathrm{Cl}(G)\) 为零均值类函数，即
\[
\sum_{g\in G}\varepsilon(g)=0.
\]
对每个 \(\rho\in\widehat G\) 定义不可约通道幅度
\[
B_G(\rho):=\sum_{g\in G}\varepsilon(g)\,\overline{\chi_\rho(g)},
\]
并定义纤维平均推前
\[
(\varphi_\ast\varepsilon)(h):=\frac1{|K|}\sum_{g\in\varphi^{-1}(h)}\varepsilon(g)
\qquad(h\in H).
\]
再记
\[
\mathrm{Inf}_\varphi(\widehat H)
:=
\{\widetilde \sigma:\ \sigma\in \widehat H,\ \widetilde \sigma=\sigma\circ\varphi\}
\subseteq \widehat G.
\]
则在归一化均匀测度对应的 \(L^2\) 范数下有严格恒等式
\[
\boxed{\;
\|\varepsilon\|_{L^2(G)}^2-\|\varphi_\ast\varepsilon\|_{L^2(H)}^2
=
\frac1{|G|^2}
\sum_{\rho\in\widehat G\setminus \mathrm{Inf}_\varphi(\widehat H)}
\bigl|B_G(\rho)\bigr|^2.\;}
\]
\end{theorem}
\begin{proof}
先由
\[
\sum_{h\in H}(\varphi_\ast\varepsilon)(h)
=
\frac1{|K|}\sum_{g\in G}\varepsilon(g)=0
\]
可知 \(\varphi_\ast\varepsilon\) 也是 \(H\) 上的零均值类函数。

对 \(G\) 应用命题 \ref{prop:kernel-chebotarev-plancherel-nonabelian}，得到
\[
\|\varepsilon\|_{L^2(G)}^2
=
\frac1{|G|^2}
\sum_{\rho\in\widehat G\setminus\{\mathbf 1_G\}}
\bigl|B_G(\rho)\bigr|^2.
\]

对 \(H\) 上的 \(\varphi_\ast\varepsilon\) 同样应用该命题。若 \(\sigma\in \widehat H\)，定义
\[
B_H(\sigma):=\sum_{h\in H}(\varphi_\ast\varepsilon)(h)\,\overline{\chi_\sigma(h)}.
\]
则由纤维平均定义，
\[
\begin{aligned}
B_H(\sigma)
&=
\sum_{h\in H}
\frac1{|K|}
\sum_{g\in\varphi^{-1}(h)}
\varepsilon(g)\,\overline{\chi_\sigma(h)}\\
&=
\frac1{|K|}
\sum_{g\in G}
\varepsilon(g)\,\overline{\chi_\sigma(\varphi(g))}\\
&=
\frac1{|K|}
\sum_{g\in G}
\varepsilon(g)\,\overline{\chi_{\widetilde \sigma}(g)}\\
&=
\frac1{|K|}B_G(\widetilde \sigma).
\end{aligned}
\]
由 \(|G|=|H||K|\) 遂得
\[
\frac1{|H|^2}\bigl|B_H(\sigma)\bigr|^2
=
\frac1{|G|^2}\bigl|B_G(\widetilde \sigma)\bigr|^2.
\]
因此
\[
\|\varphi_\ast\varepsilon\|_{L^2(H)}^2
=
\frac1{|G|^2}
\sum_{\rho\in \mathrm{Inf}_\varphi(\widehat H)\setminus\{\mathbf 1_G\}}
\bigl|B_G(\rho)\bigr|^2.
\]
从 \(\|\varepsilon\|_{L^2(G)}^2\) 的通道展开中减去上式，即得所示恒等式。证毕。
\end{proof}

\begin{theorem}[正规覆盖算子的 Perron 根完全由平凡通道控制]\label{thm:xi-time-part64-cover-perron-root-trivial-channel}
\leanverified{paper\_xi\_time\_part64\_cover\_perron\_root\_trivial\_channel}
沿用命题 \ref{prop:kernel-peter-weyl-block-diagonalization}、命题 \ref{prop:kernel-artin-factorization} 以及定义 \ref{def:kernel-rh} 的记号，并假设 \(M_K(u)\) 原始。记
\[
\lambda_1(u):=\rho\!\bigl(M_K(u)\bigr)=\rho\!\bigl(M_{K,\mathbf 1}(u)\bigr).
\]
则对每个不可约表示 \(\pi\in\widehat G\)，都有
\[
\boxed{\;
\rho\!\bigl(M_{K,\pi}(u)\bigr)\le \lambda_1(u).\;}
\]
并且整个覆盖算子的谱半径满足
\[
\boxed{\;
\rho\!\bigl(M_{K,G}(u)\bigr)=\lambda_1(u).\;}
\]
进一步，若记
\[
\eta(u):=\max_{\pi\neq \mathbf 1}\rho\!\bigl(M_{K,\pi}(u)\bigr),
\qquad
\Lambda(u):=\max\{|\lambda_2(u)|,\eta(u)\},
\]
则下列两条等价：
\begin{enumerate}
  \item 平凡通道的次特征值与全部非平凡通道谱半径共同满足平方根界，即
  \[
  \Lambda(u)\le \lambda_1(u)^{1/2};
  \]
  \item \(\textsf{RH}_K(u)\) 与 \(\textsf{GRH}_{K,G}(u)\) 同时成立。
\end{enumerate}
\end{theorem}
\begin{proof}
由命题 \ref{prop:kernel-peter-weyl-block-diagonalization}，
\[
M_{K,G}(u)
\ \cong\
\bigoplus_{\pi\in\widehat G}
\bigl(M_{K,\pi}(u)\otimes I_{\dim\pi}\bigr),
\]
故
\[
\rho\!\bigl(M_{K,G}(u)\bigr)
=
\max_{\pi\in\widehat G}\rho\!\bigl(M_{K,\pi}(u)\bigr).
\]
因此只需证明对每个 \(\pi\) 都有
\(
\rho(M_{K,\pi}(u))\le \lambda_1(u)
\)。

固定 \(\pi\in\widehat G\)。由于 \(G\) 有限，\(\pi\) 可酉化；取 \(V_\pi\) 上一个 \(G\)-不变 Hermitian 范数，使 \(\pi(g)\) 全为酉算子。设
\[
x=(x_i)_{i\in V}\in V_\pi^{\,V}
\]
是 \(M_{K,\pi}(u)\) 的特征向量，对应特征值 \(\mu\)。按定义
\[
(M_{K,\pi}(u)x)_i
=
\sum_{e:i\to j}u^{\kappa(e)}\,\pi(g(e))\,x_j.
\]
逐坐标取范数，并记
\[
h_i:=\|x_i\|\ge 0,
\qquad
h:=(h_i)_{i\in V}\neq 0,
\]
则由三角不等式与酉不变性，
\[
\begin{aligned}
|\mu|\,h_i
&=
\bigl\|(M_{K,\pi}(u)x)_i\bigr\|\\
&\le
\sum_{e:i\to j}u^{\kappa(e)}\,\|x_j\|\\
&=
(M_K(u)h)_i.
\end{aligned}
\]
故
\[
|\mu|\,h\le M_K(u)h
\]
在非负锥上逐坐标成立。由 Perron--Frobenius 的 Collatz--Wielandt 公式，
\[
|\mu|\le \rho(M_K(u))=\lambda_1(u).
\]
对 \(M_{K,\pi}(u)\) 的所有特征值取上确界，便得
\[
\rho\!\bigl(M_{K,\pi}(u)\bigr)\le \lambda_1(u).
\]

另一方面，平凡通道 \(\pi=\mathbf 1\) 上恰有
\[
\rho\!\bigl(M_{K,\mathbf 1}(u)\bigr)=\lambda_1(u),
\]
因此
\[
\rho\!\bigl(M_{K,G}(u)\bigr)=\lambda_1(u).
\]

最后，由定义
\[
\Lambda(u)=\max\{|\lambda_2(u)|,\eta(u)\},
\]
第 \(1\) 条本身即是
\[
\Lambda(u)\le \lambda_1(u)^{1/2},
\]
而这正等价于 \(\textsf{RH}_K(u)\) 与 \(\textsf{GRH}_{K,G}(u)\) 同时成立。证毕。
\end{proof}

\endinput

\paragraph{Chebotarev 商等号、最小承载商与极值见证集}
在阿贝尔与非阿贝尔两侧的 Chebotarev--Plancherel 能量恒等式、商单调性以及二主项见证定理都已建立之后，还可把“偏差在商下何时无损”“全部可见偏差究竟由哪个最小商承载”以及“最优见证集如何从极大通道中被唯一选出”进一步压缩为下列四条接口。

\begin{theorem}[阿贝尔 Chebotarev 偏差在商映射下的等号判据]\label{thm:xi-time-part64aa-abelian-quotient-equality-criterion}
\leanverified{paper\_xi\_time\_part64aa\_abelian\_quotient\_equality\_criterion}
沿用推论 \ref{cor:kernel-chebotarev-L2-monotone} 的记号。设
\[
\pi:G_2\twoheadrightarrow G_1
\]
为有限阿贝尔群满同态，并记
\[
\pi^\ast:\widehat{G_1}\hookrightarrow \widehat{G_2},
\qquad
\chi\longmapsto \chi\circ\pi.
\]
对固定 \(u\in(0,1]\) 与 \(n\ge 1\)，设细层与粗层的相对 Haar 密度分别为
\[
r_n^{(2)}(\,\cdot\,;u)\in \CC^{G_2},
\qquad
r_n^{(1)}(\,\cdot\,;u)\in \CC^{G_1},
\]
并满足
\[
r_n^{(1)}=\mathbb E_\pi(r_n^{(2)}),
\]
其中 \(\mathbb E_\pi\) 表示沿 \(\pi\) 的纤维平均。则下列三条等价：
\begin{enumerate}
  \item
  \[
  \boxed{\;
  \|r_n^{(1)}(\cdot;u)-1\|_{L^2(G_1)}
  =
  \|r_n^{(2)}(\cdot;u)-1\|_{L^2(G_2)}.\;}
  \]
  \item \(r_n^{(2)}(\cdot;u)\) 在 \(\pi\)-纤维上常值，即
  \[
  \boxed{\;
  r_n^{(2)}(c;u)=r_n^{(2)}(c';u)
  \qquad
  \text{只要 }\pi(c)=\pi(c').\;}
  \]
  \item 若把细层扭曲 primitive 谱写作
  \[
  B_n^{(2)}(\chi;u)
  \qquad
  (\chi\in \widehat{G_2}),
  \]
  则一切不来自 \(\pi^\ast(\widehat{G_1})\) 的非平凡角色通道都严格消失：
  \[
  \boxed{\;
  B_n^{(2)}(\chi;u)=0
  \qquad
  \forall\,\chi\in \widehat{G_2}\setminus \pi^\ast(\widehat{G_1}).\;}
  \]
\end{enumerate}
换言之，商映射下 \(L^2\) 偏差不发生任何损失，当且仅当全部可见 Chebotarev 能量已经完全落在该商所允许的角色通道内。
\end{theorem}
\begin{proof}
由推论 \ref{cor:kernel-chebotarev-L2-monotone} 的证明，
\[
r_n^{(1)}-1=\mathbb E_\pi(r_n^{(2)}-1).
\]
通过把 \(L^2(G_1)\) 中的函数 \(h\) 拉回为 \(h\circ\pi\)，可把 \(L^2(G_1)\) 识别为 \(G_2\) 上的纤维常值函数子空间；在此识别下，\(\mathbb E_\pi\) 正是 \(L^2(G_2)\) 到该子空间的正交投影。因此对任意 \(f\in L^2(G_2)\) 都有
\[
\|\mathbb E_\pi f\|_2=\|f\|_2
\iff
f\in \operatorname{Im}(\mathbb E_\pi).
\]
取 \(f=r_n^{(2)}-1\)，便得第 \(1\) 条与第 \(2\) 条等价。

再证第 \(2\) 条与第 \(3\) 条等价。有限阿贝尔群上，沿 \(\pi\) 纤维常值的函数恰是那些在 \(\ker\pi\) 上平凡的角色所张成的子空间；等价地，
\[
\operatorname{Im}(\mathbb E_\pi)
=
\operatorname{span}\bigl(\pi^\ast(\widehat{G_1})\bigr).
\]
另一方面，由命题 \ref{prop:kernel-chebotarev-fourier} 与
\[
r_n^{(2)}(c;u)
=
\frac{|G_2|\,B_{n,c}^{(2)}(u)}{B_n^{(2)}(\mathbf{1};u)}
\]
可得
\[
r_n^{(2)}(c;u)-1
=
\frac{1}{B_n^{(2)}(\mathbf{1};u)}
\sum_{\chi\in \widehat{G_2}\setminus\{\mathbf 1\}}
\overline{\chi(c)}\,B_n^{(2)}(\chi;u).
\]
因此 \(r_n^{(2)}-1\) 属于 \(\operatorname{span}(\pi^\ast(\widehat{G_1}))\)，当且仅当所有不属于 \(\pi^\ast(\widehat{G_1})\) 的角色系数全部为零。这正是第 \(3\) 条。证毕。
\end{proof}

\begin{theorem}[有限阿贝尔标签系统存在唯一最小偏差承载商]\label{thm:xi-time-part64aa-abelian-minimal-deviation-carrier-quotient}
\leanverified{paper\_xi\_time\_part64aa\_abelian\_minimal\_deviation\_carrier\_quotient}
在定理 \ref{thm:xi-time-part64aa-abelian-quotient-equality-criterion} 的阿贝尔设定下，固定 \(u\in(0,1]\) 与 \(n\ge 1\)。定义活跃角色集
\[
\mathcal A_{n,u}
:=
\{\chi\in \widehat G\setminus\{\mathbf 1\}:\ B_n(\chi;u)\neq 0\},
\]
并定义其核交
\[
K_{n,u}
:=
\bigcap_{\chi\in \mathcal A_{n,u}}\ker\chi,
\]
其中在 \(\mathcal A_{n,u}=\varnothing\) 时约定 \(K_{n,u}=G\)。则有：
\begin{enumerate}
  \item 存在唯一函数
  \[
  \widetilde r_{n,u}:G/K_{n,u}\to \CC
  \]
  使得
  \[
  \boxed{\;
  r_n(c;u)=\widetilde r_{n,u}(cK_{n,u})
  \qquad(\forall\,c\in G).\;}
  \]
  \item 对任意有限商
  \[
  \pi:G\twoheadrightarrow H,
  \qquad
  N:=\ker\pi,
  \]
  下列三条等价：
  \[
  \boxed{\;
  r_n(\cdot;u)\ \text{通过 }\pi\text{ 因子化}\;}
  \iff
  \boxed{\;
  N\subseteq K_{n,u}\;}
  \iff
  \boxed{\;
  \|r_n^{H}-1\|_{L^2(H)}
  =
  \|r_n^{G}-1\|_{L^2(G)}.\;}
  \]
  其中 \(r_n^{G}:=r_n\)，而 \(r_n^{H}\) 表示由 \(\pi\) 推前得到的商层相对 Haar 密度。
  \item 因而 \(G/K_{n,u}\) 是保持全部 \(n\)-步 Chebotarev 偏差的唯一最小商：任何其他保持全部偏差的商都唯一地再向下满射到它。
\end{enumerate}
\end{theorem}
\begin{proof}
先证第 \(1\) 条。若 \(k\in K_{n,u}\)，则对每个活跃角色 \(\chi\in\mathcal A_{n,u}\) 都有 \(\chi(k)=1\)。由定理 \ref{thm:xi-time-part64aa-abelian-quotient-equality-criterion} 证明中同一 Fourier 展开，
\[
r_n(ck;u)-1
=
\frac{1}{B_n(\mathbf 1;u)}
\sum_{\chi\in\mathcal A_{n,u}}
\overline{\chi(c)}\,\overline{\chi(k)}\,B_n(\chi;u)
=
r_n(c;u)-1.
\]
故 \(r_n(\cdot;u)\) 在 \(K_{n,u}\)-陪集上常值，从而通过商 \(G/K_{n,u}\) 因子化。唯一性由商映射满射立得。

再证第 \(2\) 条。若 \(r_n(\cdot;u)\) 通过 \(\pi\) 因子化，则它在每条 \(N\)-陪集上常值。由定理 \ref{thm:xi-time-part64aa-abelian-quotient-equality-criterion}，这等价于所有活跃角色都在 \(N\) 上平凡，即
\[
N\subseteq \ker\chi
\qquad
(\forall\,\chi\in\mathcal A_{n,u}),
\]
从而 \(N\subseteq K_{n,u}\)。反之若 \(N\subseteq K_{n,u}\)，则每个活跃角色都在 \(N\) 上平凡，故同一 Fourier 展开只含来自 \(\widehat{G/N}\) 的角色，于是 \(r_n(\cdot;u)\) 必通过 \(\pi\) 因子化。最后，把定理 \ref{thm:xi-time-part64aa-abelian-quotient-equality-criterion} 应用于 \(G\twoheadrightarrow H\)，便得到与 \(L^2\) 等号的等价性。

第 \(3\) 条是第 \(2\) 条的直接推论：若某商 \(G/N\) 保持全部偏差，则必有 \(N\subseteq K_{n,u}\)，因此存在唯一满同态
\[
G/N\twoheadrightarrow G/K_{n,u}.
\]
故 \(G/K_{n,u}\) 在所有保持偏差的商中是唯一最小者。证毕。
\end{proof}

\begin{theorem}[非阿贝尔 Chebotarev--Plancherel 商单调性的等号判据与最小类函数商]\label{thm:xi-time-part64aa-nonabelian-equality-criterion-minimal-class-quotient}
\leanverified{paper\_xi\_time\_part64aa\_nonabelian\_equality\_criterion\_minimal\_class\_quotient}
沿用命题 \ref{prop:kernel-chebotarev-plancherel-nonabelian} 与定理 \ref{thm:xi-time-part64-nonabelian-quotient-energy-loss-exact} 的记号。设 \(G\) 为有限群，\(\varepsilon\in\mathrm{Cl}(G)\) 为零均值类函数，并令
\[
\varphi:G\twoheadrightarrow H
\]
为满同态。把 \(\varepsilon\) 的不可约特征标展开写为
\[
\varepsilon=\sum_{\rho\in\widehat G\setminus\{\mathbf 1\}} a_\rho\,\chi_\rho.
\]
则下列三条等价：
\begin{enumerate}
  \item
  \[
  \boxed{\;
  \|\varphi_\ast\varepsilon\|_{L^2(H)}
  =
  \|\varepsilon\|_{L^2(G)}.\;}
  \]
  \item \(\varepsilon\) 在 \(\varphi\)-纤维上常值，即
  \[
  \boxed{\;
  \varepsilon(g)=\varepsilon(g')
  \qquad
  \text{只要 }\varphi(g)=\varphi(g').\;}
  \]
  \item 一切不通过 \(H\) 的不可约通道都不出现，即
  \[
  \boxed{\;
  a_\rho=0
  \qquad
  \forall\,\rho\in\widehat G
  \text{ 满足 }
  \ker\varphi\nsubseteq \ker\rho.\;}
  \]
\end{enumerate}

进一步，定义
\[
K(\varepsilon)
:=
\bigcap_{a_\rho\neq 0}\ker\rho,
\]
并在 \(\varepsilon=0\) 时约定 \(K(\varepsilon)=G\)。则 \(\varepsilon\) 唯一通过商 \(G/K(\varepsilon)\) 因子化，并且该商是保持全部类函数能量的唯一最小商。
\end{theorem}
\begin{proof}
先证第 \(1\) 条与第 \(2\) 条等价。命题 \ref{prop:kernel-chebotarev-plancherel-nonabelian} 中的 \(\varphi_\ast\) 正是类函数空间上的纤维平均条件期望；而把 \(\mathrm{Cl}(H)\) 中的类函数 \(f\) 拉回为 \(f\circ\varphi\) 后，\(\mathrm{Cl}(H)\) 恰可识别为 \(G\) 上纤维常值类函数子空间。在此识别下，
\[
\|\varphi_\ast\varepsilon\|_2=\|\varepsilon\|_2
\iff
\varepsilon\in \operatorname{Im}(\varphi_\ast),
\]
而右侧正是纤维常值性。

再证第 \(1\) 条与第 \(3\) 条等价。由定理 \ref{thm:xi-time-part64-nonabelian-quotient-energy-loss-exact}，
\[
\|\varepsilon\|_{L^2(G)}^2-\|\varphi_\ast\varepsilon\|_{L^2(H)}^2
=
\frac1{|G|^2}
\sum_{\rho\notin \mathrm{Inf}_\varphi(\widehat H)}
|B_G(\rho)|^2.
\]
又由
\[
a_\rho=\langle \varepsilon,\chi_\rho\rangle
=
\frac{1}{|G|}B_G(\rho),
\]
可知上式右端为零，当且仅当一切不来自 \(\mathrm{Inf}_\varphi(\widehat H)\) 的 \(\rho\) 都满足 \(a_\rho=0\)。而 \(\rho\) 属于 \(\mathrm{Inf}_\varphi(\widehat H)\) 的充要条件正是 \(\ker\varphi\subseteq \ker\rho\)。故第 \(1\) 条与第 \(3\) 条等价。

最后证明最小商断言。若 \(k\in K(\varepsilon)\)，则对每个出现的 \(\rho\) 都有 \(\rho(k)=I\)，故
\[
\chi_\rho(gk)=\chi_\rho(g)
\qquad(\forall\,g\in G).
\]
对全部出现通道求和便得 \(\varepsilon(gk)=\varepsilon(g)\)，从而 \(\varepsilon\) 通过 \(G/K(\varepsilon)\) 因子化。反之，若 \(\varepsilon\) 通过某个商 \(G/N\) 因子化，则由上面三条等价性可知每个出现的 \(\rho\) 都满足 \(N\subseteq \ker\rho\)，于是 \(N\subseteq K(\varepsilon)\)。因此存在唯一满同态
\[
G/N\twoheadrightarrow G/K(\varepsilon),
\]
故 \(G/K(\varepsilon)\) 是唯一最小类函数商。证毕。
\end{proof}

\leanverified{paper\_xi\_time\_part64aa\_optimal\_chebotarev\_witness\_sign\_slice}
\begin{theorem}[二主项最优见证集由极大通道特征函数的符号层给出]\label{thm:xi-time-part64aa-optimal-chebotarev-witness-sign-slice}
沿用定理 \ref{thm:kernel-chebotarev-second-main-term-witness} 的设定。固定 \(u\in(0,1]\)，并取达到
\[
\Lambda(u)=\max_{\rho\neq \mathbf 1}\operatorname{spr}(M_{K,\rho}(u))
\]
的一个非平凡通道 \(\rho_\ast\) 及其对应本征值 \(\mu_\ast\)，并假设该通道在 \(\mu_\ast\) 处具有谱隙。对每个共轭不变集合 \(A\subseteq G\)，记
\[
U_A:=\pi^{-1}(A),
\]
并把定理 \ref{thm:kernel-chebotarev-second-main-term-witness} 中的二主项系数记作 \(c_A:=c_{U_A}\)。则存在唯一零均值类函数
\[
\psi_{\ast,u}\in \mathrm{Cl}(G)
\]
使得对所有共轭不变集合 \(A\subseteq G\) 都有
\[
\boxed{\;
c_A=\langle 1_A,\psi_{\ast,u}\rangle_{\mathrm{Cl}(G)}.\;}
\]
对每个 \(\theta\in\RR\)，再定义共轭类并
\[
A_\theta
:=
\bigcup_{\substack{C\in \mathrm{Conj}(G)\\
\Re(e^{-i\theta}\psi_{\ast,u}(C))>0}} C.
\]
则最优见证振幅满足精确变分公式
\[
\boxed{\;
\sup_{A\subseteq G\ \mathrm{class\ union}}|c_A|
=
\frac12
\sup_{\theta\in\RR}
\sum_{C\in \mathrm{Conj}(G)}
\frac{|C|}{|G|}
\left|\Re\!\bigl(e^{-i\theta}\psi_{\ast,u}(C)\bigr)\right|.\;}
\]
并且每个达到上确界的共轭不变集合都可取为某个 \(A_\theta\)，在满足
\[
\Re(e^{-i\theta}\psi_{\ast,u}(C))=0
\]
的零层共轭类上可任意取舍。
\end{theorem}
\begin{proof}
定理 \ref{thm:kernel-chebotarev-second-main-term-witness} 的证明已经表明：\(c_A\) 来自 \(1_A\) 在固定通道 \(\rho_\ast\) 上的 Fourier 分量与固定左右本征向量（等价地，谱投影）的线性配对。因此映射
\[
1_A\longmapsto c_A
\]
在类函数空间 \(\mathrm{Cl}(G)\) 上线性延拓为一个连续复线性泛函 \(L_{\ast,u}\)。由于 \(\mathrm{Cl}(G)\) 是有限维 Hilbert 空间，存在唯一 \(\psi_{\ast,u}\in \mathrm{Cl}(G)\) 使
\[
L_{\ast,u}(f)=\langle f,\psi_{\ast,u}\rangle_{\mathrm{Cl}(G)}
\qquad(\forall\,f\in \mathrm{Cl}(G)).
\]
特别地，对 \(f=1_A\) 即得所示表示公式。

又因当 \(A=G\) 时，\(1_A\) 只有平凡通道分量，而 \(c_G\) 是非平凡极大通道的二主项系数，故 \(c_G=0\)。于是
\[
\langle 1,\psi_{\ast,u}\rangle_{\mathrm{Cl}(G)}=0,
\]
即 \(\psi_{\ast,u}\) 为零均值类函数。

对任意固定 \(\theta\in\RR\)，定义实值零均值类函数
\[
f_\theta:=\Re(e^{-i\theta}\psi_{\ast,u}).
\]
则
\[
|c_A|
=
\sup_{\theta\in\RR}\Re(e^{-i\theta}c_A)
=
\sup_{\theta\in\RR}\langle 1_A,f_\theta\rangle_{\mathrm{Cl}(G)}.
\]
当 \(\theta\) 固定时，\(1_A\) 在每个共轭类 \(C\) 上只能取 \(0\) 或 \(1\)。因此最大化
\[
\langle 1_A,f_\theta\rangle_{\mathrm{Cl}(G)}
=
\sum_{C\in \mathrm{Conj}(G)}\frac{|C|}{|G|}\,1_A(C)\,f_\theta(C)
\]
的最优策略，正是选取所有满足 \(f_\theta(C)>0\) 的共轭类并；这恰给出 \(A_\theta\)。对应最大值为
\[
\sum_{f_\theta(C)>0}\frac{|C|}{|G|}f_\theta(C).
\]
又因 \(f_\theta\) 零均值，
\[
\sum_{C}\frac{|C|}{|G|}f_\theta(C)=0,
\]
故正部分与负部分的总质量相等，从而
\[
\sum_{f_\theta(C)>0}\frac{|C|}{|G|}f_\theta(C)
=
\frac12
\sum_{C\in \mathrm{Conj}(G)}\frac{|C|}{|G|}\,|f_\theta(C)|.
\]
对 \(\theta\) 再取上确界，即得所示变分公式。最后，在满足 \(f_\theta(C)=0\) 的共轭类上，无论取入或取出都不改变目标值，因此一切极值集都可表示为某个 \(A_\theta\) 加上若干零层共轭类的任意选择。证毕。
\end{proof}

\endinput

\paragraph{primitive Chebotarev 类分布的 Haar 收敛率、Plancherel 偏差指数与 Artin cofactor 半径}
在非阿贝尔 Chebotarev--Plancherel 能量恒等式、正规覆盖算子的 Perron 主相刚性、二主项见证集合以及最小类函数商都已建立之后，还可把 primitive Chebotarev / Peter--Weyl 通道化链进一步压缩为下列四条直接后果。它们表明：一旦把类分布偏差写成 centered class function，平凡通道便会在归一化后完全消失；于是真正控制统计收敛率、平方能量与去主因子 cofactor 最近零点半径的，恰是非平凡通道主谱半径 \(\eta(u)\)。

\begin{theorem}[primitive Chebotarev 类分布到 Haar 分布的精确指数收敛率]\label{thm:xi-time-part64ab-primitive-chebotarev-haar-tv-exponent}
沿用定理 \ref{thm:xi-time-part64-cover-perron-root-trivial-channel}、定理
\ref{thm:kernel-weighted-prime-orbit}、定理
\ref{thm:kernel-chebotarev-second-main-term-witness} 与定理
\ref{thm:xi-time-part64aa-optimal-chebotarev-witness-sign-slice} 的记号。设
\[
\lambda_1(u):=\rho(M_{K,\mathbf 1}(u)),
\qquad
\eta(u):=\max_{\rho\neq \mathbf 1}\rho(M_{K,\rho}(u)).
\]
对每个共轭类 \(C\in \mathrm{Conj}(G)\)，记 primitive Chebotarev 计数为
\[
B_{n,C}(u),
\]
并记总计数为
\[
B_n(u):=\sum_{C\in \mathrm{Conj}(G)}B_{n,C}(u).
\]
定义共轭类分布
\[
p_{n,u}(C):=\frac{B_{n,C}(u)}{B_n(u)},
\qquad
\mu_{\mathrm{Haar}}(C):=\frac{|C|}{|G|},
\]
以及总变差距离
\[
\|p_{n,u}-\mu_{\mathrm{Haar}}\|_{\mathrm{TV}}
:=
\sup_{A\subseteq G\ \mathrm{class\ union}}
\bigl|p_{n,u}(A)-\mu_{\mathrm{Haar}}(A)\bigr|.
\]
则有
\[
\boxed{\;
\|p_{n,u}-\mu_{\mathrm{Haar}}\|_{\mathrm{TV}}
=
O\!\left(\Bigl(\frac{\eta(u)}{\lambda_1(u)}\Bigr)^n\right).\;}
\]

若进一步存在唯一非平凡通道 \(\rho_\ast\neq \mathbf 1\) 与本征值 \(\mu_\ast\in\Spec(M_{K,\rho_\ast}(u))\) 满足
\[
|\mu_\ast|=\eta(u),
\]
且该通道在 \(\mu_\ast\) 处具有谱隙，则
\[
\boxed{\;
\limsup_{n\to\infty}
\|p_{n,u}-\mu_{\mathrm{Haar}}\|_{\mathrm{TV}}^{1/n}
=
\frac{\eta(u)}{\lambda_1(u)}.\;}
\]
\end{theorem}
\begin{proof}
对任意共轭类并 \(A\subseteq G\)，记
\[
B_{n,A}(u):=\sum_{C\subseteq A}B_{n,C}(u),
\qquad
\varepsilon_n(A;u):=
B_{n,A}(u)-\mu_{\mathrm{Haar}}(A)\,B_n(u).
\]
如定理 \ref{thm:kernel-chebotarev-second-main-term-witness} 的证明所示，\(B_{n,A}(u)\) 可按 Peter--Weyl 通道分解为平凡通道与全部非平凡通道的线性组合；其中平凡通道部分恰为
\[
\mu_{\mathrm{Haar}}(A)\,B_n(u).
\]
因此 \(\varepsilon_n(A;u)\) 只由非平凡通道贡献组成。对每个非平凡 \(\rho\neq \mathbf 1\)，标准谱半径估计给出相应通道 primitive 系数为
\[
O\!\left(\frac{\eta(u)^n}{n}\right),
\]
故
\[
\varepsilon_n(A;u)=O\!\left(\frac{\eta(u)^n}{n}\right).
\]
由于共轭类并的总数有限，上述估计可在 \(A\) 上统一。

另一方面，由定理 \ref{thm:kernel-weighted-prime-orbit}，
\[
B_n(u)=\frac{\lambda_1(u)^n}{n}+O\!\left(\frac{|\lambda_2(u)|^n}{n}\right)
\sim \frac{\lambda_1(u)^n}{n}.
\]
于是对任意共轭类并 \(A\) 都有
\[
\bigl|p_{n,u}(A)-\mu_{\mathrm{Haar}}(A)\bigr|
=
\frac{|\varepsilon_n(A;u)|}{B_n(u)}
=
O\!\left(\Bigl(\frac{\eta(u)}{\lambda_1(u)}\Bigr)^n\right).
\]
对 \(A\) 取上确界，即得总变差上界。

若唯一主导通道 \(\rho_\ast\) 存在，则由定理
\ref{thm:kernel-chebotarev-second-main-term-witness} 与定理
\ref{thm:xi-time-part64aa-optimal-chebotarev-witness-sign-slice}，可取某个共轭类并
\[
A_\ast\subseteq G
\]
使得对应二主项系数 \(c_{A_\ast}\neq 0\)。于是沿某个无穷子序列有
\[
|\varepsilon_n(A_\ast;u)|
\ge
c\,\frac{\eta(u)^n}{n}
\]
成立，其中 \(c>0\)。再除以
\[
B_n(u)\sim \frac{\lambda_1(u)^n}{n}
\]
得到
\[
\bigl|p_{n,u}(A_\ast)-\mu_{\mathrm{Haar}}(A_\ast)\bigr|
\ge
c'\Bigl(\frac{\eta(u)}{\lambda_1(u)}\Bigr)^n
\]
在无穷多个 \(n\) 上成立。由总变差的定义，
\[
\|p_{n,u}-\mu_{\mathrm{Haar}}\|_{\mathrm{TV}}
\ge
\bigl|p_{n,u}(A_\ast)-\mu_{\mathrm{Haar}}(A_\ast)\bigr|,
\]
故得到相反方向的 \(\limsup\) 下界。与前面的上界合并，即得结论。证毕。
\end{proof}

\begin{theorem}[Chebotarev--Plancherel 偏差的平方能量具有精确指数率]\label{thm:xi-time-part64ab-chebotarev-plancherel-l2-exponent}
沿用定理 \ref{thm:xi-time-part64ab-primitive-chebotarev-haar-tv-exponent} 的记号。定义 centered class-count 偏差
\[
\varepsilon_n(C;u):=
B_{n,C}(u)-\mu_{\mathrm{Haar}}(C)\,B_n(u),
\]
并定义其按 Haar 权归一化的平方能量
\[
\mathcal E_n(u)
:=
\sum_{C\in\mathrm{Conj}(G)}
\frac{|\varepsilon_n(C;u)|^2}{\mu_{\mathrm{Haar}}(C)}.
\]
则有
\[
\boxed{\;
\mathcal E_n(u)
=
O\!\left(\frac{\eta(u)^{2n}}{n^2}\right).\;}
\]
若再处于定理 \ref{thm:xi-time-part64ab-primitive-chebotarev-haar-tv-exponent} 的唯一主导通道设定下，则
\[
\boxed{\;
\limsup_{n\to\infty}\mathcal E_n(u)^{1/n}
=
\eta(u)^2.\;}
\]
\end{theorem}
\begin{proof}
由定义，
\[
\varepsilon_n(C;u)=B_n(u)\bigl(p_{n,u}(C)-\mu_{\mathrm{Haar}}(C)\bigr),
\]
故
\[
\mathcal E_n(u)
=
B_n(u)^2
\sum_{C\in\mathrm{Conj}(G)}
\frac{|p_{n,u}(C)-\mu_{\mathrm{Haar}}(C)|^2}{\mu_{\mathrm{Haar}}(C)}.
\]
由于共轭类个数有限，而定理
\ref{thm:xi-time-part64ab-primitive-chebotarev-haar-tv-exponent} 已给出逐类一致估计
\[
p_{n,u}(C)-\mu_{\mathrm{Haar}}(C)
=
O\!\left(\Bigl(\frac{\eta(u)}{\lambda_1(u)}\Bigr)^n\right),
\]
遂得
\[
\sum_{C\in\mathrm{Conj}(G)}
\frac{|p_{n,u}(C)-\mu_{\mathrm{Haar}}(C)|^2}{\mu_{\mathrm{Haar}}(C)}
=
O\!\left(\Bigl(\frac{\eta(u)}{\lambda_1(u)}\Bigr)^{2n}\right).
\]
再由
\[
B_n(u)\sim \frac{\lambda_1(u)^n}{n},
\]
即得
\[
\mathcal E_n(u)
=
O\!\left(\frac{\eta(u)^{2n}}{n^2}\right).
\]

在唯一主导通道设定下，利用加权 Cauchy--Schwarz 不等式
\[
\left(\sum_{C}|x_C|\right)^2
\le
\left(\sum_C \frac{|x_C|^2}{\mu_{\mathrm{Haar}}(C)}\right)
\left(\sum_C \mu_{\mathrm{Haar}}(C)\right),
\]
并取
\[
x_C:=p_{n,u}(C)-\mu_{\mathrm{Haar}}(C),
\]
便得到
\[
4\|p_{n,u}-\mu_{\mathrm{Haar}}\|_{\mathrm{TV}}^2
\le
\sum_{C}
\frac{|p_{n,u}(C)-\mu_{\mathrm{Haar}}(C)|^2}{\mu_{\mathrm{Haar}}(C)}.
\]
因此
\[
\mathcal E_n(u)
\ge
4B_n(u)^2\,\|p_{n,u}-\mu_{\mathrm{Haar}}\|_{\mathrm{TV}}^2.
\]
再代入
\[
B_n(u)\sim \frac{\lambda_1(u)^n}{n}
\]
与定理 \ref{thm:xi-time-part64ab-primitive-chebotarev-haar-tv-exponent} 的精确 \(\limsup\) 公式，便得
\[
\limsup_{n\to\infty}\mathcal E_n(u)^{1/n}\ge \eta(u)^2.
\]
与前面上界给出的
\[
\limsup_{n\to\infty}\mathcal E_n(u)^{1/n}\le \eta(u)^2
\]
合并，即得所示等式。证毕。
\end{proof}

\begin{theorem}[Artin cofactor 行列式的最近零点半径精确等于非平凡通道主谱倒数]\label{thm:xi-time-part64ab-artin-cofactor-nearest-zero-radius}
沿用命题 \ref{prop:kernel-artin-factorization} 与定理 \ref{thm:xi-time-part64-cover-perron-root-trivial-channel} 的记号。定义去主因子 cofactor 行列式
\[
\Delta_G^{\mathrm{cof}}(z,u)
:=
\frac{\det(I-zM_{K,G}(u))}{\det(I-zM_{K,\mathbf 1}(u))}
=
\prod_{\rho\neq \mathbf 1}\det(I-zM_{K,\rho}(u))^{d_\rho}.
\]
记其离原点最近的零点半径为
\[
R_{\mathrm{cof}}(u)
:=
\min\bigl\{|z|:\ \Delta_G^{\mathrm{cof}}(z,u)=0\bigr\}.
\]
则有
\[
\boxed{\;
R_{\mathrm{cof}}(u)=\frac{1}{\eta(u)}.\;}
\]
特别地，若满足核版 GRH 型平方根界
\[
\eta(u)\le \lambda_1(u)^{1/2},
\]
则
\[
\boxed{\;
R_{\mathrm{cof}}(u)\ge \lambda_1(u)^{-1/2}.\;}
\]
\leanverified{paper\_xi\_time\_part64ab\_artin\_cofactor\_nearest\_zero\_radius}
\end{theorem}
\begin{proof}
由命题 \ref{prop:kernel-artin-factorization}，
\[
\Delta_G^{\mathrm{cof}}(z,u)
=
\prod_{\rho\neq \mathbf 1}\det(I-zM_{K,\rho}(u))^{d_\rho}.
\]
因此它的全部零点恰由所有非平凡通道矩阵 \(M_{K,\rho}(u)\) 的特征值倒数构成。故最近零点半径满足
\[
R_{\mathrm{cof}}(u)
=
\min_{\rho\neq \mathbf 1}\ \min_{\lambda\in\Spec(M_{K,\rho}(u))}
\frac{1}{|\lambda|}
=
\frac{1}{\max_{\rho\neq \mathbf 1}\rho(M_{K,\rho}(u))}
=
\frac{1}{\eta(u)}.
\]
这就证明了第一式。

若再有
\[
\eta(u)\le \lambda_1(u)^{1/2},
\]
则取倒数即得
\[
R_{\mathrm{cof}}(u)\ge \lambda_1(u)^{-1/2}.
\]
证毕。
\end{proof}

\begin{theorem}[群商不会恶化 Chebotarev 误差指数率]\label{thm:xi-time-part64ab-quotient-monotonicity-error-exponent}
\leanverified{paper\_xi\_time\_part64ab\_quotient\_monotonicity\_error\_exponent}
设
\[
\varphi:G\twoheadrightarrow H
\]
为有限群满同态。把核 \(K\) 的边标签经 \(\varphi\) 推送后得到 \(H\)-标号商核，记其非平凡通道主谱半径为
\[
\eta_G(u),\qquad \eta_H(u),
\]
而平凡通道 Perron 根统一记为
\[
\lambda_1(u).
\]
则有
\[
\boxed{\;
\eta_H(u)\le \eta_G(u).\;}
\]
因而同时成立
\[
\boxed{\;
\frac{\eta_H(u)}{\lambda_1(u)}
\le
\frac{\eta_G(u)}{\lambda_1(u)},
\qquad
R_{\mathrm{cof},H}(u)\ge R_{\mathrm{cof},G}(u).\;}
\]
换言之，群商不会恶化 primitive Chebotarev 类分布向 Haar 分布的指数收敛率，也不会缩小去主因子 cofactor 行列式的最近零点半径。
\end{theorem}
\begin{proof}
任取 \(H\) 的一个非平凡不可约表示
\[
\sigma\in\widehat H\setminus\{\mathbf 1\}.
\]
其 inflation
\[
\widetilde\sigma:=\sigma\circ\varphi
\]
是 \(G\) 的非平凡不可约表示。由通道矩阵的定义，
\[
M_{K_H,\sigma}(u)=M_{K,\widetilde\sigma}(u).
\]
因此
\[
\rho(M_{K_H,\sigma}(u))
=
\rho(M_{K,\widetilde\sigma}(u))
\le
\eta_G(u).
\]
对所有非平凡 \(\sigma\) 取上确界，便得
\[
\eta_H(u)\le \eta_G(u).
\]

另一方面，平凡表示在推送下保持平凡，故
\[
M_{K_H,\mathbf 1}(u)=M_{K,\mathbf 1}(u),
\]
从而平凡通道 Perron 根不变，即两侧都等于 \(\lambda_1(u)\)。于是立得
\[
\frac{\eta_H(u)}{\lambda_1(u)}
\le
\frac{\eta_G(u)}{\lambda_1(u)}.
\]
最后，由定理 \ref{thm:xi-time-part64ab-artin-cofactor-nearest-zero-radius}
\[
R_{\mathrm{cof}}(u)=\frac{1}{\eta(u)},
\]
可知
\[
R_{\mathrm{cof},H}(u)=\frac{1}{\eta_H(u)}
\ge
\frac{1}{\eta_G(u)}
=
R_{\mathrm{cof},G}(u).
\]
证毕。
\end{proof}

\noindent
由此，Chebotarev--Plancherel / Peter--Weyl 通道化链在 primitive 层进一步闭合：归一化共轭类分布向 Haar 分布的精确指数率不再由平凡通道内部的次谱所控制，而是完全由非平凡通道主谱半径 \(\eta(u)\) 决定；相应的平方偏差能量则以 \(\eta(u)^2\) 为精确指数率衰减；在解析侧，Artin 因子化中的去主因子 cofactor 最近零点半径也正好等于 \(\eta(u)^{-1}\)；而在函子侧，群商只会压低 \(\eta(u)\)，因而只能改进而不会恶化这些统计与解析指数。于是，primitive Chebotarev 收敛、Plancherel 偏差、Artin cofactor 半径与 quotient monotonicity 便被压缩为同一个非平凡通道主谱参数。

\endinput

\paragraph{Chebotarev 商熵、可见角色能量与平方根统计障碍}
在 primitive 同余层的严格 Fourier 反演、subgroup--Möbius 完全反演、商映射下的精确能量损失、非阿贝尔通道能量账本、最小偏差承载商以及主导通道见证集合都已建立之后，还可把群商粗化、Plancherel 通道分裂、次主奇点与 RH/GRH 型统计后果进一步压缩为下列四条结构后果。它们分别把商映射的 \(L^2\) 收缩提升为相对熵链式律，把可见偏差写成可见/隐藏角色能量的精确正交分解，把主导扭曲通道的次主解析障碍严格识别为表示维数加权极点，并把平方根谱隙条件转译为 TV/\(\chi^2\)/KL 的统一指数律及其见证性反证。

\begin{theorem}[群商下 Chebotarev 分布相对均匀基线的 KL 链式律]\label{thm:xi-time-part64ac-abelian-quotient-kl-chain}
沿用推论 \ref{cor:kernel-chebotarev-L2-monotone} 的记号。设
\[
\pi:G_2\twoheadrightarrow G_1
\]
为有限阿贝尔群满同态，并记 \(K:=\ker\pi\)。对固定 \(u\in(0,1]\) 与 \(n\ge 1\)，定义两层 primitive Chebotarev 概率分布
\[
\mu_{n,u}^{(i)}(c_i):=\frac{B_{n,c_i}^{(i)}(u)}{B_n(u)},
\qquad
B_n(u):=B_n^{(1)}(\mathbf{1};u)=B_n^{(2)}(\mathbf{1};u),
\qquad
i\in\{1,2\},
\]
并记 \(u_{G_i}\) 为 \(G_i\) 上的均匀分布。对每个满足 \(\mu_{n,u}^{(1)}(c_1)>0\) 的 \(c_1\in G_1\)，定义纤维条件分布
\[
\mu_{n,u}^{(c_1),\mathrm{fib}}(c_2)
:=
\frac{\mu_{n,u}^{(2)}(c_2)}{\mu_{n,u}^{(1)}(c_1)},
\qquad
c_2\in \pi^{-1}(c_1),
\]
并记 \(u_{\pi^{-1}(c_1)}\) 为该纤维上的均匀分布。则有严格链式恒等式
\[
\boxed{\;
D_{\mathrm{KL}}\!\bigl(\mu_{n,u}^{(2)}\Vert u_{G_2}\bigr)
=
D_{\mathrm{KL}}\!\bigl(\mu_{n,u}^{(1)}\Vert u_{G_1}\bigr)
+
\sum_{c_1\in G_1}\mu_{n,u}^{(1)}(c_1)\,
D_{\mathrm{KL}}\!\bigl(\mu_{n,u}^{(c_1),\mathrm{fib}}\Vert u_{\pi^{-1}(c_1)}\bigr)\; }.
\]
其中对 \(\mu_{n,u}^{(1)}(c_1)=0\) 的纤维约定对应项为 \(0\)。特别地，商映射所丢失的全部相对熵恰等于纤维条件非均匀性的平均 KL 成本。
\leanverified{paper\_xi\_time\_part64ac\_abelian\_quotient\_kl\_chain}
\end{theorem}
\begin{proof}
由于 \(\mu_{n,u}^{(1)}=\pi_\ast\mu_{n,u}^{(2)}\) 且 \(\pi\) 为群满同态，每条纤维大小都等于 \(|K|\)，故
\[
u_{G_2}(c_2)
=
\frac{1}{|G_2|}
=
\frac{1}{|G_1|\,|K|}
=
u_{G_1}(\pi(c_2))\,u_{\pi^{-1}(\pi(c_2))}(c_2).
\]
于是对任意满足 \(\mu_{n,u}^{(2)}(c_2)>0\) 的 \(c_2\in G_2\)，有
\[
\log\frac{\mu_{n,u}^{(2)}(c_2)}{u_{G_2}(c_2)}
=
\log\frac{\mu_{n,u}^{(1)}(\pi(c_2))}{u_{G_1}(\pi(c_2))}
+
\log\frac{\mu_{n,u}^{(\pi(c_2)),\mathrm{fib}}(c_2)}{u_{\pi^{-1}(\pi(c_2))}(c_2)}.
\]
对 \(\mu_{n,u}^{(2)}\) 取期望并按纤维分组，即得所示分解；这也是引理 \ref{lem:pom-fiberwise-kl} 在等势群商情形下的直接专门化。证毕。
\end{proof}

\begin{theorem}[可见/隐藏角色能量的精确正交分解]\label{thm:xi-time-part64ac-visible-hidden-character-energy-splitting}
沿用定理 \ref{thm:xi-time-part64ac-abelian-quotient-kl-chain} 的设定，并记
\[
r_n^{(i)}(c_i;u):=|G_i|\,\mu_{n,u}^{(i)}(c_i)
=
\frac{|G_i|\,B_{n,c_i}^{(i)}(u)}{B_n(u)},
\qquad i\in\{1,2\}.
\]
把 \(\widehat{G_1}\) 通过 \(\pi^\ast\) 识别为
\[
K^\perp:=\{\chi\in \widehat{G_2}:\ \chi|_K\equiv 1\}\subseteq \widehat{G_2},
\]
并记细层角色幅度
\[
B_n^{(2)}(\chi;u):=\sum_{c_2\in G_2}\chi(c_2)\,B_{n,c_2}^{(2)}(u).
\]
则有两组精确分解：
\[
\boxed{\;
\bigl\|r_n^{(2)}(\cdot;u)-1\bigr\|_{L^2(G_2)}^2
=
\bigl\|r_n^{(1)}(\cdot;u)-1\bigr\|_{L^2(G_1)}^2
+
\frac{1}{|G_2|}
\sum_{c_1\in G_1}\sum_{c_2\in\pi^{-1}(c_1)}
\bigl|r_n^{(2)}(c_2;u)-r_n^{(1)}(c_1;u)\bigr|^2\; }.
\]
\[
\boxed{\;
\bigl\|r_n^{(1)}(\cdot;u)-1\bigr\|_{L^2(G_1)}^2
=
\frac{1}{B_n(u)^2}
\sum_{\chi\in K^\perp\setminus\{\mathbf{1}\}}
\bigl|B_n^{(2)}(\chi;u)\bigr|^2\; }.
\]
\[
\boxed{\;
\bigl\|r_n^{(2)}(\cdot;u)-1\bigr\|_{L^2(G_2)}^2
-
\bigl\|r_n^{(1)}(\cdot;u)-1\bigr\|_{L^2(G_1)}^2
=
\frac{1}{B_n(u)^2}
\sum_{\chi\in \widehat{G_2}\setminus K^\perp}
\bigl|B_n^{(2)}(\chi;u)\bigr|^2\; }.
\]
因而，群商仅保留在 \(K\) 上平凡的可见角色通道，并精确删除全部隐藏通道能量。
\end{theorem}
\begin{proof}
由推论 \ref{cor:kernel-chebotarev-L2-monotone} 的证明，
\[
r_n^{(1)}(\cdot;u)\circ\pi
=
\mathbb{E}_\pi\!\bigl[r_n^{(2)}(\cdot;u)\bigr]
\]
正是 \(r_n^{(2)}(\cdot;u)\) 沿纤维常值子空间的条件期望。故在 \(L^2(G_2)\) 中可应用正交投影的勾股恒等式：
\[
\bigl\|r_n^{(2)}-1\bigr\|_2^2
=
\bigl\|r_n^{(1)}\circ\pi-1\bigr\|_2^2
+
\bigl\|r_n^{(2)}-r_n^{(1)}\circ\pi\bigr\|_2^2.
\]
按纤维展开第二项，即得第一式。

另一方面，推论 \ref{cor:kernel-chebotarev-plancherel-energy} 应用于细层给出
\[
\bigl\|r_n^{(2)}(\cdot;u)-1\bigr\|_{L^2(G_2)}^2
=
\frac{1}{B_n(u)^2}
\sum_{\chi\in \widehat{G_2}\setminus\{\mathbf{1}\}}
\bigl|B_n^{(2)}(\chi;u)\bigr|^2.
\]
而粗层角色正好对应于 \(K^\perp=\pi^\ast(\widehat{G_1})\)，故同理
\[
\bigl\|r_n^{(1)}(\cdot;u)-1\bigr\|_{L^2(G_1)}^2
=
\frac{1}{B_n(u)^2}
\sum_{\chi\in K^\perp\setminus\{\mathbf{1}\}}
\bigl|B_n^{(2)}(\chi;u)\bigr|^2,
\]
得到第二式。将两式相减即得第三式；这也是定理
\ref{thm:xi-time-part63b-primitive-quotient-energy-loss-exact}
在相对 Haar 密度口径下的直接重写。证毕。
\end{proof}

\begin{theorem}[主导非平凡通道的极点阶与见证振荡]\label{thm:xi-time-part64ac-dominant-channel-pole-order-witness-oscillation}
沿用定理 \ref{thm:xi-time-part64-cover-zeta-pole-order-weighted-channel-multiplicity} 与定理 \ref{thm:kernel-chebotarev-second-main-term-witness} 的记号。固定 \(u\in(0,1]\)，并假设存在唯一非平凡通道 \(\rho_\star\neq \mathbf{1}\) 与其简单本征值 \(\mu_\star\in\Spec(M_{K,\rho_\star}(u))\) 满足
\[
|\mu_\star|
=
\max_{\rho\neq \mathbf{1}}\rho\!\bigl(M_{K,\rho}(u)\bigr)
=:\eta(u),
\]
且该通道在 \(\mu_\star\) 处具有谱隙。则有：
\begin{enumerate}
  \item 去主因子覆盖 \(\zeta\) 在 \(z=\mu_\star^{-1}\) 处具有精确极点阶
  \[
  \boxed{\;
  \operatorname{ord}_{z=\mu_\star^{-1}}
  \frac{\zeta_{K,G}(z,u)}{L_{K,\mathbf{1}}(z,u)}
  =
  \dim\rho_\star\; }.
  \]
  \item 存在某个常数 \(\eta_-(u)<\eta(u)\)，使得对任意共轭不变集合 \(A\subseteq G\)，若其指示函数 \(1_A\) 在 \(\rho_\star\)-通道上的 Fourier 分量非零，则记
  \[
  B_{n,A}(u):=\sum_{C\subseteq A}B_{n,C}(u),
  \qquad
  \mu_{\mathrm{Haar}}(A):=\frac{|A|}{|G|},
  \]
  并存在 \(c_A\neq 0\) 使
  \[
  \boxed{\;
  B_{n,A}(u)
  =
  \mu_{\mathrm{Haar}}(A)\frac{\lambda_1(u)^n}{n}
  +
  \Re\!\bigl(c_A\mu_\star^n\bigr)\frac1n
  +
  O\!\left(\frac{\eta_-(u)^n}{n}\right)\; }.
  \]
  特别地，存在常数 \(c(A)>0\) 使
  \[
  \boxed{\;
  \left|B_{n,A}(u)-\mu_{\mathrm{Haar}}(A)\frac{\lambda_1(u)^n}{n}\right|
  \ge
  c(A)\frac{\eta(u)^n}{n}\; }
  \]
  在无穷多个 \(n\) 上成立。
\end{enumerate}
\end{theorem}
\begin{proof}
由 Artin 因子化
\[
\frac{\zeta_{K,G}(z,u)}{L_{K,\mathbf{1}}(z,u)}
=
\prod_{\rho\neq \mathbf{1}}L_{K,\rho}(z,u)^{\dim\rho}
\]
可知，该商在 \(z=\mu_\star^{-1}\) 处的极点阶只由非平凡通道决定。再由定理 \ref{thm:xi-time-part64-cover-zeta-pole-order-weighted-channel-multiplicity}，
\[
\operatorname{ord}_{z=\mu_\star^{-1}}
\frac{\zeta_{K,G}(z,u)}{L_{K,\mathbf{1}}(z,u)}
=
\sum_{\rho\in\widehat G\setminus\{\mathbf{1}\}}(\dim\rho)\,m_\rho(\mu_\star),
\]
其中 \(m_\rho(\mu_\star)\) 为 \(\mu_\star\) 在 \(M_{K,\rho}(u)\) 中的代数重数。由唯一主导通道假设与 \(\mu_\star\) 的简单性可知，\(\rho_\star\) 贡献 \(m_{\rho_\star}(\mu_\star)=1\)，而其余非平凡通道都不含该特征值，故极点阶等于 \(\dim\rho_\star\)。

第二部分正是定理 \ref{thm:kernel-chebotarev-second-main-term-witness} 的直接应用：唯一主导通道给出单一的次主振荡模态，谱隙使所有其余通道并入某个 \(\eta_-(u)<\eta(u)\) 的统一误差项；若 \(1_A\) 在 \(\rho_\star\)-通道上的分量非零，则对应系数 \(c_A\neq 0\)，从而沿无穷多个 \(n\) 得到 \(\Omega(\eta(u)^n/n)\) 的见证振荡下界。证毕。
\end{proof}

\begin{theorem}[GRH 型平方根谱隙的 TV/\texorpdfstring{$\chi^2$}{chi2}/KL 统计后果]\label{thm:xi-time-part64ac-grh-tv-chi2-kl-barrier}
沿用定理 \ref{thm:xi-time-part64ab-primitive-chebotarev-haar-tv-exponent} 的设定，并进一步假设 \(G\) 为有限阿贝尔群。定义
\[
\mu_{n,u}(c):=\frac{B_{n,c}(u)}{B_n(u)},
\qquad
u_G(c):=\frac{1}{|G|},
\qquad
r_n(c;u):=|G|\,\mu_{n,u}(c),
\]
以及
\[
\eta(u):=\max_{\chi\neq \mathbf{1}}\rho\!\bigl(M_{K,\chi}(u)\bigr).
\]
则有统一指数界
\[
\boxed{\;
\bigl\|\mu_{n,u}-u_G\bigr\|_{\mathrm{TV}}
=
O\!\left(\Bigl(\frac{\eta(u)}{\lambda_1(u)}\Bigr)^n\right),\;}
\]
\[
\boxed{\;
\chi^2\!\bigl(\mu_{n,u}\Vert u_G\bigr)
=
O\!\left(\Bigl(\frac{\eta(u)}{\lambda_1(u)}\Bigr)^{2n}\right),\qquad
D_{\mathrm{KL}}\!\bigl(\mu_{n,u}\Vert u_G\bigr)
=
O\!\left(\Bigl(\frac{\eta(u)}{\lambda_1(u)}\Bigr)^{2n}\right).\;}
\]
特别地，若 \(\textsf{GRH}_{K,G}(u)\) 成立，即
\[
\eta(u)\le \lambda_1(u)^{1/2},
\]
则
\[
\boxed{\;
\bigl\|\mu_{n,u}-u_G\bigr\|_{\mathrm{TV}}=O\!\bigl(\lambda_1(u)^{-n/2}\bigr),\qquad
\chi^2\!\bigl(\mu_{n,u}\Vert u_G\bigr),\
D_{\mathrm{KL}}\!\bigl(\mu_{n,u}\Vert u_G\bigr)
=
O\!\bigl(\lambda_1(u)^{-n}\bigr).\;}
\]
反之，若定理 \ref{thm:xi-time-part64ac-dominant-channel-pole-order-witness-oscillation} 的唯一主导通道假设成立，且
\[
\eta(u)>\lambda_1(u)^{1/2},
\]
则存在某个集合 \(A\subseteq G\) 与常数 \(c(A)>0\)，使
\[
\boxed{\;
\bigl|\mu_{n,u}(A)-u_G(A)\bigr|
\ge
c(A)\Bigl(\frac{\eta(u)}{\lambda_1(u)}\Bigr)^n\; }
\]
在无穷多个 \(n\) 上成立。因而此时不可能存在统一的平方根级 TV 收敛律。
\end{theorem}
\begin{proof}
阿贝尔情形下，共轭类分布即元素分布，故定理 \ref{thm:xi-time-part64ab-primitive-chebotarev-haar-tv-exponent} 直接给出 TV 上界。另一方面，
\[
\chi^2\!\bigl(\mu_{n,u}\Vert u_G\bigr)
=
\frac{1}{|G|}\sum_{c\in G}\bigl(r_n(c;u)-1\bigr)^2.
\]
由推论 \ref{cor:kernel-chebotarev-plancherel-energy} 的相对 Haar 密度版本，
\[
\chi^2\!\bigl(\mu_{n,u}\Vert u_G\bigr)
=
\frac{1}{B_n(u)^2}
\sum_{\chi\in\widehat G\setminus\{\mathbf{1}\}}
\bigl|B_n(\chi;u)\bigr|^2.
\]
对每个非平凡角色 \(\chi\neq \mathbf{1}\)，由命题 \ref{prop:kernel-chebotarev-mobius-adams} 与谱半径估计有
\[
B_n(\chi;u)=O\!\left(\frac{\eta(u)^n}{n}\right),
\]
而 \(B_n(u)\sim \lambda_1(u)^n/n\) 由定理 \ref{thm:kernel-weighted-prime-orbit} 给出。由于非平凡角色数有限，遂得
\[
\chi^2\!\bigl(\mu_{n,u}\Vert u_G\bigr)
=
O\!\left(\Bigl(\frac{\eta(u)}{\lambda_1(u)}\Bigr)^{2n}\right).
\]
再由标准不等式
\[
D_{\mathrm{KL}}(P\Vert Q)\le \chi^2(P\Vert Q),
\]
立即得到 KL 上界。若 \(\eta(u)\le \lambda_1(u)^{1/2}\)，则把该平方根界代入上述三式即得所示结论。

反向方向中，若唯一主导通道存在且 \(\eta(u)>\lambda_1(u)^{1/2}\)，则由定理
\ref{thm:xi-time-part64ac-dominant-channel-pole-order-witness-oscillation}
可取某个 \(A\subseteq G\) 使
\[
\left|B_{n,A}(u)-u_G(A)\frac{\lambda_1(u)^n}{n}\right|
\ge
c(A)\frac{\eta(u)^n}{n}
\]
在无穷多个 \(n\) 上成立。再除以
\[
B_n(u)\sim \frac{\lambda_1(u)^n}{n}
\]
即得
\[
\bigl|\mu_{n,u}(A)-u_G(A)\bigr|
\ge
c'(A)\Bigl(\frac{\eta(u)}{\lambda_1(u)}\Bigr)^n
\]
沿无穷多个 \(n\) 成立。由于总变差是事件偏差的上确界，若再有统一的 \(O(\lambda_1(u)^{-n/2})\) 上界，则会与 \(\eta(u)>\lambda_1(u)^{1/2}\) 矛盾。证毕。
\end{proof}

\noindent
因此，primitive Chebotarev 的 quotient pushforward 不仅满足 \(L^2\) 收缩，而且在相对熵层面具有严格链式记账；相应的可见偏差可精确拆分为可见/隐藏角色通道能量；一旦某个非平凡通道主导全部次主谱，它便同时决定去主因子覆盖 \(\zeta\) 的最近奇点阶数与 Chebotarev 见证集合的振荡量级；而 RH/GRH 型平方根谱隙条件则恰把这些通道参数进一步压缩为 TV/\(\chi^2\)/KL 的统一统计指数与尖锐障碍。于是，商粗化、角色 Fourier、Artin 因子化与平方根误差预算便在同一条有限维通道链上闭合。

\endinput

\paragraph{三域符号盒、最小退化逐点判据与 fold-gauge 的奇次重整化}
在三域 Frobenius 直积的完全张量分解、审计偶数窗最小退化 Fibonacci 律以及 fold-gauge 缺陷序列的 Bernoulli 剥离已经建立之后，还可把条件独立、子覆盖猜想与非线性正规形继续压缩为下列六条后果。前半部分把 \(L_4\) 因子、四条二次特征与全分裂事件收束为严格的符号盒均分与条件独立；后半部分则把最小退化扇区的几何猜想收束为逐点多重度判据，并把 fold-gauge 缺陷的重整化结构提升到五次与纯奇次两个层面。

\begin{theorem}[\texorpdfstring{$L_4$}{L4} 事件对 \texorpdfstring{$(\chi_{10},\chi_{\mathrm{LY}})$}{(chi10, chiLY)} 的完全条件独立]\label{thm:xi-time-part64a-l4-event-complete-conditional-independence}
沿用定理 \ref{thm:xi-time-part62d-triple-frobenius-classfunction-tensor-independence} 的记号。对不分歧素数 \(p\)，记
\[
\sigma_{10}(p)\in S_{10},
\qquad
\sigma_{3}(p)\in S_3,
\qquad
\sigma_{4,\mathrm{tot}}(p):=\bigl(\sigma_4(p),\varepsilon_4(p)\bigr)\in S_4\times C_2,
\]
并定义二次特征
\[
\chi_{10}(p):=\operatorname{sgn}\!\bigl(\sigma_{10}(p)\bigr),
\qquad
\chi_{\mathrm{LY}}(p):=\operatorname{sgn}\!\bigl(\sigma_3(p)\bigr).
\]
任取 \(S_4\times C_2\) 的共轭类并集 \(\mathcal E\)，定义事件
\[
E_{\mathcal E}:=\{\,p:\ \sigma_{4,\mathrm{tot}}(p)\in\mathcal E\,\}.
\]
则对任意 \((\varepsilon_1,\varepsilon_2)\in\{\pm1\}^2\)，有
\[
\delta\!\bigl(
E_{\mathcal E}\cap\{\chi_{10}=\varepsilon_1,\ \chi_{\mathrm{LY}}=\varepsilon_2\}
\bigr)
=
\frac{1}{4}\cdot \frac{|\mathcal E|}{48}.
\]
从而
\[
\delta\!\bigl(
E_{\mathcal E}\mid \chi_{10}=\varepsilon_1,\ \chi_{\mathrm{LY}}=\varepsilon_2
\bigr)
=
\frac{|\mathcal E|}{48}.
\]
特别地，若记
\[
E_{\mathrm{split}}:=\{\text{\(C(X)\bmod p\) 全分裂}\},
\]
则对任意 \((\varepsilon_1,\varepsilon_2)\in\{\pm1\}^2\)，有
\[
\delta\!\bigl(E_{\mathrm{split}}\mid \chi_{10}=\varepsilon_1,\ \chi_{\mathrm{LY}}=\varepsilon_2\bigr)=\frac{1}{48}.
\]
\end{theorem}
\begin{proof}
由定理 \ref{thm:xi-time-part62d-triple-frobenius-classfunction-tensor-independence}，
\[
\Gal(L_{10}L_{\mathrm{LY}}L_4/\QQ)
\cong
S_{10}\times S_3\times (S_4\times C_2),
\]
并且三个因子的 Frobenius 极限分布严格乘积化。

取类函数
\[
f(\tau):=\mathbf 1_{\{\operatorname{sgn}(\tau)=\varepsilon_1\}},
\qquad
g(\upsilon):=\mathbf 1_{\{\operatorname{sgn}(\upsilon)=\varepsilon_2\}},
\qquad
h(\eta):=\mathbf 1_{\mathcal E}(\eta).
\]
由于 \(A_{10}\subset S_{10}\) 与 \(A_3\subset S_3\) 均为指数 \(2\) 子群，故
\[
\langle f\rangle_{S_{10}}=\frac12,
\qquad
\langle g\rangle_{S_3}=\frac12.
\]
而
\[
\langle h\rangle_{S_4\times C_2}=\frac{|\mathcal E|}{48}.
\]
将这三式代入定理 \ref{thm:xi-time-part62d-triple-frobenius-classfunction-tensor-independence} 的乘积公式，便得
\[
\delta\!\bigl(
E_{\mathcal E}\cap\{\chi_{10}=\varepsilon_1,\ \chi_{\mathrm{LY}}=\varepsilon_2\}
\bigr)
=
\frac12\cdot\frac12\cdot\frac{|\mathcal E|}{48}
=
\frac14\cdot\frac{|\mathcal E|}{48}.
\]
再除以
\[
\delta(\chi_{10}=\varepsilon_1,\ \chi_{\mathrm{LY}}=\varepsilon_2)=\frac14
\]
即得条件密度公式。

最后，由定理 \ref{thm:xi-time-part62e-l4-fullsplit-universal-dilution} 的证明，\(C(X)\bmod p\) 全分裂恰对应第三因子的恒等元共轭类，故
\[
|\mathcal E|=1,
\]
于是
\[
\delta\!\bigl(E_{\mathrm{split}}\mid \chi_{10}=\varepsilon_1,\ \chi_{\mathrm{LY}}=\varepsilon_2\bigr)=\frac{1}{48}.
\]
证毕。
\end{proof}

\begin{theorem}[四重二次特征的 \texorpdfstring{$\{\pm1\}^4$}{pm1^4} 等分布]\label{thm:xi-time-part64a-four-quadratic-character-equidistribution}
\leanverified{paper\_xi\_time\_part64a\_four\_quadratic\_character\_equidistribution}
沿用定理 \ref{thm:xi-time-part62e-triple-product-abelianization-c2-four} 的记号，定义
\[
\chi_{4,\mathrm{sgn}}(p):=\operatorname{sgn}\!\bigl(\sigma_4(p)\bigr),
\qquad
\chi_{4,\mathrm{aux}}(p):=\iota\!\bigl(\varepsilon_4(p)\bigr)\in\{\pm1\},
\]
其中 \(\iota:C_2\xrightarrow{\sim}\{\pm1\}\) 为定理
\ref{thm:xi-time-part62e-triple-product-abelianization-c2-four} 中固定的同构。则四元组
\[
\bigl(
\chi_{10}(p),\chi_{\mathrm{LY}}(p),\chi_{4,\mathrm{sgn}}(p),\chi_{4,\mathrm{aux}}(p)
\bigr)
\]
在 \(\{\pm1\}^4\) 上等分布；亦即对任意
\[
(\varepsilon_1,\varepsilon_2,\varepsilon_3,\varepsilon_4)\in\{\pm1\}^4,
\]
都有
\[
\delta\!\Bigl(
\chi_{10}=\varepsilon_1,\ \chi_{\mathrm{LY}}=\varepsilon_2,\ 
\chi_{4,\mathrm{sgn}}=\varepsilon_3,\ \chi_{4,\mathrm{aux}}=\varepsilon_4
\Bigr)
=
\frac{1}{16}.
\]
\end{theorem}
\begin{proof}
仍由定理 \ref{thm:xi-time-part62d-triple-frobenius-classfunction-tensor-independence}，三个因子的 Frobenius 统计严格乘积化。前两个因子的符号条件各占一半，故
\[
\delta(\chi_{10}=\varepsilon_1,\ \chi_{\mathrm{LY}}=\varepsilon_2)=\frac14.
\]

再看第三因子。定义映射
\[
\Theta:S_4\times C_2\longrightarrow \{\pm1\}\times\{\pm1\},
\qquad
\Theta(\sigma,\epsilon):=\bigl(\operatorname{sgn}(\sigma),\iota(\epsilon)\bigr).
\]
该映射显然满射，且其核为
\[
\ker\Theta=A_4\times\{1\},
\]
大小为 \(12\)。故每个符号盒
\[
\Theta^{-1}(\varepsilon_3,\varepsilon_4)
\]
都恰有 \(12\) 个元素，从而
\[
\delta(\chi_{4,\mathrm{sgn}}=\varepsilon_3,\ \chi_{4,\mathrm{aux}}=\varepsilon_4)
=
\frac{12}{48}
=
\frac14.
\]
于是对全部四重符号盒，
\[
\delta\!\Bigl(
\chi_{10}=\varepsilon_1,\ \chi_{\mathrm{LY}}=\varepsilon_2,\ 
\chi_{4,\mathrm{sgn}}=\varepsilon_3,\ \chi_{4,\mathrm{aux}}=\varepsilon_4
\Bigr)
=
\frac14\cdot\frac14
=
\frac{1}{16}.
\]
证毕。
\end{proof}

\begin{corollary}[三域混合事件的显式密度闭式]\label{cor:xi-time-part64a-triple-mixed-event-explicit-density}
沿用推论 \ref{cor:killo-leyang-joint-density-explicit} 的记号，定义
\[
A:=\{\text{\(P_{10}\bmod p\) 含一次因子}\},
\qquad
B_{\mathrm{root}}:=\{\text{\(g\bmod p\) 至少有一根}\},
\]
并令
\[
C_{\mathrm{split}}:=\{\text{\(C(X)\bmod p\) 全分裂}\}.
\]
则有显式密度
\[
\delta(A\cap B_{\mathrm{root}}\cap C_{\mathrm{split}})
=
\frac{28319}{3225600}.
\]
\end{corollary}
\begin{proof}
由推论 \ref{cor:killo-leyang-joint-density-explicit}，
\[
\delta(A\cap B_{\mathrm{root}})=\frac{28319}{67200}.
\]
另一方面，事件 \(A\cap B_{\mathrm{root}}\) 只依赖 \(S_{10}\times S_3\) 两个因子，而
\(C_{\mathrm{split}}\) 只依赖第三因子。故由定理
\ref{thm:xi-time-part62e-l4-fullsplit-universal-dilution}，
\[
\delta(A\cap B_{\mathrm{root}}\cap C_{\mathrm{split}})
=
\frac{1}{48}\,\delta(A\cap B_{\mathrm{root}}).
\]
代入上式即得
\[
\delta(A\cap B_{\mathrm{root}}\cap C_{\mathrm{split}})
=
\frac{1}{48}\cdot \frac{28319}{67200}
=
\frac{28319}{3225600}.
\]
证毕。
\end{proof}

\begin{proposition}[最小退化扇区猜想的逐点等价判据]\label{prop:xi-time-part64a-minsector-pointwise-equivalence}
沿用定义 \ref{def:fold-bin-degeneracy-sectors} 与猜想
\ref{conj:fold-bin-min-sector-regular-subcover} 的记号。设 \(m\ge 6\) 为偶数，并假设
\[
d_{\min}(m)=F_{m/2},
\qquad
\bigl|\mathcal S_{m,d_{\min}(m)}\bigr|=F_m=|X_{m-2}|.
\]
定义候选嵌入
\[
\iota_{m-2}:X_{m-2}\hookrightarrow X_m,
\qquad
\iota_{m-2}(v):=(v,0,1).
\]
则猜想
\[
\mathcal S_{m,d_{\min}(m)}=\iota_{m-2}(X_{m-2})
\]
与下述逐点陈述完全等价：
\[
d_m^{\mathrm{bin}}(v,0,1)=F_{m/2}
\qquad
(\forall\,v\in X_{m-2}).
\]
\end{proposition}
\begin{proof}
若
\[
\mathcal S_{m,d_{\min}(m)}=\iota_{m-2}(X_{m-2}),
\]
则对每个 \(v\in X_{m-2}\)，像点 \(\iota_{m-2}(v)=(v,0,1)\) 属于最小退化扇区，故
\[
d_m^{\mathrm{bin}}(v,0,1)=d_{\min}(m)=F_{m/2}.
\]

反之，若对所有 \(v\in X_{m-2}\) 都有
\[
d_m^{\mathrm{bin}}(v,0,1)=F_{m/2}=d_{\min}(m),
\]
则
\[
\iota_{m-2}(X_{m-2})\subseteq \mathcal S_{m,d_{\min}(m)}.
\]
又 \(\iota_{m-2}\) 显然为单射，因此
\[
\bigl|\iota_{m-2}(X_{m-2})\bigr|
=
|X_{m-2}|
=
F_m
=
\bigl|\mathcal S_{m,d_{\min}(m)}\bigr|.
\]
包含关系与基数相同遂推出
\[
\mathcal S_{m,d_{\min}(m)}=\iota_{m-2}(X_{m-2}).
\]
证毕。
\end{proof}

\begin{theorem}[fold-gauge 缺陷的五次重整化正规形]\label{thm:xi-time-part64a-fold-gauge-quintic-renormalization-normal-form}
沿用定理 \ref{thm:xi-time-part60ab-gauge-constant-recursive-bernoulli-stripping} 与
\ref{thm:xi-fold-gauge-defect-cubic-renormalization-normal-form} 的记号，记
\[
E_m:=\log C_m-\log(2\pi),
\qquad
q:=\frac{\varphi}{2}.
\]
则当 \(m\to\infty\) 时有精确提升
\[
E_{m+1}
=
qE_m
+\frac{3(2+\sqrt5)}{32}\,E_m^3
-\frac{27(655+294\sqrt5)}{4480}\,E_m^5
+O(E_m^7).
\]
特别地，
\[
\lim_{m\to\infty}
\frac{
E_{m+1}
-qE_m
-\frac{3(2+\sqrt5)}{32}\,E_m^3
}{
E_m^5
}
=
-\frac{27(655+294\sqrt5)}{4480}.
\]
\end{theorem}
\begin{proof}
由定理 \ref{thm:xi-time-part60ab-gauge-constant-recursive-bernoulli-stripping} 在 \(r=1,2,3\) 处的系数公式，
\[
E_m
=
aq^m+bq^{3m}+cq^{5m}+O(q^{7m}),
\]
其中
\[
a=\frac{1}{3\varphi},
\qquad
b=
\frac{B_4}{2\cdot 3}\,(\varphi^{-1}+\varphi)
=
-\frac{\sqrt5}{180},
\]
\[
c=
\frac{B_6}{3\cdot 5}\,(\varphi^{-1}+\varphi^3)
=
\frac{1}{42\cdot 15}\cdot 3\varphi
=
\frac{\varphi}{210}.
\]
现取 \(\alpha,\beta\) 使
\[
E_{m+1}-qE_m-\alpha E_m^3-\beta E_m^5=O(q^{7m}),
\]
并比较 \(q^{3m}\) 与 \(q^{5m}\) 的系数。

首先，
\[
E_{m+1}-qE_m
=
b(q^3-q)q^{3m}+c(q^5-q)q^{5m}+O(q^{7m}).
\]
另一方面，
\[
E_m^3=a^3q^{3m}+3a^2b\,q^{5m}+O(q^{7m}),
\qquad
E_m^5=a^5q^{5m}+O(q^{7m}).
\]
比较 \(q^{3m}\) 项可得
\[
\alpha=\frac{b(q^3-q)}{a^3}
=
\frac{3(2+\sqrt5)}{32}.
\]
再比较 \(q^{5m}\) 项，得到
\[
\beta
=
\frac{c(q^5-q)-3\alpha a^2b}{a^5}
=
-\frac{27(655+294\sqrt5)}{4480}.
\]
于是
\[
E_{m+1}
=
qE_m
+\frac{3(2+\sqrt5)}{32}\,E_m^3
-\frac{27(655+294\sqrt5)}{4480}\,E_m^5
+O(q^{7m}).
\]
又因 \(a>0\)，有
\[
E_m\sim aq^m,
\]
故 \(q^{7m}=O(E_m^7)\)。这就得到所示五次正规形。

最后，将等式两边减去前三项并除以 \(E_m^5\)，再利用 \(E_m^7= o(E_m^5)\)，便得到极限公式。证毕。
\end{proof}

\begin{theorem}[fold-gauge 重整化的纯奇次性]\label{thm:xi-time-part64a-fold-gauge-purely-odd-renormalization}
沿用定理 \ref{thm:xi-time-part60ab-gauge-constant-recursive-bernoulli-stripping} 的记号。存在唯一形式幂级数
\[
\mathcal R(z)=\frac{\varphi}{2}z+\sum_{j\ge 1}c_{2j+1}z^{2j+1}
\in z\mathbb R[[z^2]]
\]
使得对任意 \(N\ge 1\)，若记其截断式
\[
\mathcal R_N(z):=\frac{\varphi}{2}z+\sum_{j=1}^{N}c_{2j+1}z^{2j+1},
\]
则当 \(m\to\infty\) 时有
\[
E_{m+1}=\mathcal R_N(E_m)+O(E_m^{2N+3}).
\]
换言之，fold-gauge 缺陷的全部有限阶重整化正规形只含奇次幂，而一切偶次非线性项在所有阶上全部消失。
\end{theorem}
\begin{proof}
由定理 \ref{thm:xi-time-part60ab-gauge-constant-recursive-bernoulli-stripping}，
\[
E_m=\sum_{r\ge 1}a_r q^{(2r-1)m},
\qquad
q:=\frac{\varphi}{2},
\]
其中
\[
a_1=\frac{1}{3\varphi}\neq 0.
\]
令
\[
F(y):=\sum_{r\ge 1}a_r y^{2r-1}\in y\mathbb R[[y^2]].
\]
若记 \(y_m:=q^m\)，则
\[
E_m=F(y_m).
\]

由于 \(F\) 的线性项系数 \(a_1\) 非零，\(F\) 在形式幂级数意义下存在唯一复合逆
\[
F^{-1}(z)\in z\mathbb R[[z]].
\]
又 \(F\) 为奇级数，故 \(F^{-1}\) 亦为奇级数，从而
\[
F^{-1}(z)\in z\mathbb R[[z^2]].
\]
定义
\[
\mathcal R(z):=F\!\bigl(q\,F^{-1}(z)\bigr).
\]
因为 \(F\) 与 \(F^{-1}\) 都是奇级数，故
\[
\mathcal R(z)\in z\mathbb R[[z^2]].
\]

现验证该级数给出所求重整化。由 \(y_{m+1}=qy_m\)，有
\[
E_{m+1}
=
F(y_{m+1})
=
F(qy_m)
=
F\!\bigl(q\,F^{-1}(E_m)\bigr)
=
\mathcal R(E_m).
\]
这说明 \(\mathcal R\) 就是所求的形式重整化映射。

再证唯一性。若另有 \(\widetilde{\mathcal R}(z)\in z\mathbb R[[z]]\) 满足
\[
F(qy)=\widetilde{\mathcal R}(F(y))
\]
作为形式幂级数恒等式成立，则以 \(y=F^{-1}(z)\) 代入得到
\[
\widetilde{\mathcal R}(z)=F\!\bigl(qF^{-1}(z)\bigr)=\mathcal R(z),
\]
故 \(\mathcal R\) 唯一。

最后，对任意 \(N\ge 1\)，令 \(\mathcal R_N\) 为 \(\mathcal R\) 的 \(2N+1\) 阶截断。则
\[
\mathcal R(z)-\mathcal R_N(z)=O(z^{2N+3})
\qquad(z\to 0).
\]
由于 \(E_m\to 0\)，代入 \(z=E_m\) 即得
\[
E_{m+1}
=
\mathcal R(E_m)
=
\mathcal R_N(E_m)+O(E_m^{2N+3}).
\]
证毕。
\end{proof}

\noindent
上述六条结果把三条先前分立的闭合链进一步压缩到两个统一结构层：在算术侧，三域 Frobenius 数据被四维 \(C_2\)-符号盒完全组织，并使 \(L_4\) 因子仅以无偏条件化与 \(1/48\) 稀释的方式介入；在 bin-fold/gauge 侧，最小退化子覆盖猜想被还原为尾部固定 \(01\) 的逐点多重度判据，而 fold-gauge 缺陷的全部高阶重整化则被唯一的纯奇形式自映射严格控制。

\endinput

\paragraph{部分可逆尾量、二原子信息塌缩、乘法外置与条件正交的横向闭合}
在预言机容量闭式、bin-fold 相对均匀基线的两点 \(f\)-散度极限、稳定地址半环的 PTIME 双向可解释性、一般 Ostrowski 数字层的规范化函子性、三轴 CRT 最小窗以及十次分歧/Lee--Yang 三次的 Chebotarev 乘法律已经建立之后，还可把部分可逆误差、信息几何塌缩、数字层局域性、三轴最小窗及跨指纹条件独立进一步压缩为下列八条结论。前四条把最优部分读出的全部损失先后还原为纤维超额尾量、Neyman--Pearson 对偶与矩谱驱动的非渐近失败界；中间三条分别把渐近 \(f\)-几何压缩到唯一二原子似然比测度、把 Zeckendorf/Ostrowski 层的局域可加性与乘法账本不可有限化严格分离，并指出三轴最小窗在最小层面即不可纯化；最后一条则表明十次线性因子指纹对 Lee--Yang 的符号/根/全分裂粗读出保持完全条件不变。

\begin{theorem}[最优部分可逆率的溢出恒等式]\label{thm:xi-time-part64b-optimal-partial-invertibility-overflow}
\leanverified{paper\_xi\_time\_part64b\_optimal\_partial\_invertibility\_overflow}
沿用定理 \ref{thm:fold-oracle-capacity-closed-form} 的记号。记
\[
d_m(x):=\bigl|\Fold_m^{-1}(x)\bigr|,
\qquad
T:=2^B,
\]
\[
C_m(B):=\mathcal C_m(B)=\sum_{x\in X_m}\min\bigl(d_m(x),T\bigr),
\qquad
\mathrm{Succ}_m(B):=\frac{C_m(B)}{2^m}.
\]
则有严格恒等式
\[
\boxed{\;
1-\mathrm{Succ}_m(B)
=
2^{-m}\sum_{x\in X_m}\bigl(d_m(x)-T\bigr)_+.\;}
\]
换言之，全部不可逆损失恰等于超过阈值 \(T\) 的纤维超额体积之归一化总和。
\end{theorem}
\begin{proof}
对任意整数 \(d\ge 0\)，都有逐点恒等式
\[
\min(d,T)=d-(d-T)_+.
\]
将其代入容量闭式得
\[
C_m(B)
=
\sum_{x\in X_m}\Bigl[d_m(x)-\bigl(d_m(x)-T\bigr)_+\Bigr].
\]
又
\[
\sum_{x\in X_m}d_m(x)=|\Omega_m|=2^m,
\]
故
\[
C_m(B)=2^m-\sum_{x\in X_m}\bigl(d_m(x)-T\bigr)_+.
\]
两边同除以 \(2^m\)，便得到
\[
1-\mathrm{Succ}_m(B)
=
2^{-m}\sum_{x\in X_m}\bigl(d_m(x)-T\bigr)_+.
\]
证毕。
\end{proof}

\begin{corollary}[近完美读出的尖锐必要条件]\label{cor:xi-time-part64b-near-perfect-readout-sharp-necessary-condition}
沿用定理 \ref{thm:xi-time-part64b-optimal-partial-invertibility-overflow} 的记号，并记
\[
M_m:=\max_{x\in X_m}d_m(x).
\]
若对某个 \(\varepsilon\ge 0\) 有
\[
\mathrm{Succ}_m(B)\ge 1-\varepsilon,
\]
则必有
\[
\boxed{\;
2^B\ge M_m-\varepsilon\,2^m.\;}
\]
特别地：
\begin{enumerate}
  \item 当 \(\varepsilon=0\) 时，
  \[
  2^B\ge M_m
  \]
  是零误差全可逆的必要条件；
  \item 若 \(2^B<M_m\)，则必有严格损失下界
  \[
  \boxed{\;
  1-\mathrm{Succ}_m(B)\ge \frac{M_m-2^B}{2^m}.\;}
  \]
\end{enumerate}
\end{corollary}
\begin{proof}
由定理 \ref{thm:xi-time-part64b-optimal-partial-invertibility-overflow}，
\[
1-\mathrm{Succ}_m(B)
=
2^{-m}\sum_{x\in X_m}\bigl(d_m(x)-2^B\bigr)_+.
\]
由于对某个达到最大值的 \(x_\ast\in X_m\) 有 \(d_m(x_\ast)=M_m\)，故
\[
1-\mathrm{Succ}_m(B)
\ge
2^{-m}\bigl(M_m-2^B\bigr)_+.
\]
若 \(\mathrm{Succ}_m(B)\ge 1-\varepsilon\)，则
\[
\varepsilon\ge \frac{\bigl(M_m-2^B\bigr)_+}{2^m}.
\]
若 \(M_m\le 2^B\)，所示结论显然成立；若 \(M_m>2^B\)，则上式化为
\[
M_m-2^B\le \varepsilon\,2^m,
\]
即
\[
2^B\ge M_m-\varepsilon\,2^m.
\]
其余两条断言分别对应 \(\varepsilon=0\) 与 \(M_m>2^B\) 的特例。证毕。
\end{proof}

\begin{theorem}[预言机容量的 Neyman--Pearson 对偶：推前分布与稳定层均匀基线的精确差异表达]\label{thm:xi-time-part64b-oracle-capacity-neyman-pearson-duality}
沿用定理 \ref{thm:fold-oracle-capacity-closed-form} 的记号。定义推前分布与稳定层均匀基线
\[
\mu_m(x):=\frac{d_m(x)}{2^m},
\qquad
u_m(x):=\frac{1}{|X_m|}
\qquad (x\in X_m),
\]
并令阈值 \(T:=2^B\)。再定义差异泛函
\[
\mathfrak D_m(B):=
\sup_{A\subseteq X_m}
\Bigl(
\mu_m(A)-2^{B-m}|X_m|\,u_m(A)
\Bigr).
\]
则有严格恒等式
\[
\mathrm{Succ}_m(B)=1-\mathfrak D_m(B),
\]
并且上确界恰由阈值集
\[
A_B^\star:=\{x\in X_m:\ d_m(x)>2^B\}
\]
取得。
\end{theorem}
\begin{proof}
由于
\[
|X_m|\,u_m(A)=|A|,
\]
故对任意 \(A\subseteq X_m\)，有
\[
\mu_m(A)-2^{B-m}|X_m|\,u_m(A)
=
\sum_{x\in A}\left(\frac{d_m(x)}{2^m}-\frac{2^B}{2^m}\right)
=
\frac{1}{2^m}\sum_{x\in A}\bigl(d_m(x)-2^B\bigr).
\]
上式在且仅在取遍所有正项时达到最大，因此
\[
\mathfrak D_m(B)
=
\frac{1}{2^m}
\sum_{x:\,d_m(x)>2^B}\bigl(d_m(x)-2^B\bigr),
\]
且最大集正是 \(A_B^\star\)。

另一方面，由定理 \ref{thm:xi-time-part64b-optimal-partial-invertibility-overflow} 的逐点恒等式
\[
\min(d,2^B)=d-(d-2^B)_+
\]
可得
\[
C_m(B)
=
\sum_{x\in X_m}\min\bigl(d_m(x),2^B\bigr)
=
2^m-\sum_{x\in X_m}\bigl(d_m(x)-2^B\bigr)_+.
\]
因此
\[
\mathfrak D_m(B)
=
\frac{2^m-C_m(B)}{2^m}
=
1-\mathrm{Succ}_m(B).
\]
证毕。
\end{proof}

\begin{theorem}[幂和诱导的非渐近失败率指数界与可计算预算下界]\label{thm:xi-time-part64b-power-sum-failure-bound-budget-lower}
对任意 \(q>1\)，定义幂和
\[
S_q(m):=\sum_{x\in X_m}d_m(x)^q.
\]
则对任意实预算 \(B\in\mathbb R\) 都有失败率上界
\[
1-\mathrm{Succ}_m(B)
\le
\frac{S_q(m)}{2^m\,2^{(q-1)B}}.
\]
因此，若希望
\[
\mathrm{Succ}_m(B)\ge 1-\varepsilon
\qquad (\varepsilon\in(0,1)),
\]
则充分条件为
\[
B\ge
\frac{\log_2 S_q(m)-m-\log_2\varepsilon}{q-1}.
\]
对右端再对 \(q>1\) 取下确界，便得到完全由矩谱驱动的最强可计算预算下界。
\end{theorem}
\begin{proof}
记
\[
T:=2^B.
\]
由定理 \ref{thm:xi-time-part64b-optimal-partial-invertibility-overflow}，
\[
2^m-C_m(B)=\sum_{x\in X_m}(d_m(x)-T)_+.
\]
对任意 \(d\ge 0\) 与 \(q>1\)，成立初等不等式
\[
(d-T)_+\le \frac{d^q}{T^{q-1}}.
\]
确实，若 \(d\le T\)，左端为 \(0\)；若 \(d>T\)，则
\[
d-T\le d\le d\left(\frac{d}{T}\right)^{q-1}=\frac{d^q}{T^{q-1}}.
\]
逐项求和即得
\[
2^m-C_m(B)
\le
\frac{1}{T^{q-1}}\sum_{x\in X_m}d_m(x)^q
=
\frac{S_q(m)}{2^{(q-1)B}}.
\]
两边同除以 \(2^m\)，便有
\[
1-\mathrm{Succ}_m(B)\le \frac{S_q(m)}{2^m\,2^{(q-1)B}}.
\]
若右端不超过 \(\varepsilon\)，则 \(\mathrm{Succ}_m(B)\ge 1-\varepsilon\)。对上式整理即得
\[
B\ge
\frac{\log_2 S_q(m)-m-\log_2\varepsilon}{q-1}.
\]
证毕。
\end{proof}

\begin{theorem}[折叠信息几何的二点塌缩原理]\label{thm:xi-time-part64b-fold-information-geometry-two-point-collapse}
沿用命题 \ref{prop:fold-bin-f-divergence-collapse} 的记号。定义
\[
a:=\frac{\varphi^2}{\sqrt5},
\qquad
b:=\frac{\varphi}{\sqrt5},
\qquad
\nu_\varphi:=\frac1\varphi\,\delta_a+\frac1{\varphi^2}\,\delta_b.
\]
则对任意凸函数 \(f:(0,\infty)\to\RR\) 满足 \(f(1)=0\)，其对应的 Csisz\'ar \(f\)-散度满足
\[
\boxed{\;
\lim_{m\to\infty}D_f(\mu_m\Vert u_m)
=
\int_{(0,\infty)}f(t)\,\dd\nu_\varphi(t)
=
\frac1\varphi\,f\!\left(\frac{\varphi^2}{\sqrt5}\right)
+\frac1{\varphi^2}\,f\!\left(\frac{\varphi}{\sqrt5}\right).\;}
\]
因此，\((\mu_m,u_m)\) 的全部渐近 \(f\)-信息几何并不依赖于完整纤维谱，而完全塌缩到两个似然比原子 \(a,b\) 上；换言之，折叠多重度在信息论极限中仅保留二相可辨结构。
\end{theorem}
\begin{proof}
命题 \ref{prop:fold-bin-f-divergence-collapse} 已直接给出
\[
D_f(\mu_m\Vert u_m)
=
\frac1\varphi\,f\!\left(\frac{\varphi^2}{\sqrt5}\right)
+\frac1{\varphi^2}\,f\!\left(\frac{\varphi}{\sqrt5}\right)
+O_f\bigl((\varphi/2)^m\bigr),
\]
故所示极限公式立即成立。

另一方面，同一命题的证明还给出一致展开
\[
\frac{\mu_m(w)}{u_m(w)}
=
\frac{\varphi^{2-w_m}}{\sqrt5}+O\bigl((\varphi/2)^m\bigr)
\qquad (w\in X_m),
\]
因而极限似然比仅可能落在
\[
\frac{\varphi^2}{\sqrt5}=a,
\qquad
\frac{\varphi}{\sqrt5}=b
\]
两点上。再由
\[
u_m(w_m=0)=\frac{F_{m+1}}{F_{m+2}}\longrightarrow \frac1\varphi,
\qquad
u_m(w_m=1)=\frac{F_m}{F_{m+2}}\longrightarrow \frac1{\varphi^2},
\]
可知极限似然比测度恰为
\[
\nu_\varphi=\frac1\varphi\,\delta_a+\frac1{\varphi^2}\,\delta_b.
\]
于是任意由似然比驱动并在该一致收敛下稳定的 \(f\)-统计，都只能因子化通过 \(\nu_\varphi\)。证毕。
\end{proof}

\begin{theorem}[局域可加性与非局域可乘性的严格偏振]\label{thm:xi-time-part64b-local-additivity-nonlocal-multiplicativity-polarization}
沿用定义 \ref{def:stable-add}、\ref{def:stable-mul}，定理 \ref{thm:normalize-then-truncate-functorial}，注 \ref{rem:zeckendorf-normalize-fst} 以及命题 \ref{prop:killo-implementation-vs-structural-half-circle-dimension} 的记号。则在本文构造的 Zeckendorf/Ostrowski 数字层中，规范化与加法具有纯数字串实现：可在有限状态换能器或局部重写框架内完成，并在分辨率口径上线性、在比特口径上为多项式时间。

然而，若坚持“乘法 \(\mapsto\) 有限寄存器加法”的严格同态语义，则不存在任何有限秩压缩：既不存在单射幺半群同态
\[
\NN_{\times}\hookrightarrow (\NN^k,+)
\qquad(k<\infty),
\]
亦不存在有限生成交换群 \(L\) 使得
\[
\Phi(ab)=\Phi(a)+\Phi(b),
\qquad
\Phi:\NN_{\times}\to \ZZ\oplus L
\]
为单射。

并且
\[
\boxed{\;
\cdim^{\mathrm{hom}}_{+}(\NN_{\times})=+\infty,
\qquad
\cdim^{\mathrm{impl}}(\NN_{\times})=\frac12.\;}
\]
因此，该体系中的精确算术必然分裂为两层：一层是局域、有限状态、纯数字的可加层；另一层是任何有限加法账本都无法吸收的外置乘法层。
\end{theorem}
\begin{proof}
先看数字层的局域实现。定理 \ref{thm:normalize-then-truncate-functorial} 已把一般 Ostrowski 数字层的“先规范形再截断”固定为定义性的局部机制，而注 \ref{rem:zeckendorf-normalize-fst} 进一步指出，在黄金分支上 Zeckendorf/Ostrowski 规范化与在线加法可由纯数字串层的有限状态换能器或局部重写实现，从而在分辨率口径上具有线性复杂度。另一方面，定理 \ref{thm:xi-stable-semiring-ptime-biinterpretability} 又给出值层回拉实现的多项式时间界：稳定地址层与自然数半环在 PTIME 内双向可解释，故在比特口径上，加法、规范化以及经由值层回拉的乘法都具有多项式时间实现。

再看“乘法 \(\mapsto\) 有限寄存器加法”的严格同态化是否可能。推论 \ref{cor:killo-no-finite-additive-register-linearization} 断言：对任意有限 \(k\)，不存在单射幺半群同态
\[
\NN_{\times}\hookrightarrow (\NN^k,+).
\]
定理 \ref{thm:killo-godel-compression-not-finite-rank-homologizable} 又断言：若 \(L\) 有限生成，则不存在单射
\[
\Phi:\NN_{\times}\to \ZZ\oplus L
\]
满足
\[
\Phi(ab)=\Phi(a)+\Phi(b).
\]
因此，一切有限秩加法账本上的严格乘法同态化都被排除。

最后，由命题 \ref{prop:killo-implementation-vs-structural-half-circle-dimension}，
\[
\cdim^{\mathrm{hom}}_{+}(\NN_{\times})=+\infty,
\qquad
\cdim^{\mathrm{impl}}(\NN_{\times})=\frac12.
\]
这正表明：作为集合，\(\NN_{\times}\) 仍可嵌入单寄存器实现层；但作为乘法幺半群，其结构层同态半圆维不可有限化。故局域数字可加性与乘法外置性形成严格偏振。证毕。
\end{proof}

\begin{theorem}[三轴最小窗的非纯化定理]\label{thm:xi-time-part64b-three-axis-minimum-window-non-purification}
设 \(m_\star\) 为命题 \ref{prop:crt-235-min-depth} 给出的三轴窗最小分辨率。则
\[
\boxed{\;
m_\star=58,
\qquad
m_\star+2=60.\;}
\]
其对应底层 Fibonacci 模数满足
\[
\boxed{\;
F_{60}
=
2^4\cdot 3^2\cdot 5\cdot 11\cdot 31\cdot 41\cdot 61\cdot 2521.\;}
\]
因此，第一次三轴共现并非仅引入目标三条算术方向 \(2,3,5\)；它被严格迫使同时携带五条额外素轴 \(11,31,41,61,2521\)。换言之，三轴统一在最小可见层上即不可纯化。
\end{theorem}
\begin{proof}
命题 \ref{prop:crt-235-min-depth} 已给出最小三轴窗的精确位置
\[
m_\star=58,
\qquad
m_\star+2=60.
\]
另一方面，注 \ref{rem:crt-235-m58-factorization} 已给出对应 Fibonacci 模数的精确分解
\[
F_{60}
=
1548008755920
=
2^4\cdot 3^2\cdot 5\cdot 11\cdot 31\cdot 41\cdot 61\cdot 2521.
\]
故在最小三轴层中，除目标三轴 \(2,3,5\) 外，还不可避免地并入五条额外素轴
\[
11,\ 31,\ 41,\ 61,\ 2521.
\]
这就说明三轴最小窗并不能被“纯三轴”算术所单独承载，而天然带有额外的剩余素数方向。证毕。
\end{proof}

\begin{theorem}[十次指纹对 Lee--Yang 粗读出的条件不变性]\label{thm:xi-time-part64b-p10-fingerprint-leyang-coarse-readout-invariance}
记
\[
A:=\{\text{\(P_{10}\bmod p\) 含一次因子}\},
\]
并定义 Lee--Yang 侧三个粗事件
\[
E_1:=\{\chi_{\mathrm{LY}}(p)=+1\},
\qquad
E_2:=\{\chi_{\mathrm{LY}}(p)=-1\},
\qquad
E_3:=B_{\mathrm{root}}:=\{\text{\(g\bmod p\) 至少有一根}\}.
\]
再记
\[
B_{\mathrm{split}}:=\{\text{\(g\bmod p\) 全分裂}\}.
\]
则自然密度满足
\[
\boxed{\;
\delta(A\mid E_1)
=
\delta(A\mid E_2)
=
\delta(A\mid E_3)
=
\delta(A)
=
\frac{28319}{44800}.\;}
\]
并且
\[
\boxed{\;
\delta(A\mid B_{\mathrm{split}})
=
\frac{28319}{44800}.\;}
\]
因此，十次模分解指纹与 Lee--Yang 的符号读出、存在根读出及全分裂读出在统计上完全正交：任何此类粗 Lee--Yang 观测都不改变十次线性因子事件的密度。
\end{theorem}
\begin{proof}
由命题 \ref{prop:fold-gauge-anomaly-P10-leyang-linear-disjointness} 与推论
\ref{cor:fold-gauge-anomaly-P10-leyang-chebotarev-product-law}，任意只依赖 \(S_{10}\) 因子的共轭类事件与任意只依赖 \(S_3\) 因子的共轭类事件，其联合自然密度都等于边际密度之积。

事件 \(A\) 只依赖 \(P_{10}\) 分裂域的 \(S_{10}\) 因子；而
\[
E_1,\ E_2,\ B_{\mathrm{root}},\ B_{\mathrm{split}}
\]
都只依赖 Lee--Yang 三次分裂域的 \(S_3\) 因子。故对任意
\[
E\in\{E_1,E_2,B_{\mathrm{root}},B_{\mathrm{split}}\}
\]
都成立
\[
\delta(A\cap E)=\delta(A)\,\delta(E).
\]

再由推论 \ref{cor:fold-gauge-anomaly-P10-modp-splitting-densities}，
\[
\delta(A)=\frac{28319}{44800}.
\]
另一方面，在 \(S_3\) 中，偶置换与奇置换各占 \(3\) 个元素，故
\[
\delta(E_1)=\delta(E_2)=\frac12.
\]
又由结论 \ref{con:xi-terminal-zm-leyang-cubic-modp-splitting-density}，\(g\bmod p\) 的分裂型与 \(S_3\) 的循环型一一对应：恒等类对应全分裂，换位类对应恰有一根，三循环类对应不可约。三类大小分别为 \(1,3,2\)，故
\[
\delta(B_{\mathrm{split}})=\frac16,
\qquad
\delta(B_{\mathrm{root}})=\frac{1+3}{6}=\frac23.
\]

由于上述四个条件事件的密度均为正，分别相除即可得到
\[
\delta(A\mid E_1)=\delta(A),
\qquad
\delta(A\mid E_2)=\delta(A),
\]
\[
\delta(A\mid B_{\mathrm{root}})=\delta(A),
\qquad
\delta(A\mid B_{\mathrm{split}})=\delta(A).
\]
将 \(\delta(A)=28319/44800\) 代回，即得全部公式。证毕。
\end{proof}

\noindent
上述六条结果把容量亏损、极限信息几何、数字层实现、三轴最小窗与 Chebotarev 条件化五条原先分立的链进一步压缩为同一组有限证书：部分可逆误差由纤维超额尾量完全记录，渐近似然比只剩两原子，乘法的外置性由同态半圆维无穷障碍强制，最小三轴窗天然携带额外素轴，而十次指纹对 Lee--Yang 粗读出保持严格条件不变。于是，恢复论、信息论、实现论、模窗算术与分裂型统计之间的横向耦合便获得了统一的有限闭合。

\endinput

\paragraph{stop-loss 容量曲线、幂散度单标量刚性、平衡剖面与最小扇区 Witten 电荷}
在临界预言机容量双折线、二原子极限似然比、\(\chi^2\)-型极限常数对幂散度族的坐标化、折叠多重度的 Robin--Hood 调平机制以及审计偶窗最小退化扇区的同伦相位都已建立之后，还可把这些链条进一步压缩为下列七条彼此衔接的对象层后果。

\begin{theorem}[临界 oracle 曲线恰为极限似然比的 stop-loss 变换]\label{thm:xi-time-part64ba-critical-oracle-stoploss-transform}
沿用定理 \ref{thm:xi-foldbin-dyadic-capacity-critical-window-limit} 与定理 \ref{thm:xi-time-part64b-fold-information-geometry-two-point-collapse} 的记号。定义
\[
r_0:=\frac{\varphi^2}{\sqrt5},
\qquad
r_1:=\frac{\varphi}{\sqrt5},
\]
并令 \(R_\infty\) 为二原子随机变量
\[
R_\infty=
\begin{cases}
r_0,& \text{概率 } \varphi^{-1},\\[1mm]
r_1,& \text{概率 } \varphi^{-2}.
\end{cases}
\]
再记临界尺度下的最优成功率极限曲线为
\[
S(\tau):=\mathcal S_{\mathrm{bin}}(\tau)
\qquad (\tau\ge 0),
\]
即
\[
S(\tau)=
\begin{cases}
\dfrac{\varphi^2}{\sqrt5}\,\tau,& 0\le \tau\le \varphi^{-1},\\[3mm]
\dfrac{1}{\varphi\sqrt5}+\dfrac{\varphi}{\sqrt5}\,\tau,& \varphi^{-1}\le \tau\le 1,\\[3mm]
1,& \tau\ge 1.
\end{cases}
\]
则对一切 \(\tau\ge 0\)，成立精确恒等式
\[
\boxed{\;
S(\tau)=\EE\!\left[\min\{R_\infty,\tau r_0\}\right]
=
\int_{0}^{\tau r_0}\PP(R_\infty>t)\,\dd t.\;}
\]
特别地，当 \(\tau\notin\{\varphi^{-1},1\}\) 时，
\[
\boxed{\;
S'(\tau)=r_0\,\PP(R_\infty>\tau r_0).\;}
\]
因此，临界 oracle 曲线的全部折点恰由极限似然比尾分布的 stop-loss 变换完全产生。
\end{theorem}
\begin{proof}
记
\[
\lambda:=\tau r_0.
\]
分三段讨论。

若 \(0\le \tau\le \varphi^{-1}\)，则
\[
\lambda\le r_1,
\]
故
\[
\EE[\min\{R_\infty,\lambda\}]=\lambda=\frac{\varphi^2}{\sqrt5}\tau.
\]

若 \(\varphi^{-1}\le \tau\le 1\)，则
\[
r_1\le \lambda\le r_0,
\]
于是
\[
\EE[\min\{R_\infty,\lambda\}]
=
\varphi^{-2}r_1+\varphi^{-1}\lambda
=
\frac{1}{\varphi\sqrt5}+\frac{\varphi}{\sqrt5}\tau.
\]

若 \(\tau\ge 1\)，则
\[
\lambda\ge r_0,
\]
故
\[
\EE[\min\{R_\infty,\lambda\}]
=
\EE[R_\infty]
=
\varphi^{-1}r_0+\varphi^{-2}r_1
=1.
\]
这与 \(S(\tau)\) 的三段闭式完全一致。

第二个等式是 stop-loss 变换的标准层蛋糕恒等式
\[
\EE[\min\{X,\lambda\}]=\int_0^\lambda \PP(X>t)\,\dd t
\qquad (X\ge 0),
\]
应用于 \(X=R_\infty\) 即得。对 \(\tau\) 在非折点处求导，链式法则给出
\[
S'(\tau)=r_0\,\PP(R_\infty>\tau r_0).
\]
证毕。
\end{proof}

\begin{theorem}[单个 \texorpdfstring{$\chi^2$}{chi2} 型极限常数决定整个 power-divergence 谱，并满足二阶常系数微分方程]\label{thm:xi-time-part64ba-chi2-rigid-powerdivergence-ode}
\leanverified{paper\_xi\_time\_part64ba\_chi2\_rigid\_powerdivergence\_ode}
记
\[
\chi:=D_{f_2}^\infty.
\]
则对每个 \(\alpha>1\)，全部极限 power-divergence 常数都可由 \(\chi\) 显式恢复：
\[
\boxed{\;
D_{f_\alpha,\infty}
=
5^{-\alpha/2}
\left[
\left(\frac{5\chi+3}{2}\right)^{2\alpha-1}
+
\left(\frac{5\chi+3}{2}\right)^{\alpha-2}
\right]-1.\;}
\]
若再令
\[
E(\alpha):=D_{f_\alpha,\infty}+1,
\]
\[
c_\chi:=\log\frac{(5\chi+3)^2}{4\sqrt5},
\qquad
d_\chi:=\log\frac{5\chi+3}{2\sqrt5},
\]
则 \(E\) 满足精确线性常微分方程
\[
\boxed{\;
E''(\alpha)-(c_\chi+d_\chi)E'(\alpha)+c_\chi d_\chi\,E(\alpha)=0.\;}
\]
因此，整个极限 power-divergence 族并非自由函数，而是由单个二阶统计常数 \(\chi\) 刚性锁定的二维指数簇。
\end{theorem}
\begin{proof}
由结论 \ref{con:xi-time-part9l-chi2-generates-power-divergence-and-cumulants}，
\[
D_{f_2}^\infty=\frac{2\varphi-3}{5},
\qquad
\varphi=\frac{5D_{f_2}^\infty+3}{2}.
\]
代入 \(\chi=D_{f_2}^\infty\) 后得到
\[
\varphi=\frac{5\chi+3}{2}.
\]
同一结论又给出一般公式
\[
D_{f_\alpha}^\infty
=
\frac{1}{5^{\alpha/2}}\bigl(\varphi^{2\alpha-1}+\varphi^{\alpha-2}\bigr)-1.
\]
把 \(\varphi=(5\chi+3)/2\) 代入，即得第一条闭式。

现看 \(E(\alpha)=D_{f_\alpha,\infty}+1\)。由上式，
\[
E(\alpha)
=
5^{-\alpha/2}\bigl(\varphi^{2\alpha-1}+\varphi^{\alpha-2}\bigr)
=
\varphi^{-1}\left(\frac{\varphi^2}{\sqrt5}\right)^\alpha
+
\varphi^{-2}\left(\frac{\varphi}{\sqrt5}\right)^\alpha.
\]
再用
\[
\varphi=\frac{5\chi+3}{2}
\]
重写两条指数率，便得
\[
\frac{\varphi^2}{\sqrt5}
=
\frac{(5\chi+3)^2}{4\sqrt5}
=
e^{c_\chi},
\qquad
\frac{\varphi}{\sqrt5}
=
\frac{5\chi+3}{2\sqrt5}
=
e^{d_\chi}.
\]
故
\[
E(\alpha)=A e^{c_\chi\alpha}+B e^{d_\chi\alpha},
\qquad
A=\varphi^{-1},
\qquad
B=\varphi^{-2}.
\]
任意两指数和都满足
\[
(\partial_\alpha-c_\chi)(\partial_\alpha-d_\chi)E=0,
\]
亦即
\[
E''-(c_\chi+d_\chi)E'+c_\chi d_\chi E=0.
\]
证毕。
\end{proof}

\begin{corollary}[黄金几何审计常数由同一个 \texorpdfstring{$\chi^2$}{chi2} 型极限标量完全恢复]\label{cor:xi-time-part64ba-chi2-recovers-golden-geometric-audits}
仍记 \(\chi=D_{f_2}^\infty\)。则黄金分支上既有 Wasserstein/星差异审计常数可全部写成 \(\chi\) 的有理分式：
\[
\boxed{\;
\lim_{\substack{n\to\infty\\ n\ \mathrm{even}}}F_n D_n^\ast
=
1+\frac{1}{5\chi+2}.\;}
\]
\[
\boxed{\;
\lim_{\substack{n\to\infty\\ n\ \mathrm{even}}}F_n T_n
=
\frac12+\frac{1}{2(5\chi+2)}.\;}
\]
\[
\boxed{\;
\lim_{\substack{n\to\infty\\ n\ \mathrm{odd}}}F_n T_n
=
\frac{17}{30}-\frac{1}{2(5\chi+2)}.\;}
\]
故同一个 \(\chi^2\)-型极限常数不仅决定信息散度族，而且刚性控制黄金分支的几何审计常数。
\end{corollary}
\begin{proof}
文稿在黄金分支上已给出
\[
\lim_{\substack{n\to\infty\\ n\ \mathrm{even}}}F_nT_n
=
\frac12+\frac{1}{2\sqrt5},
\]
\[
\lim_{\substack{n\to\infty\\ n\ \mathrm{odd}}}F_nT_n
=
\frac12-\frac{1}{2\sqrt5}+\frac{1}{15}
=
\frac{17}{30}-\frac{1}{2\sqrt5},
\]
并且偶数支上满足
\[
T_n=\frac12 D_n^\ast.
\]
另一方面，由结论 \ref{con:xi-time-part9l-chi2-generates-power-divergence-and-cumulants}，
\[
\sqrt5=5\chi+2.
\]
把该恒等式代入上面三条公式，便得到所示全部闭式。证毕。
\end{proof}

\begin{theorem}[每一级折叠多重度剖面都可由有限 Robin--Hood 链严格平衡化；平衡剖面位于其双随机轨道中]\label{thm:xi-time-part64ba-fold-multiplicity-majorization-balancing}
\leanverified{paper\_xi\_time\_part64ba\_fold\_multiplicity\_majorization\_balancing}
记
\[
F:=F_{m+2},
\qquad
M:=2^m,
\qquad
M=aF+\rho,
\qquad
0\le \rho<F.
\]
将折叠多重度剖面写为
\[
d^{(m)}:=(d_m(u))_{u\in \ZZ/(F\ZZ)}\in \ZZ_{\ge 0}^{F},
\]
并定义平衡剖面
\[
b^{(m)}:=
(\underbrace{a+1,\dots,a+1}_{\rho\text{ 个}},
\underbrace{a,\dots,a}_{F-\rho\text{ 个}}).
\]
则在按非增序重排后，有主序关系
\[
\boxed{\;
d^{(m),\downarrow}\succ b^{(m)}.\;}
\]
等价地，存在一个 \(F\times F\) 双随机矩阵 \(\Pi_m\) 使
\[
\boxed{\;
b^{(m)}=\Pi_m\,d^{(m)}.\;}
\]
更强地，可由有限次平衡化操作
\[
(v_{u_1},v_{u_2})\longmapsto (v_{u_1}-1,v_{u_2}+1)
\qquad (v_{u_1}\ge v_{u_2}+2)
\]
把 \(d^{(m)}\) 明确送到 \(b^{(m)}\)。
\end{theorem}
\begin{proof}
结论 \ref{con:xi-time-part9l-fold-fourier-energy-arithmetic-lower-bound} 的证明已经给出如下严格调平机制：若某个整数向量 \(v\in\ZZ_{\ge 0}^{F}\) 满足
\[
\sum_u v_u=M
\]
且存在指标 \(u_1,u_2\) 使
\[
v_{u_1}\ge v_{u_2}+2,
\]
则执行一步 Robin--Hood 操作
\[
(v_{u_1},v_{u_2})\mapsto (v_{u_1}-1,v_{u_2}+1)
\]
会严格降低平方偏差和，因此有限步之后必然终止于某个任意两坐标相差至多 \(1\) 的向量。

由于总和固定为
\[
M=aF+\rho,
\]
满足“任意两坐标相差至多 \(1\)” 的整数向量只能是恰有 \(\rho\) 个分量等于 \(a+1\)、其余 \(F-\rho\) 个分量等于 \(a\) 的排列，也就是 \(b^{(m)}\) 的一个置换。把这些分量按非增序排列，唯一得到
\[
b^{(m)}.
\]
因此 \(d^{(m)}\) 经过有限次 Robin--Hood 调平后确实到达 \(b^{(m)}\)。

而 Hardy--Littlewood--Pólya 判别表明：一个向量可经有限次 Robin--Hood 操作被送到另一个向量，当且仅当前者 majorize 后者。故
\[
d^{(m),\downarrow}\succ b^{(m)}.
\]
再由 Hardy--Littlewood--Pólya 定理与 Birkhoff--von Neumann 定理，主序关系等价于存在双随机矩阵 \(\Pi_m\) 使
\[
b^{(m)}=\Pi_m d^{(m)}.
\]
证毕。
\end{proof}

\begin{corollary}[平衡剖面给出全阶 Rényi 熵的统一上包络]\label{cor:xi-time-part64ba-balanced-profile-universal-renyi-envelope}
在定理 \ref{thm:xi-time-part64ba-fold-multiplicity-majorization-balancing} 的记号下，定义
\[
p_m(u):=\frac{d_m(u)}{2^m},
\qquad
p_m^{\mathrm{bal}}:=\frac{b^{(m)}}{2^m}.
\]
则对任意 \(\alpha>0\)，Rényi 熵满足
\[
\boxed{\;
H_\alpha(p_m)\le H_\alpha(p_m^{\mathrm{bal}}).\;}
\]
其中当 \(\alpha\neq 1\) 时，
\[
H_\alpha(p):=\frac{1}{1-\alpha}\log\sum_u p(u)^\alpha,
\]
而上包络显式为
\[
\boxed{\;
H_\alpha(p_m)
\le
\frac{1}{1-\alpha}
\log\!\left(
\frac{(F-\rho)a^\alpha+\rho(a+1)^\alpha}{(aF+\rho)^\alpha}
\right).\;}
\]
在 Shannon 极限 \(\alpha\to 1\) 下，
\[
\boxed{\;
H(p_m)\le
-\frac{(F-\rho)a}{M}\log\frac{a}{M}
-\frac{\rho(a+1)}{M}\log\frac{a+1}{M}.\;}
\]
并且上述各式取等当且仅当多重度剖面本身已经平衡。
\end{corollary}
\begin{proof}
由定理 \ref{thm:xi-time-part64ba-fold-multiplicity-majorization-balancing}，
\[
d^{(m),\downarrow}\succ b^{(m)}.
\]
归一化后仍有
\[
p_m\succ p_m^{\mathrm{bal}}.
\]
而 Rényi 熵对一切 \(\alpha>0\) 都是 Schur--凹函数，故
\[
H_\alpha(p_m)\le H_\alpha(p_m^{\mathrm{bal}}).
\]
把平衡剖面
\[
p_m^{\mathrm{bal}}=
\left(
\underbrace{\frac{a+1}{M},\dots,\frac{a+1}{M}}_{\rho\text{ 个}},
\underbrace{\frac{a}{M},\dots,\frac{a}{M}}_{F-\rho\text{ 个}}
\right)
\]
直接代入 Rényi 熵定义，便得到
\[
H_\alpha(p_m^{\mathrm{bal}})
=
\frac{1}{1-\alpha}
\log\!\left(
\frac{(F-\rho)a^\alpha+\rho(a+1)^\alpha}{M^\alpha}
\right).
\]
其中 \(M=aF+\rho\)，故即为题述上界。令 \(\alpha\to 1\) 即得 Shannon 形式。

最后，由定理 \ref{thm:xi-time-part64ba-fold-multiplicity-majorization-balancing} 中 Robin--Hood 降平过程的严格性可知，若 \(d^{(m)}\) 不是平衡剖面，则主序关系严格；而 Rényi 熵在此情形下严格 Schur--凹，故等号仅当 \(d^{(m)}\) 本身已平衡时成立。证毕。
\end{proof}

\begin{theorem}[已审计偶窗的最小退化扇区携带一条显式可求的总 Witten 电荷]\label{thm:xi-time-part64ba-audited-even-minsector-total-witten-charge}
对审计偶窗
\[
m\in\{6,8,10,12\},
\]
定义最小退化扇区
\[
S_{m,\min}:=\{x\in X_m:\ d_m(x)=d_{\min}(m)\},
\]
以及交错 Witten 指数
\[
W_x:=\sum_{S\in \mathsf{Ind}(\Gamma_x)}(-1)^{|S|}.
\]
再定义总 Witten 电荷
\[
\mathcal W_{\min}(m):=\sum_{x\in S_{m,\min}}W_x.
\]
则有精确公式
\[
\boxed{\;
\mathcal W_{\min}(m)=
\begin{cases}
0,& m=6,12,\\[1mm]
-F_m,& m=8,10.
\end{cases}\;}
\]
特别地，
\[
\boxed{\;
\frac{\mathcal W_{\min}(m)}{|S_{m,\min}|}=
\begin{cases}
0,& m=6,12,\\[1mm]
-1,& m=8,10.
\end{cases}\;}
\]
因此，最小退化扇区在审计偶窗上并不只是经历“可缩/不可缩”二分，而且携带一条从 \(0\) 到 \(-F_m\) 的整扇区总电荷跃迁。
\end{theorem}
\begin{proof}
由定理 \ref{thm:fold-bin-min-degeneracy-fib-audited}，
\[
d_{\min}(6)=2,\qquad d_{\min}(8)=3,\qquad d_{\min}(10)=5,\qquad d_{\min}(12)=8,
\]
且
\[
|S_{m,\min}|=F_m.
\]

另一方面，附录定理 \ref{thm:derived-fiber-indcomplex-alternating-witten-parity} 给出
\[
W_x=
\begin{cases}
0,& d_m(x)\ \text{为偶},\\[1mm]
(-1)^{\tau(x)},& d_m(x)\ \text{为奇}.
\end{cases}
\]
故当 \(m=6,12\) 时，\(d_{\min}(m)=2,8\) 为偶，于是对每个
\[
x\in S_{m,\min}
\]
都有
\[
W_x=0.
\]
因此
\[
\mathcal W_{\min}(m)=0.
\]

当 \(m=8,10\) 时，\(d_{\min}(m)=3,5\) 为奇。又由定理 \ref{thm:xi-audited-even-minsector-homotopy-mod6-phase} 的证明，在这两种情形中都必有
\[
\tau(x)=1
\qquad (x\in S_{m,\min}),
\]
因为
\[
3=F_4,\qquad 5=F_5
\]
分别只对应单一非平凡 Fibonacci 因子。于是
\[
W_x=(-1)^{\tau(x)}=-1
\qquad (x\in S_{m,\min}).
\]
故
\[
\mathcal W_{\min}(m)= -|S_{m,\min}|=-F_m.
\]
最后再除以 \(|S_{m,\min}|=F_m\)，便得到平均电荷公式。证毕。
\end{proof}

\begin{corollary}[最小退化扇区的非平凡拓扑质量密度在审计偶窗上仅取两值 \(0\) 与 \(1\)]\label{cor:xi-time-part64ba-audited-even-minsector-topological-mass-density}
定义
\[
\mathfrak m_{\min}(m)
:=
\frac{1}{|S_{m,\min}|}\sum_{x\in S_{m,\min}}|W_x|.
\]
则对审计偶窗 \(m\in\{6,8,10,12\}\) 有
\[
\boxed{\;
\mathfrak m_{\min}(m)=
\begin{cases}
0,& m=6,12,\\[1mm]
1,& m=8,10.
\end{cases}\;}
\]
换言之，在这些审计偶窗上，最小退化扇区要么整体拓扑平凡，要么每条纤维都携带单位绝对 Witten 电荷；不存在混合相。
\end{corollary}
\begin{proof}
由附录定理 \ref{thm:derived-fiber-indcomplex-alternating-witten-parity}，
\[
|W_x|=
\begin{cases}
0,& d_m(x)\ \text{为偶},\\[1mm]
1,& d_m(x)\ \text{为奇}.
\end{cases}
\]
再结合定理 \ref{thm:xi-time-part64ba-audited-even-minsector-total-witten-charge} 中对四个审计偶窗的逐窗判定，便立即得到所示两值结论。证毕。
\end{proof}

\noindent
由此，临界 oracle 曲线、幂散度谱、折叠多重度剖面与最小退化扇区的拓扑相位不再是彼此平行的局部现象：一方面，临界容量双折线正是极限似然比的 stop-loss 影像，而整个 power-divergence 族又由单个 \(\chi^2\)-型极限常数刚性锁定；另一方面，Robin--Hood 调平把任意多重度剖面推向唯一平衡端点，并使这一平衡端点同时成为全阶 Rényi 熵的统一上包络；最后，审计偶窗上的最小退化扇区不只在同伦型上区分为可缩相与球面相，而且在交错 Witten 电荷层面进一步裂解为零电荷相与满密度单位电荷相。于是，统计 stop-loss、二阶信息坐标、双随机平衡化与拓扑电荷跃迁最终被压缩进同一条可审计主线。

\endinput

\paragraph{双分歧场、同余三次专化与 Lee--Yang/\texorpdfstring{$19$}{19} 次层的乘法化 Frobenius 统计}
前述链条已经分别建立了十次分歧域 \(L_{10}\)、九次分歧域 \(L_9\)、同余三次专化域 \(M_t\)、Lee--Yang 三次域与十九次指数域 \(L_H\) 的满对称伽罗瓦性及其两两独立性。
下面把这些分散的算术接口进一步压缩为两条互不重叠而可交叉编译的统计主轴：一条是
\[
S_{10}\times S_9\times S_3
\]
上的三重分解型与符号盒规律，另一条是
\[
S_{10}\times S_3\times S_{19}
\]
上的 Lee--Yang 根出现事件与十九次固定点统计的完全张量化。

\begin{theorem}[双分歧场与同余三次专化场的三重直积伽罗瓦群]\label{thm:xi-time-part64bb-p10-p9-trigonal-threefold-galois-independence}
沿用定理 \ref{thm:fold-gauge-anomaly-P10-P9-linear-disjointness} 与定理
\ref{thm:fold-gauge-anomaly-second-trigonal-ht-specialization-s3-disjoint} 的记号。
设 \(L_{10}/\QQ\) 与 \(L_9/\QQ\) 分别为十次分歧多项式 \(P_{10}\) 与九次分歧多项式 \(P_9\) 的分裂域。
对整数
\[
t\ge 5,
\qquad
t\not\equiv 1\pmod 3,
\]
令 \(h_t(u)\in\QQ[u]\) 为三次专化多项式，其分裂域记为 \(M_t\)。
则
\[
M_t\cap L_{10}L_9=\QQ,
\qquad
\Gal(L_{10}L_9M_t/\QQ)\cong S_{10}\times S_9\times S_3.
\]
特别地，
\[
[L_{10}L_9M_t:\QQ]=10!\cdot 9!\cdot 6.
\]
\end{theorem}
\begin{proof}
由定理 \ref{thm:fold-gauge-anomaly-P10-P9-linear-disjointness}，
\[
L_{10}\cap L_9=\QQ,
\qquad
\Gal(L_{10}L_9/\QQ)\cong S_{10}\times S_9.
\]
而由定理 \ref{thm:fold-gauge-anomaly-second-trigonal-ht-specialization-s3-disjoint}，
\[
\Gal(M_t/\QQ)\cong S_3,
\qquad
M_t\cap L_{10}=\QQ,
\qquad
M_t\cap L_9=\QQ.
\]
故只须排除 \(M_t\) 与合成域 \(L_{10}L_9\) 之间出现新的公共正规子域。

记
\[
E:=M_t\cap L_{10}L_9.
\]
由于两边都是 \(\QQ\) 上伽罗瓦扩张，故 \(E/\QQ\) 亦为伽罗瓦扩张。
又因为 \(\Gal(M_t/\QQ)\cong S_3\)，\(M_t\) 的唯一非平凡真正规子域只能是其判别式二次域
\[
K_t:=\QQ\!\left(\sqrt{\Disc_u(h_t)}\right).
\]
同一定理已给出
\[
\Disc_u(h_t)=-(t-1)P_9(t)<0,
\]
故 \(K_t\) 为虚二次域。

另一方面，
\[
\Gal(L_{10}L_9/\QQ)\cong S_{10}\times S_9
\]
的交换化为
\[
(S_{10}\times S_9)^{\mathrm{ab}}\cong C_2\times C_2.
\]
因此 \(L_{10}L_9\) 的二次正规子域恰对应三条非平凡符号特征，分别为
\[
\QQ(\sqrt{66269989}),
\qquad
\QQ(\sqrt{33605193}),
\qquad
\QQ(\sqrt{66269989\cdot 33605193}),
\]
它们全都是实二次域。
于是 \(K_t\) 不可能嵌入 \(L_{10}L_9\)，从而 \(E\neq K_t\)。
故只能 \(E=\QQ\)。

交域平凡且两边皆为伽罗瓦扩张，遂得到线性互素与直积型伽罗瓦群同构
\[
\Gal(L_{10}L_9M_t/\QQ)\cong \Gal(L_{10}L_9/\QQ)\times \Gal(M_t/\QQ)
\cong S_{10}\times S_9\times S_3.
\]
次数公式随之立即成立。证毕。
\end{proof}

\begin{corollary}[三重分解型的完全乘法密度与 \texorpdfstring{$1/270$}{1/270} 定律]\label{cor:xi-time-part64bb-p10-p9-trigonal-chebotarev-1over270}
在定理 \ref{thm:xi-time-part64bb-p10-p9-trigonal-threefold-galois-independence} 的设定下，
对任意分拆
\[
\lambda\vdash 10,
\qquad
\mu\vdash 9,
\qquad
\nu\vdash 3,
\]
记 \(C_\lambda\subset S_{10}\)、\(C_\mu\subset S_9\)、\(C_\nu\subset S_3\) 为对应循环型共轭类。
则除去有限个分歧素数后，满足
\[
\Frob_p\in C_\lambda\times C_\mu\times C_\nu
\]
的素数集合具有自然密度
\[
\delta(\lambda,\mu,\nu)
=
\frac{|C_\lambda|}{10!}\cdot
\frac{|C_\mu|}{9!}\cdot
\frac{|C_\nu|}{6}.
\]
特别地，同时满足
\[
P_{10}\bmod p\ \text{不可约},
\qquad
P_9\bmod p\ \text{不可约},
\qquad
h_t\bmod p\ \text{不可约}
\]
的素数集合密度为
\[
\frac{1}{10}\cdot\frac{1}{9}\cdot\frac{1}{3}
=
\frac{1}{270},
\]
而三者同时全分裂的密度为
\[
\frac{1}{10!}\cdot \frac{1}{9!}\cdot \frac{1}{6}.
\]
\end{corollary}
\begin{proof}
由定理 \ref{thm:xi-time-part64bb-p10-p9-trigonal-threefold-galois-independence}，
\[
\Gal(L_{10}L_9M_t/\QQ)\cong S_{10}\times S_9\times S_3.
\]
对直积群应用 Chebotarev 密度定理，便得到一般共轭类的乘法密度公式。

对特例，只需注意：\(n\) 次多项式模 \(p\) 不可约对应 \(S_n\) 中的 \(n\)-循环，其比例为 \(1/n\)；全分裂对应恒等元类，其比例为 \(1/n!\)。
而在 \(S_3\) 中，不可约对应 \(3\)-循环类，大小为 \(2\)，故比例为 \(2/6=1/3\)。
代入即可得到 \(1/270\) 与全分裂密度公式。证毕。
\end{proof}

\begin{theorem}[三重判别特征立方体的等分布与三次分裂律的条件不变性]\label{thm:xi-time-part64bb-p10-p9-trigonal-signcube-equidistribution}
在定理 \ref{thm:xi-time-part64bb-p10-p9-trigonal-threefold-galois-independence} 的设定下，对不分歧素数 \(p\) 定义
\[
\chi_{10}(p):=\sgn(\sigma_{10}(p))\in\{\pm 1\},
\qquad
\chi_9(p):=\sgn(\sigma_9(p))\in\{\pm 1\},
\qquad
\chi_t(p):=\sgn(\sigma_t(p))\in\{\pm 1\},
\]
其中
\[
\sigma_{10}(p)\in S_{10},
\qquad
\sigma_9(p)\in S_9,
\qquad
\sigma_t(p)\in S_3
\]
分别表示相应分裂域上的 Frobenius 元。
则有：
\begin{enumerate}
  \item 三元符号向量
  \[
  \bigl(\chi_{10}(p),\chi_9(p),\chi_t(p)\bigr)
  \]
  在 \(\{\pm1\}^3\) 上完全等分布；每一符号模式的自然密度均为
  \[
  \frac{1}{8}.
  \]
  \item 对任意给定的双符号条件
  \[
  \chi_{10}(p)=\varepsilon_1,
  \qquad
  \chi_9(p)=\varepsilon_2
  \qquad
  (\varepsilon_1,\varepsilon_2\in\{\pm1\}),
  \]
  三次专化 \(h_t\bmod p\) 的三种分解型条件密度保持不变：
  \[
  \delta\bigl((1)(1)(1)\mid \chi_{10}=\varepsilon_1,\chi_9=\varepsilon_2\bigr)=\frac16,
  \]
  \[
  \delta\bigl((1)(2)\mid \chi_{10}=\varepsilon_1,\chi_9=\varepsilon_2\bigr)=\frac12,
  \]
  \[
  \delta\bigl((3)\mid \chi_{10}=\varepsilon_1,\chi_9=\varepsilon_2\bigr)=\frac13.
  \]
\end{enumerate}
换言之，十次分歧通道与九次分歧通道的全部二次判别偏置，对三次专化层的局部分裂统计完全透明。
\end{theorem}
\begin{proof}
由定理 \ref{thm:xi-time-part64bb-p10-p9-trigonal-threefold-galois-independence}，
\[
\Gal(L_{10}L_9M_t/\QQ)\cong S_{10}\times S_9\times S_3.
\]
对三个因子分别取符号特征，得到满射
\[
S_{10}\times S_9\times S_3\twoheadrightarrow C_2\times C_2\times C_2.
\]
Chebotarev 定理在该阿贝尔商上给出八个符号盒的等分布，故每种符号模式的密度都是 \(1/8\)。

再看条件分布。
给定 \((\chi_{10},\chi_9)\) 只会约束 \(S_{10}\times S_9\) 坐标，而 \(h_t\bmod p\) 的分解型仅由 \(S_3\) 坐标决定。
由于直积群上的 Frobenius 分布完全乘法化，条件化后 \(S_3\) 的共轭类比例不变。
而 \(S_3\) 的三个共轭类大小分别为 \(1,3,2\)，对应分解型
\[
(1)(1)(1),\qquad (1)(2),\qquad (3),
\]
故条件密度分别为
\[
\frac16,\qquad \frac12,\qquad \frac13.
\]
证毕。
\end{proof}

\begin{theorem}[十九次固定点统计与 \texorpdfstring{$(10,3)$}{(10,3)} 分裂指纹的三向独立]\label{thm:xi-time-part64bb-p10-leyang-h19-fixedpoint-tensorization}
沿用定理 \ref{thm:fold-gauge-anomaly-P10-leyang-H-triple-product} 的记号。
令 \(L_H/\QQ\) 为十九次指数域对应的 \(S_{19}\)-扩张，并对不分歧素数 \(p\) 记事件
\[
E_r(p):=\text{Frobenius 在 \(19\) 个根上恰有 \(r\) 个不动点}
\qquad
(0\le r\le 19),
\]
等价地，十九次多项式 \(H\bmod p\) 恰有 \(r\) 个一次因子。
再定义
\[
I_{10}(p):=\{\text{\(P_{10}\bmod p\) 不可约}\},
\]
\[
R(p):=\{\text{\(g(y)=256y^3+411y^2+165y+32\bmod p\) 至少有一根}\},
\]
\[
S(p):=\{\text{\(g\bmod p\) 全分裂}\}.
\]
则有精确密度公式
\[
\delta\bigl(I_{10}\cap R\cap E_r\bigr)
=
\frac1{10}\cdot \frac23\cdot \frac{\binom{19}{r}\,!(19-r)}{19!},
\]
以及
\[
\delta\bigl(I_{10}\cap S\cap E_r\bigr)
=
\frac1{10}\cdot \frac16\cdot \frac{\binom{19}{r}\,!(19-r)}{19!}.
\]
这表明：十次分歧的不可约性、Lee--Yang 三次层的根出现事件、以及十九次指数层的一次因子数，在自然密度意义下完全张量化。
\end{theorem}
\begin{proof}
由定理 \ref{thm:fold-gauge-anomaly-P10-leyang-H-triple-product}，
\[
\Gal(L_{10}L_{\mathrm{LY}}L_H/\QQ)\cong S_{10}\times S_3\times S_{19},
\]
而推论 \ref{cor:fold-gauge-anomaly-P10-leyang-H-chebotarev-triple-product} 给出三因子的 Frobenius 类分布完全乘法化。

另一方面，由命题 \ref{prop:fold-gauge-anomaly-H-rencontres-fixed-point-density}，
\[
\delta(E_r)=\frac{\binom{19}{r}\,!(19-r)}{19!}.
\]
在 \(S_{10}\) 中，\(I_{10}\) 对应 \(10\)-循环类，故
\[
\delta(I_{10})=\frac1{10}.
\]
在 \(S_3\) 中，“至少有一根”对应恒等类与换位类之并，故
\[
\delta(R)=\frac16+\frac12=\frac23,
\]
而“全分裂”对应恒等类，故
\[
\delta(S)=\frac16.
\]
将这三个边际密度相乘，即得
\[
\delta\bigl(I_{10}\cap R\cap E_r\bigr)
=
\delta(I_{10})\,\delta(R)\,\delta(E_r)
=
\frac1{10}\cdot \frac23\cdot \frac{\binom{19}{r}\,!(19-r)}{19!},
\]
以及
\[
\delta\bigl(I_{10}\cap S\cap E_r\bigr)
=
\delta(I_{10})\,\delta(S)\,\delta(E_r)
=
\frac1{10}\cdot \frac16\cdot \frac{\binom{19}{r}\,!(19-r)}{19!}.
\]
证毕。
\end{proof}

\endinput


\paragraph{bin-fold 满秩化、原子有限部位移与碰撞势内域大偏差}
在 bin-fold 两态一致渐近、规范群 Stirling 账本、真实输入 \(40\) 状态核的原子 Euler 剥离以及碰撞位势的解析/大偏差链已经建立之后，还可进一步抽取出下列十三条直接后果。它们分别刻画 fold-gauge 直积分解的中心与阿贝尔化精确结构、阿贝尔化的最终满秩化、中心的最终消失与纤维奇偶荷最小完备秩的严格稳定、规范复杂度相对异常维数的指数超载、体积密度对平均纤维信息项的线性响应斜率、纯碰撞误差指数缺口的 \(\ell^p\) 临界幂、原子因子的 primitive 与 Abel 有限部刚性位移、平方根归一化慢模的三歧相变与 \(u=1\) 邻域的统一亚临界性、RH-sharp 临界点上一阶滤子的精确消模、误差指数 \(\beta=\tfrac12\) 的余维一尖点接触，以及 RH 失效在 Taylor 解析域内部的可见性与均值点速率函数的显式二次胚。

\begin{theorem}[fold-gauge 直积分解的中心与阿贝尔化精确结构]\label{thm:xi-time-part65-binfold-gauge-center-abelianization-exact}
沿用命题 \ref{prop:fold-bin-gauge-decomposition} 的记号，记
\[
n_2(m):=\#\{w\in X_m:\ d_m^{\mathrm{bin}}(w)=2\},
\qquad
n_{\ge 2}(m):=\#\{w\in X_m:\ d_m^{\mathrm{bin}}(w)\ge 2\}.
\]
则有同构
\[
Z\!\bigl(G_m^{\mathrm{bin}}\bigr)\cong (\ZZ/2\ZZ)^{n_2(m)},
\qquad
\bigl(G_m^{\mathrm{bin}}\bigr)^{\mathrm{ab}}\cong (\ZZ/2\ZZ)^{n_{\ge 2}(m)}.
\]
特别地，每一条二点纤维都贡献一个中心 sheet-flip，而所有 \(d_m^{\mathrm{bin}}(w)\ge 3\) 的对称群因子都具有平凡中心。
\end{theorem}
\begin{proof}
由命题 \ref{prop:fold-bin-gauge-decomposition}，
\[
G_m^{\mathrm{bin}}
\cong
\prod_{w\in X_m}\mathfrak S_{d_m^{\mathrm{bin}}(w)}.
\]
对称群满足
\[
Z(\mathfrak S_1)=1,
\qquad
Z(\mathfrak S_2)=\mathfrak S_2\cong \ZZ/2\ZZ,
\qquad
Z(\mathfrak S_d)=1\quad(d\ge 3),
\]
以及
\[
\mathfrak S_1^{\mathrm{ab}}=1,
\qquad
\mathfrak S_d^{\mathrm{ab}}\cong \ZZ/2\ZZ\quad(d\ge 2).
\]
再由直积群的中心与阿贝尔化分别按因子直积分解，立得
\[
Z\!\bigl(G_m^{\mathrm{bin}}\bigr)\cong (\ZZ/2\ZZ)^{n_2(m)},
\qquad
\bigl(G_m^{\mathrm{bin}}\bigr)^{\mathrm{ab}}\cong (\ZZ/2\ZZ)^{n_{\ge 2}(m)}.
\]
证毕。
\end{proof}
\leanverified{paper\_xi\_time\_part65\_binfold\_gauge\_center\_abelianization\_exact}

\begin{theorem}[bin-fold 规范群阿贝尔化的最终满秩化]\label{thm:xi-time-part65-binfold-abel-eventual-full-rank}
沿用定理 \ref{thm:xi-time-part65-binfold-gauge-center-abelianization-exact} 与
\ref{thm:fold-bin-two-state-asymptotic} 的记号，并记
\[
N_m^{(2)}:=\#\{w\in X_m:\ d_m^{\mathrm{bin}}(w)\ge 2\}.
\]
则存在 \(m_0\)，使得对所有 \(m\ge m_0\) 都有
\[
\boxed{\;
\bigl(G_m^{\mathrm{bin}}\bigr)^{\mathrm{ab}}
\cong
(\ZZ/2\ZZ)^{F_{m+2}}.\;}
\]
等价地，bin-fold 规范群的全部稳定类型在充分高分辨率下都产生彼此独立的 \((\ZZ/2\ZZ)\)-奇偶通道，而阿贝尔化秩最终不再遗漏任何稳定类型。
\end{theorem}
\begin{proof}
由定理 \ref{thm:fold-bin-two-state-asymptotic} 的一致对数展开，
\[
\log d_m^{\mathrm{bin}}(w)
=
m\log\!\Bigl(\frac{2}{\varphi}\Bigr)-w_m\log\varphi
+O\!\Bigl(\Bigl(\frac{\varphi}{2}\Bigr)^m\Bigr)
\qquad (w\in X_m).
\]
由于 \(w_m\in\{0,1\}\)，故对全部 \(w\in X_m\) 一致有
\[
\log d_m^{\mathrm{bin}}(w)
\ge
m\log\!\Bigl(\frac{2}{\varphi}\Bigr)-\log\varphi
+O\!\Bigl(\Bigl(\frac{\varphi}{2}\Bigr)^m\Bigr).
\]
而 \(\log(2/\varphi)>0\)，故右端随 \(m\to\infty\) 一致趋于 \(+\infty\)。于是存在 \(m_0\) 使当 \(m\ge m_0\) 时，对所有 \(w\in X_m\) 都有 \(d_m^{\mathrm{bin}}(w)>1\)。由于 \(d_m^{\mathrm{bin}}(w)\) 为正整数，遂得
\[
d_m^{\mathrm{bin}}(w)\ge 2
\qquad(\forall\,w\in X_m,\ m\ge m_0),
\]
即
\[
N_m^{(2)}=|X_m|=F_{m+2}.
\]
由定理 \ref{thm:xi-time-part65-binfold-gauge-center-abelianization-exact}，
\[
\bigl(G_m^{\mathrm{bin}}\bigr)^{\mathrm{ab}}
\cong
(\ZZ/2\ZZ)^{n_{\ge 2}(m)}.
\]
再代入
\[
n_{\ge 2}(m)=N_m^{(2)}=F_{m+2},
\]
即得结论。证毕。
\end{proof}

\begin{theorem}[bin-fold 规范群中心的最终消失与纤维奇偶荷最小完备秩的严格稳定]\label{thm:xi-time-part65-binfold-center-vanishing-signature-stability}
沿用定理 \ref{thm:xi-time-part65-binfold-gauge-center-abelianization-exact} 与
\ref{thm:xi-time-part65-binfold-abel-eventual-full-rank} 的记号。则存在 \(m_0\)，使得对所有 \(m\ge m_0\) 都有
\[
\boxed{\;
Z\!\bigl(G_m^{\mathrm{bin}}\bigr)=1,
\qquad
\bigl(G_m^{\mathrm{bin}}\bigr)^{\mathrm{ab}}
\cong
(\ZZ/2\ZZ)^{F_{m+2}}.\;}
\]
从而
\[
\boxed{\;
\operatorname{rank}_{\FF_2}\!\bigl((G_m^{\mathrm{bin}})^{\mathrm{ab}}\bigr)=F_{m+2}.\;}
\]
并且定理 \ref{thm:xi-foldbin-parity-charge-split-exact-minrank} 中的纤维奇偶荷最小完备秩在同一范围内亦严格等于 \(F_{m+2}\)。
\end{theorem}
\begin{proof}
由定理 \ref{thm:fold-bin-two-state-asymptotic} 的一致两态展开，存在常数 \(C>0\) 使对充分大的 \(m\) 与任意 \(w\in X_m\) 都有
\[
d_m^{\mathrm{bin}}(w)
\ge
\left(\frac{2}{\varphi}\right)^m
\left(\varphi^{-1}-C\left(\frac{\varphi}{2}\right)^m\right).
\]
由于
\[
\left(\frac{2}{\varphi}\right)^m
\left(\varphi^{-1}-C\left(\frac{\varphi}{2}\right)^m\right)
=
\varphi^{-1}\left(\frac{2}{\varphi}\right)^m-C
\longrightarrow +\infty,
\]
故存在 \(m_0\)，使得当 \(m\ge m_0\) 时，对所有 \(w\in X_m\) 都有
\[
d_m^{\mathrm{bin}}(w)\ge 3.
\]

于是
\[
n_2(m)=0,
\qquad
n_{\ge 2}(m)=|X_m|=F_{m+2}.
\]
由定理 \ref{thm:xi-time-part65-binfold-gauge-center-abelianization-exact}，
\[
Z\!\bigl(G_m^{\mathrm{bin}}\bigr)\cong (\ZZ/2\ZZ)^{n_2(m)}=1,
\]
\[
\bigl(G_m^{\mathrm{bin}}\bigr)^{\mathrm{ab}}
\cong
(\ZZ/2\ZZ)^{n_{\ge 2}(m)}
\cong
(\ZZ/2\ZZ)^{F_{m+2}}.
\]
这就给出了中心消失与交换化维数公式。

最后，由定理 \ref{thm:xi-foldbin-parity-charge-split-exact-minrank}，纤维奇偶荷的最小完备秩恰等于
\[
N_m:=\#\{w\in X_m:\ d_m^{\mathrm{bin}}(w)\ge 2\}.
\]
而此处对所有 \(w\in X_m\) 已有 \(d_m^{\mathrm{bin}}(w)\ge 3\)，故
\[
N_m=|X_m|=F_{m+2}.
\]
于是最小完备秩在同一范围内严格稳定为 \(F_{m+2}\)。证毕。
\end{proof}

\begin{theorem}[规范群复杂度相对异常维数的指数超载律]\label{thm:xi-time-part65-gauge-complexity-anomaly-overload}
沿用定理 \ref{thm:xi-time-part65-binfold-abel-eventual-full-rank} 与定理 \ref{thm:xi-foldbin-gauge-density-exact-thermodynamic-constant} 的记号。则当 \(m\to\infty\) 时，
\[
\boxed{\;
\frac{\log_2|G_m^{\mathrm{bin}}|}
{\dim_{\FF_2}\!\bigl((G_m^{\mathrm{bin}})^{\mathrm{ab}}\bigr)}
=
\frac{\sqrt5}{\varphi^2\log 2}
\left(\frac{2}{\varphi}\right)^m
\left(
m\log\!\Bigl(\frac{2}{\varphi}\Bigr)
-1-\frac{\log\varphi}{1+\varphi^2}
+o(1)
\right).\;}
\]
特别地，
\[
\boxed{\;
\frac{\log_2|G_m^{\mathrm{bin}}|}
{\dim_{\FF_2}\!\bigl((G_m^{\mathrm{bin}})^{\mathrm{ab}}\bigr)}
\longrightarrow +\infty,\;}
\]
并且其主导增长尺度恰为 \((2/\varphi)^m\)。
\end{theorem}
\begin{proof}
由定理 \ref{thm:xi-foldbin-gauge-density-exact-thermodynamic-constant}，
\[
g_m
=
m\log\!\Bigl(\frac{2}{\varphi}\Bigr)
-1-\frac{\log\varphi}{1+\varphi^2}
+O\!\Bigl(m\Bigl(\frac{\varphi}{2}\Bigr)^m\Bigr).
\]
由于 \(g_m=2^{-m}\log|G_m^{\mathrm{bin}}|\)，故
\[
\log_2|G_m^{\mathrm{bin}}|
=
\frac{2^m}{\log 2}
\left(
m\log\!\Bigl(\frac{2}{\varphi}\Bigr)
-1-\frac{\log\varphi}{1+\varphi^2}
+o(1)
\right).
\]

另一方面，由定理 \ref{thm:xi-time-part65-binfold-abel-eventual-full-rank}，对充分大的 \(m\) 有
\[
\dim_{\FF_2}\!\bigl((G_m^{\mathrm{bin}})^{\mathrm{ab}}\bigr)=F_{m+2}.
\]
再用 Binet 公式
\[
F_{m+2}=\frac{\varphi^{m+2}}{\sqrt5}+o(\varphi^m),
\]
便得到
\[
\frac{\log_2|G_m^{\mathrm{bin}}|}
{\dim_{\FF_2}\!\bigl((G_m^{\mathrm{bin}})^{\mathrm{ab}}\bigr)}
=
\frac{\sqrt5}{\varphi^2\log 2}
\left(\frac{2}{\varphi}\right)^m
\left(
m\log\!\Bigl(\frac{2}{\varphi}\Bigr)
-1-\frac{\log\varphi}{1+\varphi^2}
+o(1)
\right).
\]
由于 \(\log(2/\varphi)>0\) 且 \(2/\varphi>1\)，右端趋于 \(+\infty\)，并且主导增长尺度正是 \((2/\varphi)^m\)。证毕。
\end{proof}

\begin{theorem}[规范体积密度与纤维信息项之间的线性响应斜率]\label{thm:xi-time-part65-binfold-gauge-linear-response-slope}
沿用命题 \ref{prop:fold-bin-gauge-density-two-scale} 的记号。则
\[
\boxed{\;
\lim_{m\to\infty}
\frac{g_m-\kappa_m+1}{m(\varphi/2)^m}
=
\frac{\varphi^2}{2\sqrt5}\log\!\Bigl(\frac{2}{\varphi}\Bigr).\;}
\]
并且利用 \(|X_m|=F_{m+2}\sim \varphi^{m+2}/\sqrt5\) 可进一步改写为
\[
\boxed{\;
\frac{2^m}{|X_m|}\,\bigl(g_m-\kappa_m+1\bigr)
=
\frac12\,m\log\!\Bigl(\frac{2}{\varphi}\Bigr)+O(1).\;}
\]
特别地，对充分大的 \(m\) 有
\[
g_m>\kappa_m-1.
\]
\end{theorem}
\begin{proof}
由命题 \ref{prop:fold-bin-gauge-density-two-scale}，
\[
g_m-\kappa_m+1
=
\frac{(\varphi/2)^m}{2\sqrt5}
\Bigl(
\varphi^2 m\log\!\Bigl(\frac{2}{\varphi}\Bigr)
+\varphi^2\log(2\pi)-\log\varphi
\Bigr)
+O\!\Bigl(\Bigl(\frac{\varphi}{2}\Bigr)^m\Bigr).
\]
两边同除以 \(m(\varphi/2)^m\)，得到
\[
\frac{g_m-\kappa_m+1}{m(\varphi/2)^m}
=
\frac{1}{2\sqrt5}
\left(
\varphi^2\log\!\Bigl(\frac{2}{\varphi}\Bigr)
+\frac{\varphi^2\log(2\pi)-\log\varphi}{m}
\right)
+O\!\left(\frac1m\right),
\]
从而首个极限成立。

另一方面，由 Binet 公式，
\[
\frac{|X_m|}{2^m}
=
\frac{F_{m+2}}{2^m}
=
\frac{\varphi^2}{\sqrt5}\Bigl(\frac{\varphi}{2}\Bigr)^m\bigl(1+o(1)\bigr),
\]
故
\[
\frac{2^m}{|X_m|}
=
\frac{\sqrt5}{\varphi^2}\Bigl(\frac{2}{\varphi}\Bigr)^m\bigl(1+o(1)\bigr).
\]
将其与上式相乘，便得
\[
\frac{2^m}{|X_m|}\,\bigl(g_m-\kappa_m+1\bigr)
=
\frac12\,m\log\!\Bigl(\frac{2}{\varphi}\Bigr)+O(1).
\]
由于 \(\log(2/\varphi)>0\)，右端对充分大 \(m\) 为正，因此 \(g_m-\kappa_m+1>0\)，即
\(
g_m>\kappa_m-1
\)。
证毕。
\end{proof}

\begin{theorem}[纯碰撞误差指数缺口的 \texorpdfstring{$\ell^p$}{lp} 临界幂为 \texorpdfstring{$p=2$}{p=2}]\label{thm:xi-time-part65-collision-error-gap-lp-threshold}
沿用推论 \ref{cor:real-input-40-collision-error-phase} 的记号，并把误差指数 \(\beta(u)\) 限制在整数参数 \(u=n\ge 3\) 上。则对任意 \(p>0\)，级数
\[
\sum_{n\ge 3}\bigl(1-\beta(n)\bigr)^p
\]
满足精确判别
\[
\boxed{\;
\sum_{n\ge 3}\bigl(1-\beta(n)\bigr)^p
\begin{cases}
<\infty,& p\ge 2,\\[2pt]
=\infty,& 0<p<2.
\end{cases}}
\]
换言之，临界幂指数恰为 \(p=2\)，并且边界点的额外对数增益已经足以恢复可和性。
\end{theorem}
\begin{proof}
由推论 \ref{cor:real-input-40-collision-error-phase}，
\[
\beta(u)
=
1-\frac{2}{\sqrt u\,\log u}
+O\!\left(\frac{1}{u\log u}\right)
\qquad (u\to\infty).
\]
因此存在常数 \(A,B>0\) 及充分大的 \(n_0\)，使得对所有 \(n\ge n_0\)，
\[
\frac{A}{\sqrt n\,\log n}
\le
1-\beta(n)
\le
\frac{B}{\sqrt n\,\log n}.
\]
从而
\[
\bigl(1-\beta(n)\bigr)^p
\asymp
\frac{1}{n^{p/2}(\log n)^p}.
\]
于是由积分判别或 Cauchy 凝聚判别，
\[
\sum_{n\ge n_0}\frac{1}{n^{p/2}(\log n)^p}
\]
在 \(p>2\) 时显然收敛，在 \(0<p<2\) 时显然发散；而在边界点 \(p=2\) 时，它等价于
\[
\sum_{n\ge n_0}\frac{1}{n(\log n)^2},
\]
该级数收敛。故所示判别成立。证毕。
\end{proof}

\begin{theorem}[原子 Euler 因子的 primitive 单点刚性与 Abel 有限部的精确乘法位移]\label{thm:xi-time-part65-atomic-primitive-abel-shift}
沿用命题 \ref{prop:atomic-2-orbit-split-logM} 与定理
\ref{thm:xi-atomic-witt-cycle-primitive-exact-peeling} 的记号。写
\[
P(z,u)=\sum_{n\ge 1}p_n(u)z^n,
\qquad
P_{\mathrm{core}}(z,u)=\sum_{n\ge 1}p_n^{\mathrm{core}}(u)z^n.
\]
则对一切 \(n\ge 1\) 都有
\[
\boxed{\;
p_n(u)=p_n^{\mathrm{core}}(u)+u\,\mathbf 1_{\{n=2\}}.\;}
\]
特别地，在未加权情形 \(u=1\) 下，
\[
\boxed{\;
p_n=p_n^{\mathrm{core}}+\mathbf 1_{\{n=2\}}.\;}
\]
进一步，若 \(\mathfrak M\) 与 \(\mathfrak M_{\mathrm{core}}\) 分别表示 \(u=1\) 时 full/core 的 Abel 有限部常数，则
\[
\boxed{\;
\log\mathfrak M
=
\log\mathfrak M_{\mathrm{core}}+\varphi^{-4},\qquad
\mathfrak M=e^{\varphi^{-4}}\mathfrak M_{\mathrm{core}}.\;}
\]
因此，原子因子 \((1-uz^2)^{-1}\) 在 primitive 层只作用于长度 \(2\) 的单一坐标，而在 Abel 有限部上只留下一个完全闭式的乘法位移。
\end{theorem}
\begin{proof}
由定理 \ref{thm:xi-atomic-witt-cycle-primitive-exact-peeling}，
\[
P(z,u)=u z^2+P_{\mathrm{core}}(z,u).
\]
逐项比较 \(z^n\) 的系数，便得
\[
p_n(u)=p_n^{\mathrm{core}}(u)+u\,\mathbf 1_{\{n=2\}},
\]
从而在 \(u=1\) 时有
\(
p_n=p_n^{\mathrm{core}}+\mathbf 1_{\{n=2\}}
\)。

另一方面，命题 \ref{prop:atomic-2-orbit-split-logM} 已给出
\[
\log\mathfrak M=\log\mathfrak M_{\mathrm{core}}+z_\star^2.
\]
在本核的未加权情形下，\(z_\star=\varphi^{-2}\)，故
\[
z_\star^2=\varphi^{-4}.
\]
将其代回并指数化，便得
\[
\mathfrak M=e^{\varphi^{-4}}\mathfrak M_{\mathrm{core}}.
\]
证毕。
\end{proof}

\begin{theorem}[碰撞位势的 RH 失效在原点 Taylor 解析域内可直接见证]\label{thm:xi-time-part65-collision-taylor-visible-nonrh-window}
沿用定理 \ref{thm:xi-rh-critical-before-nearest-complex-branch} 的记号。则有严格包含
\[
\boxed{\;
(\theta_R,R_\theta)
\subset
\Bigl\{
\theta\in\RR_{>0}:\ 
P_2(\theta)=\sum_{n\ge 0}\frac{P_2^{(n)}(0)}{n!}\theta^n,
\ \textsf{RH}(e^\theta)\ \text{失效}
\Bigr\}.\;}
\]
更一般地，对任意紧区间 \(K\subset(\theta_R,R_\theta)\)，上述 Taylor 级数都在 \(K\) 上绝对且一致收敛，而 kernel RH 条件在 \(K\) 上处处失效。
\end{theorem}
\begin{proof}
由定理 \ref{thm:xi-rh-critical-before-nearest-complex-branch}，
\[
\theta_R<R_\theta
\]
且对每个 \(\theta\in(\theta_R,R_\theta)\)，一方面 \(|\theta|<R_\theta\)，另一方面 \(u=e^\theta>u_R\)，故 \(\textsf{RH}(e^\theta)\) 已失效。

由同一定理，\(P_2\) 在 \(\theta=0\) 的 Taylor 收敛半径正是 \(R_\theta\)。因此对任意 \(|\theta|<R_\theta\)，都有
\[
P_2(\theta)=\sum_{n\ge 0}\frac{P_2^{(n)}(0)}{n!}\theta^n.
\]
特别地，这对全部 \(\theta\in(\theta_R,R_\theta)\) 成立，从而得到所示包含关系。

若 \(K\subset(\theta_R,R_\theta)\) 为紧区间，则可取 \(r\) 使
\[
\sup_{\theta\in K}|\theta|<r<R_\theta.
\]
于是幂级数在闭区间 \(K\) 上由半径 \(r\) 的 Cauchy 估计控制，从而绝对且一致收敛。与此同时，由 \(K\subset(\theta_R,R_\theta)\) 仍知 \(\textsf{RH}(e^\theta)\) 在 \(K\) 上处处失效。证毕。
\end{proof}

\begin{theorem}[RH 窗口内 Taylor 截断深度的 Cauchy 指数界]\label{thm:xi-time-part65-collision-taylor-depth-exp-bound}
沿用定理 \ref{thm:xi-rh-critical-before-nearest-complex-branch} 的记号。对任意
\[
\theta_R<r<R_\theta,
\]
记
\[
M(r):=\sup_{|\theta|=r}|P_2(\theta)|.
\]
若
\[
T_N(\theta):=\sum_{n=0}^N\frac{P_2^{(n)}(0)}{n!}\theta^n
\]
为 \(P_2\) 在 \(\theta=0\) 处的 \(N\) 阶 Taylor 多项式，则对一切 \(N\ge 0\) 有一致余项界
\[
\sup_{|\theta|\le \theta_R}|P_2(\theta)-T_N(\theta)|
\le
\frac{M(r)}{1-\theta_R/r}\left(\frac{\theta_R}{r}\right)^{N+1}.
\]
因此，为使 RH 窗口内的截断误差不超过任意给定的 \(\varepsilon\in(0,1)\)，只需取
\[
N\ge
\frac{\log(\varepsilon^{-1})}{\log(r/\theta_R)}+O_r(1).
\]
在形式极限 \(r\uparrow R_\theta\) 下，上述深度需求呈现严格指数级缩放
\[
N=
\Theta\!\left(
\frac{\log(\varepsilon^{-1})}{\log(R_\theta/\theta_R)}
\right).
\]
\end{theorem}
\begin{proof}
由定理 \ref{thm:xi-rh-critical-before-nearest-complex-branch}，\(P_2\) 在圆盘 \(|\theta|<R_\theta\) 内解析。故当 \(r<R_\theta\) 时，Cauchy 估计给出 Taylor 系数满足
\[
\left|\frac{P_2^{(n)}(0)}{n!}\right|
\le
\frac{M(r)}{r^n}.
\]
于是对任意 \(|\theta|\le \theta_R<r\)，有
\[
|P_2(\theta)-T_N(\theta)|
\le
\sum_{n\ge N+1}\frac{M(r)}{r^n}|\theta|^n
\le
\sum_{n\ge N+1}M(r)\left(\frac{\theta_R}{r}\right)^n
=
\frac{M(r)}{1-\theta_R/r}\left(\frac{\theta_R}{r}\right)^{N+1}.
\]
这就得到一致余项界。

若要求右端不超过 \(\varepsilon\)，对数化并整理即可得
\[
N\ge
\frac{\log(\varepsilon^{-1})}{\log(r/\theta_R)}+O_r(1).
\]
当 \(r\) 固定时，这是关于 \(\log(\varepsilon^{-1})\) 的线性增长；而当 \(r\uparrow R_\theta\) 时，主导比例常数逼近 \(\log(R_\theta/\theta_R)^{-1}\)，故得到所述 \(\Theta\)-级深度律。证毕。
\end{proof}

\begin{theorem}[碰撞位势大偏差率函数在均值点的显式二次胚]\label{thm:xi-time-part65-collision-rate-quadratic-germ}
沿用定理 \ref{thm:xi-real-input-40-collision-rate-cramer-edgeworth} 的记号，记
\[
\mu:=P_2'(0)=\frac{3-\sqrt5}{10},
\]
并把归一化 Legendre--Fenchel 对偶写为
\[
I_2(x):=\sup_{\theta\in\RR}\Bigl(x\theta-\bigl(P_2(\theta)-P_2(0)\bigr)\Bigr).
\]
则在 \(h\to 0\) 时，
\[
\boxed{\;
I_2(\mu+h)
=
\frac{125}{12\sqrt5-10}\,h^2+O(h^3).\;}
\]
等价地，若采用未归一化对偶
\[
I(x):=\sup_{\theta\in\RR}\bigl(x\theta-P_2(\theta)\bigr),
\]
则
\[
\boxed{\;
I(\mu+h)
=
-2\log\varphi
+\frac{125}{12\sqrt5-10}\,h^2+O(h^3).\;}
\]
\end{theorem}
\begin{proof}
由定理 \ref{thm:xi-real-input-40-collision-rate-cramer-edgeworth}，
\[
I_2(\mu+h)
=
\frac{h^2}{2\kappa_2}+O(h^3),
\qquad
\kappa_2=P_2''(0)=\frac{6\sqrt5-5}{125}.
\]
因此
\[
\frac{1}{2\kappa_2}
=
\frac{125}{12\sqrt5-10},
\]
从而得到第一条展开。

另一方面，
\[
I(x)=I_2(x)-P_2(0).
\]
而在 \(u=1\) 时，\(\rho(1)=\varphi^2\)，故
\[
P_2(0)=\log \rho(1)=\log \varphi^2=2\log\varphi.
\]
将其代回，便得未归一化版本
\[
I(\mu+h)
=
-2\log\varphi+\frac{125}{12\sqrt5-10}\,h^2+O(h^3).
\]
证毕。
\end{proof}

\paragraph{平方根归一化的三歧相变与 \texorpdfstring{$u=1$}{u=1} 审计窗口的统一亚临界性}
在纯碰撞三次核的 RH/Ramanujan 窗口、负谱支控制以及实参数审计窗口已经建立之后，还可把平方根归一化下的慢模行为与 \(u=1\) 邻域的统一亚临界性压缩为下列两条结果。

\begin{theorem}[平方根归一化的三歧相变]\label{thm:xi-collision-square-root-normalization-trichotomy}
\leanverified{paper\_xi\_collision\_square\_root\_normalization\_trichotomy}
沿用定理 \ref{thm:xi-time-part62d-collision-rh-negative-branch-control} 与命题 \ref{prop:real-input-40-collision-rhk-window} 的记号。记 Perron 根为 \(\lambda_+(u)\)，次谱半径为
\[
\Lambda(u):=\max\{|\lambda_0(u)|,|\lambda_-(u)|,1\},
\]
并定义平方根归一化的慢模比
\[
\mathfrak S_n(u):=
\frac{\Lambda(u)^n}{\lambda_+(u)^{n/2}}
\qquad
(n\ge 1).
\]
则存在唯一临界参数 \(u_R>1\)，满足
\[
\Lambda(u)\le \lambda_+(u)^{1/2}
\iff
0<u\le u_R,
\qquad
u_R\approx 3.59262181344.
\]
从而：
\begin{enumerate}
  \item 若 \(0<u<u_R\)，则 \(\mathfrak S_n(u)\) 以指数速率趋于 \(0\)；
  \item 若 \(u=u_R\)，则 \(\mathfrak S_n(u)\equiv 1\) 对一切 \(n\ge 1\) 成立；
  \item 若 \(u>u_R\)，则 \(\mathfrak S_n(u)\) 以指数速率发散到 \(+\infty\)。
\end{enumerate}
\end{theorem}
\begin{proof}
定义
\[
r(u):=\frac{\Lambda(u)}{\lambda_+(u)^{1/2}}.
\]
命题 \ref{prop:real-input-40-collision-rhk-window} 已给出
\[
r(u)\le 1
\iff
0<u\le u_R,
\]
且临界参数 \(u_R\) 在正实轴上唯一。于是对 \(u<u_R\)、\(u=u_R\) 与 \(u>u_R\) 分别有
\[
r(u)<1,
\qquad
r(u)=1,
\qquad
r(u)>1.
\]
另一方面，
\[
\mathfrak S_n(u)=r(u)^n.
\]
三种情形遂立即成立。证毕。
\end{proof}

\begin{corollary}[\texorpdfstring{$u=1$}{u=1} 审计窗口中的统一亚临界性]\label{cor:xi-collision-u1-audit-window-uniform-subcriticality}
注 \ref{rem:real-input-40-output-potential-degeneracy-phase} 在实参数方向上给出无退化阈值
\[
\theta_0=0.0533477827\ldots,
\qquad
|\theta|<\theta_0
\Longrightarrow
u=e^\theta\in(0.932837\ldots,1.054796\ldots).
\]
记
\[
I_0:=[e^{-\theta_0},e^{\theta_0}].
\]
又因
\[
u_R\approx 3.59262181344,
\]
故
\[
I_0\Subset (0,u_R).
\]
于是存在常数 \(q_0\in(0,1)\)，使得对一切 \(u\in I_0\) 与一切 \(n\ge 1\)，都有
\[
\Lambda(u)^n\le q_0^n\,\lambda_+(u)^{n/2}.
\]
\end{corollary}
\begin{proof}
由定理 \ref{thm:xi-collision-square-root-normalization-trichotomy}，函数
\[
r(u):=\frac{\Lambda(u)}{\lambda_+(u)^{1/2}}
\]
在区间 \((0,u_R)\) 上严格小于 \(1\)。另一方面，\(\Lambda(u)\) 与 \(\lambda_+(u)\) 都由有限维矩阵的特征值模给出，因此 \(r(u)\) 在正实轴上连续。

由 \(\theta_0=0.0533477827\ldots\) 可得
\[
e^{\theta_0}=1.054796\ldots<u_R,
\]
且显然 \(e^{-\theta_0}>0\)，故紧区间 \(I_0=[e^{-\theta_0},e^{\theta_0}]\) 满足
\[
I_0\Subset (0,u_R).
\]
于是可定义
\[
q_0:=\sup_{u\in I_0}r(u).
\]
由连续函数在紧集上取到上确界且 \(r(u)<1\) 于 \(I_0\) 上处处成立，遂得
\[
0<q_0<1.
\]
因此对任意 \(u\in I_0\) 与任意 \(n\ge 1\)，
\[
\Lambda(u)^n
=
r(u)^n\lambda_+(u)^{n/2}
\le
q_0^n\,\lambda_+(u)^{n/2}.
\]
证毕。
\end{proof}

\endinput

在平方根归一化的三歧相变之外，RH-sharp 临界点本身还允许进一步压缩出一个单步代数滤子与一个余维一的误差指数尖点。

\begin{theorem}[临界平方根模的一阶湮灭]\label{thm:xi-time-part65-critical-sqrt-mode-first-order-annihilation}
\leanverified{paper\_xi\_time\_part65\_critical\_sqrt\_mode\_first\_order\_annihilation}
沿用命题 \ref{prop:real-input-40-collision-rhk-window} 的记号，设
\[
u_\ast:=u_R,
\qquad
s:=-\lambda_-(u_\ast)>0.
\]
则
\[
\lambda_+(u_\ast)=s^2,
\qquad
\lambda_-(u_\ast)=-s.
\]
进一步，对任意形如
\[
a_n=A\,\lambda_+(u_\ast)^n+B\,\lambda_-(u_\ast)^n+r_n
\qquad (n\ge 0)
\]
的标量序列，只要存在 \(\eta<s\) 使
\[
r_n=O(\eta^n),
\]
则一阶滤子
\[
\mathcal L_s a_n:=a_{n+1}+s\,a_n
\]
满足
\[
\mathcal L_s a_n
=(s^2+s)A\,s^{2n}+\widetilde r_n,
\qquad
\widetilde r_n=O(\eta^n)=o(s^n).
\]
特别地，负谱模 \(\lambda_-(u_\ast)^n\) 被 \(\mathcal L_s\) 单步精确湮灭，而保留下来的主导项只来自临界 Perron 模 \(s^{2n}\)。
\end{theorem}
\begin{proof}
注 \ref{rem:collision-cubic-rhsharp-quintic} 已给出
\[
\lambda_+(u_\ast)=s^2,
\qquad
\lambda_-(u_\ast)=-s.
\]
因此
\[
\mathcal L_s\bigl(\lambda_-(u_\ast)^n\bigr)
=(-s)^{n+1}+s(-s)^n
=0,
\]
而
\[
\mathcal L_s\bigl(\lambda_+(u_\ast)^n\bigr)
=s^{2n+2}+s^{2n+1}
=(s^2+s)s^{2n}.
\]
若 \(r_n=O(\eta^n)\) 且 \(\eta<s\)，则
\[
\mathcal L_s r_n=r_{n+1}+s\,r_n=O(\eta^n)=o(s^n).
\]
将上述三式按线性叠加，便得到所示展开。证毕。
\end{proof}

\begin{corollary}[平方根律的尖点接触与非厚相性]\label{cor:xi-time-part65-sqrt-law-cusp-nonthick-phase}
沿用推论 \ref{cor:real-input-40-collision-error-phase} 的记号，定义
\[
\beta(u):=\frac{\log \Lambda(u)}{\log \lambda_+(u)}.
\]
则
\[
\beta(u)=0
\qquad (0<u\le 1),
\]
\[
\beta(u)\le \frac12
\qquad (1<u\le u_R),
\]
且等号仅在 \(u=u_R\) 发生；此外，
\[
\beta(u)>\frac12
\qquad (u>u_R),
\qquad
\beta(u)\to 1
\qquad (u\to\infty).
\]
因此，\(\beta=\frac12\) 的平方根级误差并不形成参数上的开相区，而只在唯一临界参数 \(u_R\) 处与相图发生余维一接触。
\end{corollary}
\begin{proof}
推论 \ref{cor:real-input-40-collision-error-phase} 已逐段给出 \(\beta(u)\) 在
\((0,1]\)、\((1,u_R]\) 与 \((u_R,\infty)\) 上的精确刻画，并证明
\(\beta(u)\to 1\) 当 \(u\to\infty\)。将其直接合并，便得到全部结论。证毕。
\end{proof}

\endinput


\noindent
这十二条结果把 bin-fold 末位两态律、规范群 Stirling 账本、原子 Euler 手术与碰撞位势的解析几何再度压缩为同一条闭合主线：一方面，阿贝尔化在高分辨率极限中达到全部稳定类型的最终满秩，而中心在同一极限中彻底消失，从而纤维奇偶荷的最小完备秩被严格锁定为 \(F_{m+2}\)；进一步，完整规范群复杂度相对于可见异常维数并不同阶，而是以 \((2/\varphi)^m\) 呈现严格的指数超载；另一方面，其与平均纤维信息项之间仍保留一条由黄金比锁定的线性响应斜率，纯碰撞误差指数的缺口则只以临界幂 \(p=2\) 的精确边界可和方式衰减；再者，原子因子在 primitive 与 Abel 有限部两层都表现为完全局域的单点位移，而平方根归一化的慢模比不仅在唯一临界参数 \(u_R\) 处发生严格三歧相变，而且该临界点本身已退化为一个可由一阶滤子单步消去的负谱节点，\(\beta=\tfrac12\) 也只以余维一尖点方式触及相图；最后，RH 失效不仅先于最近复分歧发生，而且整个失效区间都仍由 \(\theta=0\) 的 Taylor 胚直接控制，而均值点的大偏差率函数曲率常数同样被 \(\sqrt5\) 完全闭式锁定。由此，bin-fold 的群论刚性、原子谱的局域手术与碰撞位势的局部热力学几何便被统一收束到同一组可审计的终端后果之中。

\paragraph{Lee--Yang 三次闭包、Latt\`es 六次正交、\texorpdfstring{$\chi^2$}{chi2} 缺口与 Jacobian 半单刚性}
在 Lee--Yang 分支不变量的共同三次数域、bin-fold 碰撞缺口的 Fibonacci 共振强制，以及 \(S_4\)--Hurwitz Jacobian 的可见等型块分解已经确立之后，还可进一步抽取出下列一组彼此兼容的刚性后果。

\begin{theorem}[Lee--Yang 分支几何的单参数有理统一化]\label{thm:xi-time-part65b-leyang-branch-b-birational-uniformization}
\leanverified{paper\_killo\_leyang\_b\_square\_cubic\_and\_rational\_intertwine}
沿用定理 \ref{thm:xi-terminal-zm-b2-kappa2-rational-interconstruction}、定理 \ref{thm:xi-terminal-zm-stokes-t-elliptic-branch-coordinate} 与定理 \ref{thm:xi-terminal-zm-leyang-elliptic-structure} 的记号。令 \((\lambda_*,y_*)\) 为非平凡 Lee--Yang 分歧点，并记
\[
u:=\kappa_*^2,
\qquad
b:=B^2.
\]
则 \(u\)、\(\lambda_*\) 与 \(y_*\) 都可写成 \(b\) 的有理函数；具体地，
\[
u=\frac{2294-8463b-918b^2}{14388},
\]
\[
\lambda_*
=
\frac{918b^2+8463b+4900}{2(918b^2+8463b+7298)},
\]
\[
y_*
=
-\frac{3597\,b\,(918b^2+8463b+4900)^2}
{(918b^2+8463b-2294)(918b^2+8463b+7298)^2}.
\]
因此
\[
\QQ(b)=\QQ(u)=\QQ(\lambda_*)=\QQ(y_*).
\]
换言之，Lee--Yang 分支点的局部几何不变量 \(u,\lambda_*,y_*,b\) 并非只是在同一三次数域中同居，而可由单一参数 \(b\) 作显式有理统一化。
\end{theorem}
\begin{proof}
由定理 \ref{thm:xi-terminal-zm-b2-kappa2-rational-interconstruction}，
\[
b
=
-2\cdot
\frac{24708u^2-28628u+7733}{51084u^2-44412u+8735}.
\]
同一定理并已给出 \(\QQ(b)=\QQ(u)\)，故 \(u\in \QQ(b)\) 必可唯一写成
\[
u=x_0+x_1b+x_2b^2
\qquad (x_0,x_1,x_2\in \QQ).
\]
把这一待定表达代回上式，并结合 \(b\) 的三次关系
\[
162b^3+1593b^2+1998b+1184=0
\]
在 \(\QQ(b)\) 中约化，即唯一解得
\[
u=\frac{2294-8463b-918b^2}{14388}.
\]

再由定理 \ref{thm:xi-terminal-zm-kappa-square-cubic-field-s3}，
\[
\lambda_*=\frac{3(2u-1)}{4(3u-2)}.
\]
把前式代入并化简，便得到所示 \(\lambda_*(b)\) 的表达。

另一方面，定理 \ref{thm:xi-terminal-zm-stokes-t-elliptic-branch-coordinate} 给出
\[
4\lambda_*y_*=4\lambda_*^3-3\lambda_*^2-2\lambda_*+1.
\]
再把 \(\lambda_*(b)\) 代入并整理，即得所示 \(y_*(b)\) 的表达。

由定理 \ref{thm:xi-terminal-zm-b2-kappa2-rational-interconstruction} 已知
\[
\QQ(b)=\QQ(u)=\QQ(\lambda_*).
\]
又由于 \(y_*\) 是定理 \ref{thm:xi-time-part9g-leyang-cubic-discriminant} 中不可约三次
\[
P_{\mathrm{LY}}(y)=256y^3+411y^2+165y+32
\]
的根，故 \([\QQ(y_*):\QQ]=3\)。而上式已证明 \(y_*\in \QQ(b)\)，于是只能有
\[
\QQ(y_*)=\QQ(b).
\]
遂得全部域相等式。证毕。
\end{proof}

\begin{theorem}[Lee--Yang 分支不变量链的共同二次分辨子域与统一 \texorpdfstring{$S_3$}{S3} 闭包]\label{thm:xi-time-part65b-leyang-common-s3-closure}
沿用定理 \ref{thm:xi-time-part65b-leyang-branch-b-birational-uniformization} 的记号，记
\[
u:=\kappa_*^2,
\qquad
b:=B^2,
\qquad
K:=\QQ(u)=\QQ(b)=\QQ(\lambda_*)=\QQ(y_*).
\]
设 \(u\) 与 \(b\) 的三次最小多项式分别为
\[
324u^3-540u^2+333u-74,
\qquad
162b^3+1593b^2+1998b+1184,
\]
其判别式分别记为
\[
\Delta_u=-2^6\cdot 3^9\cdot 37,
\qquad
\Delta_b=-2^2\cdot 3^9\cdot 11^2\cdot 37\cdot 109^2.
\]
若 \(L_{\mathrm{LY}}\) 记为三次数域 \(K/\QQ\) 的伽罗瓦闭包，则有
\[
\mathrm{sf}(\Delta_u)=\mathrm{sf}(\Delta_b)=-111,
\]
并且
\[
\Gal(L_{\mathrm{LY}}/\QQ)\cong S_3,
\qquad
L_{\mathrm{LY}}^{A_3}=\QQ(\sqrt{-111}).
\]
等价地，\(u\)、\(b\)、\(\lambda_*\) 与 \(y_*\) 都落在同一非伽罗瓦三次数域 \(K\) 中，其共同伽罗瓦闭包唯一的非平凡二次正规子域恰为 \(\QQ(\sqrt{-111})\)。
\end{theorem}
\begin{proof}
由定理 \ref{thm:xi-time-part65b-leyang-branch-b-birational-uniformization} 与定理 \ref{thm:xi-terminal-zm-stokes-leyang-shared-artin-representation}，
\[
\QQ(u)=\QQ(b)=\QQ(\lambda_*)=\QQ(y_*)=:K,
\qquad
\Gal(L_{\mathrm{LY}}/\QQ)\cong S_3.
\]
因此 \(u\)、\(b\)、\(\lambda_*\) 与 \(y_*\) 都生成同一非伽罗瓦三次数域 \(K\)，其分裂域必同为 \(L_{\mathrm{LY}}\)。

另一方面，
\[
\Delta_u
=
-2^6\cdot 3^9\cdot 37
=
(2^3\cdot 3^4)^2\cdot (-111),
\]
\[
\Delta_b
=
-2^2\cdot 3^9\cdot 11^2\cdot 37\cdot 109^2
=
(2\cdot 3^4\cdot 11\cdot 109)^2\cdot (-111),
\]
故
\[
\mathrm{sf}(\Delta_u)=\mathrm{sf}(\Delta_b)=-111.
\]
对不可约三次域，其伽罗瓦闭包的唯一非平凡二次正规子域由判别式的平方自由核决定，因此
\[
L_{\mathrm{LY}}^{A_3}=\QQ(\sqrt{-111}).
\]
证毕。
\end{proof}

\begin{proposition}[Lee--Yang 三次数域的两组显式整生成元]\label{prop:xi-time-part65b-leyang-u-b-integral-generators}
沿用定理 \ref{thm:xi-time-part65b-leyang-common-s3-closure} 的记号，定义
\[
x:=18u,
\qquad
y:=18b.
\]
则 \(x\) 与 \(y\) 都是代数整数，并分别满足首一整系数三次方程
\[
x^3-30x^2+333x-1332=0,
\]
\[
y^3+177y^2+3996y+42624=0.
\]
特别地，
\[
\QQ(x)=\QQ(u)=K,
\qquad
\QQ(y)=\QQ(b)=K.
\]
\end{proposition}
\begin{proof}
由定理 \ref{thm:xi-time-part65b-leyang-common-s3-closure}，
\[
324u^3-540u^2+333u-74=0.
\]
代入 \(u=x/18\) 并清分母，得到
\[
x^3-30x^2+333x-1332=0.
\]
同理，由
\[
162b^3+1593b^2+1998b+1184=0
\]
代入 \(b=y/18\) 并清分母，得到
\[
y^3+177y^2+3996y+42624=0.
\]
两式皆为首一整系数多项式，故 \(x\) 与 \(y\) 都是代数整数。
又因 \(18\in \QQ^\times\)，缩放不改变所生成的 \(\QQ\)-子域，因此
\[
\QQ(x)=\QQ(u)=K,
\qquad
\QQ(y)=\QQ(b)=K.
\]
证毕。
\end{proof}

\begin{corollary}[共同 \texorpdfstring{$S_3$}{S3} 闭包与 Latt\`es 六次域的线性互素性]\label{cor:xi-time-part65b-leyang-lattes-linear-disjointness}
设
\[
C(X)=X^6-8X^5+24X^4-36X^3+34X^2-20X+4
\]
的分裂域为 \(L_4\)，并满足
\[
\Disc(C)=-2^{19}\cdot 37^3,
\qquad
\Gal(L_4/\QQ)\cong S_4\times C_2.
\]
则定理 \ref{thm:xi-time-part65b-leyang-common-s3-closure} 中的共同闭包 \(L_{\mathrm{LY}}\) 与 \(L_4\) 线性互素：
\[
L_{\mathrm{LY}}\cap L_4=\QQ.
\]
因而
\[
\Gal(L_{\mathrm{LY}}L_4/\QQ)
\cong
S_3\times (S_4\times C_2).
\]
\end{corollary}
\begin{proof}
由定理 \ref{thm:xi-time-part65b-leyang-common-s3-closure}，
\[
L_{\mathrm{LY}}^{A_3}=\QQ(\sqrt{-111}),
\]
这是 \(L_{\mathrm{LY}}/\QQ\) 的唯一非平凡二次正规子域。于是交域
\[
F:=L_{\mathrm{LY}}\cap L_4
\]
作为 \(L_{\mathrm{LY}}/\QQ\) 的伽罗瓦子扩张，只能是 \(\QQ\) 或 \(\QQ(\sqrt{-111})\)。

若 \(F=\QQ(\sqrt{-111})\)，则素数 \(3\) 在 \(F/\QQ\) 中分歧；然而由
\[
\Disc(C)=-2^{19}\cdot 37^3
\]
可知 \(L_4/\QQ\) 只在 \(\{2,37\}\) 处分歧，故 \(3\) 在 \(L_4\) 的任一子扩张中都不分歧，矛盾。故只能有
\[
L_{\mathrm{LY}}\cap L_4=\QQ.
\]
最后，两边皆为 \(\QQ\) 上的伽罗瓦扩张，故复合域的伽罗瓦群为直积：
\[
\Gal(L_{\mathrm{LY}}L_4/\QQ)
\cong
\Gal(L_{\mathrm{LY}}/\QQ)\times \Gal(L_4/\QQ)
\cong
S_3\times (S_4\times C_2).
\]
证毕。
\end{proof}

\begin{theorem}[临界共振常数强制 \texorpdfstring{$\chi^2$}{chi2} 碰撞缺口具有正的极限下界]\label{thm:xi-time-part65b-cphi-chi2-gap-floor}
沿用定义 \ref{def:fold-normalized-amplitude}、命题 \ref{prop:fold-bin-chi2-col} 与定理 \ref{thm:fold-critical-resonance-constant} 的记号。记
\[
a_m(k):=\frac{|\widehat d_m(k)|}{2^m},
\qquad
C_\varphi:=\lim_{m\to\infty}a_m(F_m)=\lim_{m\to\infty}a_m(F_{m+1})>0.
\]
若 \(p_m^{\mathrm{bin}}\) 为 bin-fold 推前分布，\(u_m\) 为 \(X_m\) 上的均匀分布，则
\[
\liminf_{m\to\infty}\chi^2\!\bigl(p_m^{\mathrm{bin}}\|u_m\bigr)\ge 2C_\varphi^2.
\]
等价地，
\[
\liminf_{m\to\infty}
F_{m+2}\Bigl(\mathsf{Col}_m-\frac{1}{F_{m+2}}\Bigr)
\ge
2C_\varphi^2.
\]
\end{theorem}
\begin{proof}
由命题 \ref{prop:fold-bin-chi2-col}，
\[
\chi^2\!\bigl(p_m^{\mathrm{bin}}\|u_m\bigr)
=
|X_m|\,\mathsf{Col}_m-1
=
F_{m+2}\Bigl(\mathsf{Col}_m-\frac{1}{F_{m+2}}\Bigr).
\]
另一方面，结论 \ref{con:xi-time-part9m-cphi-forces-collision-gap-and-peak-bias} 已给出
\[
\liminf_{m\to\infty}
F_{m+2}\Bigl(\mathsf{Col}_m-\frac{1}{F_{m+2}}\Bigr)
\ge
2C_\varphi^2.
\]
将二者合并即得
\[
\liminf_{m\to\infty}\chi^2\!\bigl(p_m^{\mathrm{bin}}\|u_m\bigr)\ge 2C_\varphi^2.
\]
第二式只是上述恒等式的等价改写。证毕。
\end{proof}

\begin{theorem}[可见等型块代数的显式半单嵌入与 N\'eron--Severi 秩下界]\label{thm:xi-time-part65b-jacobian-visible-semisimple-ns-floor}
沿用定理 \ref{thm:killo-s4-genus49-jacobian-computable-isogeny} 的 \(\QQ\)-同源分解
\[
J(X)\sim_{\QQ}J(H)\times J(Z)^2\times J(C_F)^3\times P^3.
\]
则 \(\operatorname{End}_{\QQ}^0(J(X))\) 含有显式半单 \(\QQ\)-子代数
\[
\mathcal A_X:=\QQ\times M_2(\QQ)\times M_3(\QQ)\times M_3(\QQ)
\hookrightarrow
\operatorname{End}_{\QQ}^0(J(X)),
\]
从而
\[
\dim_{\QQ}\operatorname{End}_{\QQ}^0(J(X))\ge 1+4+9+9=23.
\]
进一步，在右端乘积上取任意乘积极化，并把相应 Rosati 对合记为 \(\dagger_{\mathrm{vis}}\)。则
\[
\dim_{\QQ}(\mathcal A_X)^{\dagger_{\mathrm{vis}}}
=
1+\frac{2\cdot 3}{2}+\frac{3\cdot 4}{2}+\frac{3\cdot 4}{2}
=
16.
\]
特别地，
\[
\operatorname{rank}\operatorname{NS}(J(X)_{\overline{\QQ}})\ge 16.
\]
\end{theorem}
\begin{proof}
定理 \ref{thm:xi-time-part9n1b-jacobian-projector-idempotent-algebra} 已经给出显式嵌入
\[
\mathcal A_X
\hookrightarrow
\operatorname{End}_{\QQ}^0(J(X))
\]
以及
\[
\dim_{\QQ}\operatorname{End}_{\QQ}^0(J(X))\ge 23.
\]
故只需证明 Rosati 不动维数与 N\'eron--Severi 秩下界。

取一个实现上述 \(\QQ\)-同源分解的同源映射，并在乘积
\[
A\times B^2\times C^3\times D^3
\qquad
\bigl(A:=J(H),\ B:=J(Z),\ C:=J(C_F),\ D:=P\bigr)
\]
上选取乘积极化 \(\lambda_{\mathrm{prod}}\)。在该模型下，\(\mathcal A_X\) 恰由四个可见等型块的重数空间上的标量矩阵作用给出，因此 \(\lambda_{\mathrm{prod}}\) 的 Rosati 对合在 \(\QQ\) 因子上恒等，在 \(M_r(\QQ)\) 因子上与转置等价。于是其不动子空间维数分别为
\[
1,\qquad \frac{2\cdot 3}{2}=3,\qquad \frac{3\cdot 4}{2}=6,\qquad \frac{3\cdot 4}{2}=6,
\]
从而
\[
\dim_{\QQ}(\mathcal A_X)^{\dagger_{\mathrm{vis}}}=1+3+6+6=16.
\]

最后，同源映射的拉回在 N\'eron--Severi 群上于有理化后给出同构，而对任意极化阿贝尔簇 \(A\) 均有标准识别
\[
\operatorname{NS}(A_{\overline{\QQ}})\otimes_{\ZZ}\QQ
\cong
\operatorname{End}^0_{\overline{\QQ}}(A)^{\dagger=1}.
\]
故上述 \(16\) 维 Rosati 不动子空间传回 \(J(X)\) 后仍给出
\[
\operatorname{rank}\operatorname{NS}(J(X)_{\overline{\QQ}})\ge 16.
\]
证毕。
\end{proof}

\noindent
以上五条结果把 Lee--Yang 分支几何的单参数有理统一化、共同 \(S_3\)-闭包、Latt\`es 六次闭包、黄金共振二阶缺口与 \(S_4\)--Hurwitz Jacobian 的可见块分解进一步压缩到同一层级：一方面，\(b\) 不仅与 \(u\)、\(\lambda_*\) 及 \(y_*\) 生成同一三次数域，而且把全部分支点局部几何显式统一化；该共同三次数域的伽罗瓦闭包唯一的非平凡二次正规子域仍被 \(\QQ(\sqrt{-111})\) 完全锁定，并与 \(L_4\) 严格正交；另一方面，临界共振常数把 bin-fold 与均匀基线之间的 \(\chi^2\) 偏离锁定为严格正的极限底噪，而 Jacobian 的可见等型块则单独强制出一个 \(23\) 维有理半单子代数与至少 \(16\) 维的 N\'eron--Severi 秩下界。

\endinput


\paragraph{escort 冻结图、归一化 residue law 与输出位势的一阶非高斯非对称}
在冻结相压力谱的线性支正、escort 态向最大纤维层的指数集中、真实输入 \(40\) 状态核的单原子 primitive/residue 剥离，以及输出位势的准幂极限定律与局部 Legendre 正规形已经建立之后，还可把这些链进一步压缩为一组互相衔接的结构后果：首先，escort--Rényi 熵率不仅在 \((a,s)\)-平面上给出双参数冻结边界，而且在更强阈值上提升为由 \(\log K_m\) 控制的有限尺度主项，并进一步诱导出最大纤维扇区在预言机容量中的显式指数贡献与侧信息临界斜率；随后，单原子的秩一局域手术被提升为归一化 periodic residue law 的精确指数、正则单纯形首项轨道与 centered residue 各向同性；最后，中心极限层的一阶 Edgeworth 偏斜与局部大偏差层的三次不对称项仍由同一输出位势常数 \(\kappa_3\) 严格锁定。

\leanverified{paper\_xi\_time\_part65c\_escort\_renyi\_two\_parameter\_freezing\_line}
\begin{theorem}[escort--Rényi 熵的双参数冻结线]\label{thm:xi-time-part65c-escort-renyi-two-parameter-freezing-line}
沿用定理 \ref{thm:xi-fixed-freezing-escort-renyi-spectrum-collapse} 与定理 \ref{thm:fold-pressure-freezing-threshold} 的记号。对整数 \(a\ge 1\) 定义 escort 分布
\[
\pi_{m,a}(x):=\frac{d_m(x)^a}{S_a(m)},
\qquad
S_a(m):=\sum_{x\in X_m}d_m(x)^a.
\]
对整数 \(s\ge 2\) 定义其 Rényi 熵
\[
H_s(\pi_{m,a})
:=
\frac{1}{1-s}\log\sum_{x\in X_m}\pi_{m,a}(x)^s.
\]
若整数压力极限
\[
P_q:=\lim_{m\to\infty}\frac{1}{m}\log S_q(m)
\qquad(q\ge 1)
\]
存在，则 escort--Rényi 熵率
\[
\mathfrak h_s(a):=
\lim_{m\to\infty}\frac{1}{m}H_s(\pi_{m,a})
\]
存在且满足严格闭式
\[
\boxed{\;
\mathfrak h_s(a)=\frac{sP_a-P_{as}}{s-1}.\;}
\]
特别地：
\begin{enumerate}
  \item 若 \(a>a_{\mathrm c}\)，则对每个整数 \(s\ge 2\) 都有
  \[
  \boxed{\;\mathfrak h_s(a)=g_*.\;}
  \]
  \item 更一般地，若 \(a\le a_{\mathrm c}\) 但 \(as>a_{\mathrm c}\)，则
  \[
  \boxed{\;
  \mathfrak h_s(a)
  =
  \frac{sP_a-(as\,\alpha_*+g_*)}{s-1}.\;}
  \]
\end{enumerate}
因此，原始压力谱的冻结阈值位于 \(a=a_{\mathrm c}\)，而 escort--Rényi 熵的相应双参数冻结边界则位于
\[
\boxed{\;as=a_{\mathrm c}.\;}
\]
\end{theorem}
\begin{proof}
由定义，
\[
\sum_{x\in X_m}\pi_{m,a}(x)^s
=
\sum_{x\in X_m}\frac{d_m(x)^{as}}{S_a(m)^s}
=
\frac{S_{as}(m)}{S_a(m)^s}.
\]
故
\[
H_s(\pi_{m,a})
=
\frac{1}{1-s}\bigl(\log S_{as}(m)-s\log S_a(m)\bigr).
\]
两边除以 \(m\) 并令 \(m\to\infty\)，即得
\[
\mathfrak h_s(a)=\frac{sP_a-P_{as}}{s-1}.
\]

若 \(a>a_{\mathrm c}\)，则由于 \(s\ge 2\)，自动有 \(as>a>a_{\mathrm c}\)。由定理 \ref{thm:fold-pressure-freezing-threshold}，
\[
P_a=a\alpha_*+g_*,
\qquad
P_{as}=as\,\alpha_*+g_*.
\]
代回上式得到
\[
\mathfrak h_s(a)
=
\frac{s(a\alpha_*+g_*)-(as\,\alpha_*+g_*)}{s-1}
=
g_*.
\]

若 \(a\le a_{\mathrm c}\) 但 \(as>a_{\mathrm c}\)，则仍有
\[
P_{as}=as\,\alpha_*+g_*,
\]
而 \(P_a\) 保持一般形式，代回即得第二条。最后，由 \(P_{as}\) 是否进入冻结支正只取决于 \(as\) 是否超过 \(a_{\mathrm c}\)，故双参数冻结边界正是直线 \(as=a_{\mathrm c}\)。证毕。
\end{proof}

\begin{corollary}[冻结相内全部整数阶 Rényi 熵率塌缩到同一常数]\label{cor:xi-time-part65c-escort-all-integer-renyi-collapse}
\leanverified{paper\_xi\_time\_part65c\_escort\_all\_integer\_renyi\_collapse}
沿用定理 \ref{thm:xi-time-part65c-escort-renyi-two-parameter-freezing-line} 的记号。若 \(a>a_{\mathrm c}\)，则对每个整数 \(s\ge 2\) 都有
\[
\boxed{\;
\lim_{m\to\infty}\frac{1}{m}H_s(\pi_{m,a})=g_*.\;}
\]
再与定理 \ref{thm:xi-fixed-freezing-escort-renyi-spectrum-collapse} 在 \(q=\infty\) 的结论合并，得到：在冻结区 \(a>a_{\mathrm c}\) 内，escort 分布 \(\pi_{m,a}\) 的全部整数阶 Rényi 熵率，包括 \(s=\infty\) 的 \(\min\)-entropy 率，都在极限中塌缩为同一常数 \(g_*\)。
\end{corollary}
\begin{proof}
对任意整数 \(s\ge 2\)，由 \(a>a_{\mathrm c}\) 可知 \(as>a_{\mathrm c}\)。于是定理 \ref{thm:xi-time-part65c-escort-renyi-two-parameter-freezing-line} 的第一条直接给出
\[
\lim_{m\to\infty}\frac{1}{m}H_s(\pi_{m,a})=g_*.
\]
另一方面，定理 \ref{thm:xi-fixed-freezing-escort-renyi-spectrum-collapse} 的第三条已给出
\[
\lim_{m\to\infty}\frac{1}{m}H_\infty(\pi_{m,a})=g_*.
\]
二者合并即得结论。证毕。
\end{proof}

\begin{theorem}[冻结相 escort--Rényi 熵的有限尺度主项由基态简并度精确控制]\label{thm:xi-time-part65c-escort-renyi-finite-size-mainterm}
\leanverified{paper\_xi\_time\_part65c\_escort\_renyi\_finite\_size\_mainterm}
沿用结论 \ref{con:xi-time-part57ba-freezing-escort-groundstate-fdiv-closed}、推论 \ref{cor:xi-fixed-freezing-escort-groundstate-tv-identity} 与定理 \ref{thm:xi-fixed-freezing-power-sum-ratio-locking} 的记号。对任意实数 \(s>1\)，定义更强阈值
\[
a_{\mathrm c,s}:=
\frac{\log\varphi-\frac{1}{s}g_*}{\alpha_*-\alpha_*^{(2)}}.
\]
则当 \(a>a_{\mathrm c,s}\) 时，随着 \(m\to\infty\)，有
\[
\boxed{\;
\sum_{x\in X_m}\pi_{m,a}(x)^s
=
K_m^{1-s}\bigl(1+o(1)\bigr),\;}
\]
\[
\boxed{\;
H_s(\pi_{m,a})=\log K_m+o(1).\;}
\]
换言之，在半直线 \(a>a_{\mathrm c,s}\) 上，冻结 escort--Rényi 熵不仅在指数率上塌缩到 \(g_*\)，其完整有限尺度主项也被基态简并度 \(\log K_m\) 精确控制。
\end{theorem}
\begin{proof}
由结论 \ref{con:xi-time-part57ba-freezing-escort-groundstate-fdiv-closed}，可写
\[
\pi_{m,a}=p_{m,a}\sigma_m+(1-p_{m,a})\eta_{m,a},
\]
其中 \(\sigma_m=\mathrm{Unif}(A_m^{\max})\)，且 \(\sigma_m\) 与 \(\eta_{m,a}\) 的支撑不交。故
\[
\sum_{x\in X_m}\pi_{m,a}(x)^s
=
p_{m,a}^s\sum_x \sigma_m(x)^s
+(1-p_{m,a})^s\sum_x \eta_{m,a}(x)^s.
\]
由于 \(\sigma_m\) 在 \(A_m^{\max}\) 上均匀，故
\[
\sum_x \sigma_m(x)^s=K_m^{1-s}.
\]
又因 \(s>1\)，任意概率分布都满足 \(\sum_x \eta_{m,a}(x)^s\le 1\)。于是
\[
p_{m,a}^sK_m^{1-s}
\le
\sum_{x\in X_m}\pi_{m,a}(x)^s
\le
p_{m,a}^sK_m^{1-s}+(1-p_{m,a})^s.
\]
将上式除以 \(K_m^{1-s}\) 得
\[
p_{m,a}^s
\le
K_m^{s-1}\sum_{x\in X_m}\pi_{m,a}(x)^s
\le
p_{m,a}^s+K_m^{s-1}(1-p_{m,a})^s.
\]
因此只需验证
\[
K_m^{s-1}(1-p_{m,a})^s\to 0.
\]
由推论 \ref{cor:xi-fixed-freezing-escort-groundstate-tv-identity}，
\[
1-p_{m,a}\le \exp\!\bigl(-m\Delta(a)+o(m)\bigr),
\qquad
\Delta(a):=a(\alpha_*-\alpha_*^{(2)})-(\log\varphi-g_*)>0.
\]
另一方面，定理 \ref{thm:xi-fixed-freezing-power-sum-ratio-locking} 给出
\[
\log K_m=g_*m+o(m).
\]
故
\[
\log\!\Bigl(K_m^{s-1}(1-p_{m,a})^s\Bigr)
\le
m\bigl((s-1)g_*-s\Delta(a)\bigr)+o(m).
\]
当且仅当
\[
s\Delta(a)>(s-1)g_*
\]
时，右端趋于 \(-\infty\)。将 \(\Delta(a)\) 的定义代回并整理，恰得
\[
a>\frac{\log\varphi-\frac{1}{s}g_*}{\alpha_*-\alpha_*^{(2)}}=a_{\mathrm c,s}.
\]
于是
\[
K_m^{s-1}\sum_{x\in X_m}\pi_{m,a}(x)^s\to 1,
\]
即
\[
\sum_{x\in X_m}\pi_{m,a}(x)^s
=
K_m^{1-s}\bigl(1+o(1)\bigr).
\]
两边取对数并乘以 \((1-s)^{-1}\)，便得到
\[
H_s(\pi_{m,a})=\log K_m+o(1).
\]
证毕。
\end{proof}

\begin{theorem}[冻结相最大纤维扇区给出的预言机容量指数下界]\label{thm:xi-time-part65ca-frozen-max-sector-oracle-capacity-lower-bound}
沿用结论 \ref{con:xi-time-part57ba-freezing-escort-groundstate-fdiv-closed}、
推论 \ref{cor:xi-fixed-freezing-escort-groundstate-tv-identity} 与
定理 \ref{thm:xi-fixed-freezing-power-sum-ratio-locking} 的记号，并固定
\[
a>a_{\mathrm c}.
\]
记
\[
M_m:=\max_{x\in X_m}d_m(x),
\qquad
A_m^{\max}:=\{x\in X_m:\ d_m(x)=M_m\},
\qquad
K_m:=|A_m^{\max}|.
\]
若整数预算序列 \((B_m)_{m\ge 1}\) 满足
\[
B_m=\beta m+o(m)
\qquad(\beta\ge 0),
\]
则离散预言机容量
\[
C_m(B_m):=\sum_{x\in X_m}\min\bigl(d_m(x),2^{B_m}\bigr)
\]
满足严格指数下界
\[
\boxed{\;
\liminf_{m\to\infty}\frac{1}{m}\log C_m(B_m)
\ge
g_*+\min\{\alpha_*,\beta\log 2\}.\;}
\]
\end{theorem}
\begin{proof}
由定义，对每个 \(m\) 都有
\[
C_m(B_m)
\ge
\sum_{x\in A_m^{\max}}\min\bigl(M_m,2^{B_m}\bigr)
=
K_m\min\bigl(M_m,2^{B_m}\bigr).
\]
另一方面，定理 \ref{thm:xi-fixed-freezing-power-sum-ratio-locking} 已给出
\[
K_m=e^{g_*m+o(m)},
\qquad
M_m=e^{\alpha_*m+o(m)}.
\]
又由
\[
B_m=\beta m+o(m)
\]
可得
\[
2^{B_m}=e^{(\beta\log 2)m+o(m)}.
\]
故
\[
\min\bigl(M_m,2^{B_m}\bigr)
=
\exp\!\Bigl(
\min\{\alpha_*,\beta\log 2\}\,m+o(m)
\Bigr).
\]
将上述估计代回容量下界并取对数、除以 \(m\)、再取 \(\liminf\)，即得结论。
\end{proof}

\begin{corollary}[最大纤维扇区的侧信息临界斜率]\label{cor:xi-time-part65ca-frozen-max-sector-sideinfo-critical-slope}
沿用定理 \ref{thm:xi-time-part65ca-frozen-max-sector-oracle-capacity-lower-bound} 的记号，定义最优成功率
\[
\mathrm{Succ}_m(B):=\frac{C_m(B)}{2^m},
\]
以及最大纤维扇区单独贡献的成功率下界
\[
\mathrm{Succ}_m^{\max}(B)
:=
2^{-m}K_m\min\bigl(M_m,2^B\bigr)
\le
\mathrm{Succ}_m(B).
\]
则对满足 \(B_m=\beta m+o(m)\) 的任意整数预算序列，都有
\[
\boxed{\;
\liminf_{m\to\infty}\frac{1}{m}\log \mathrm{Succ}_m^{\max}(B_m)
\ge
-\bigl[\log 2-g_*-\min\{\alpha_*,\beta\log 2\}\bigr]_+,\;}
\]
其中 \([t]_+:=\max\{t,0\}\)。

特别地，若记
\[
\beta_c^{\max}
:=
\begin{cases}
\dfrac{\log 2-g_*}{\log 2},
& \log 2-g_*\le \alpha_*,\\[8pt]
+\infty,
& \log 2-g_*>\alpha_*,
\end{cases}
\]
则：
\begin{enumerate}
  \item 若 \(\beta>\beta_c^{\max}\)，则最大纤维扇区单独给出的成功率下界不再呈指数衰减；
  \item 若 \(\beta<\beta_c^{\max}\)，则该扇区的单独贡献仍保持指数级稀薄；
  \item 若 \(\log 2-g_*>\alpha_*\)，则即便侧信息按线性速率增长，最大纤维扇区本身也不足以闭合微态熵缺口。
\end{enumerate}
\end{corollary}
\begin{proof}
由定理 \ref{thm:xi-time-part65ca-frozen-max-sector-oracle-capacity-lower-bound}，
\[
\liminf_{m\to\infty}\frac{1}{m}\log \mathrm{Succ}_m^{\max}(B_m)
\ge
g_*+\min\{\alpha_*,\beta\log 2\}-\log 2.
\]
另一方面，由
\[
K_mM_m\le \sum_{x\in X_m}d_m(x)=2^m
\]
可得
\[
g_*+\alpha_*\le \log 2,
\]
故上式右端恒不为正，于是等价地可改写为所示的
\([\,\cdot\,]_+\) 形式。

再看临界斜率。若
\[
\log 2-g_*\le \alpha_*,
\]
则使
\[
g_*+\min\{\alpha_*,\beta\log 2\}-\log 2
\]
首次达到 \(0\) 的唯一有限值为
\[
\beta=\frac{\log 2-g_*}{\log 2}.
\]
若
\[
\log 2-g_*>\alpha_*,
\]
则对一切有限 \(\beta\) 都有
\[
g_*+\min\{\alpha_*,\beta\log 2\}-\log 2
\le
g_*+\alpha_*-\log 2<0,
\]
因而不存在有限临界斜率，只能记为 \(+\infty\)。三条解释正是这两种情形的直接改写。
\end{proof}

\endinput


\begin{theorem}[单原子 centered residue 扰动向量构成正则单纯形]\label{thm:xi-time-part65c-atomic-centered-residue-simplex}
\leanverified{paper\_xi\_time\_part65c\_atomic\_centered\_residue\_simplex}
沿用定理 \ref{thm:xi-atomic-witt-energy-residue-l2-law} 与结论 \ref{con:xi-time-part57ba-core-full-residue-rankone-perturbation} 的记号。固定模数 \(m\ge 2\)，并对每个 \(\ell\in \ZZ/m\ZZ\) 定义 centered residue 扰动向量
\[
\Delta^{(\ell)}\in\RR^{\ZZ/m\ZZ},
\qquad
\Delta^{(\ell)}_r:=2\,\mathbf 1_{\{r=\ell\}}-\frac{2}{m}
\qquad (r\in\ZZ/m\ZZ).
\]
等价地，若 even primitive/residue 注入发生在剩余类 \(r\equiv \ell\pmod m\)，则 \(\Delta^{(\ell)}\) 正是相应未中心化原子 \(2e_\ell\) 去掉全局平均 \(2m^{-1}\mathbf 1\) 后得到的 centered residue 差。
若记 \(\omega_m:=e^{2\pi i/m}\)，并采用离散 Fourier 变换
\[
\widehat x(j):=\sum_{r\in\ZZ/m\ZZ}x_r\,\omega_m^{jr},
\qquad j\in\ZZ/m\ZZ,
\]
则
\[
\widehat{\Delta^{(\ell)}}(0)=0,
\qquad
\widehat{\Delta^{(\ell)}}(j)=2\,\omega_m^{j\ell}
\quad (j\neq 0).
\]
换言之，centered residue 向量在非平凡 Fourier 模上恰给出等模相位轨道。

记零均值超平面
\[
\mathcal H_m:=\left\{x\in\RR^{\ZZ/m\ZZ}:\ \sum_{r\in\ZZ/m\ZZ}x_r=0\right\}.
\]
则成立：
\begin{enumerate}
  \item 对每个 \(\ell\) 都有 \(\Delta^{(\ell)}\in\mathcal H_m\)；
  \item 内积满足精确公式
  \[
  \boxed{\;
  \bigl\langle \Delta^{(\ell)},\Delta^{(\ell')}\bigr\rangle
  =
  4\left(\delta_{\ell,\ell'}-\frac{1}{m}\right).\;}
  \]
  \item 因而族 \(\{\Delta^{(\ell)}\}_{\ell\in\ZZ/m\ZZ}\) 在 \(\mathcal H_m\) 中构成一个边长恒定的正则单纯形顶点集；
  \item 它们张成整个 \((m-1)\)-维超平面 \(\mathcal H_m\)。
\end{enumerate}
\end{theorem}
\begin{proof}
由定义，
\[
\sum_{r\in\ZZ/m\ZZ}\Delta_r^{(\ell)}
=
2-\frac{2}{m}\cdot m
=
0,
\]
故 \(\Delta^{(\ell)}\in\mathcal H_m\)。

再计算内积。若 \(\ell=\ell'\)，则
\[
\bigl\|\Delta^{(\ell)}\bigr\|_2^2
=
\left(2-\frac{2}{m}\right)^2
+(m-1)\left(\frac{2}{m}\right)^2
=
4\left(1-\frac{1}{m}\right).
\]
若 \(\ell\neq \ell'\)，则恰有两个坐标分别取值
\[
2-\frac{2}{m}
\quad\text{与}\quad
-\frac{2}{m},
\]
其余 \(m-2\) 个坐标都等于 \(-2/m\)。因此
\[
\bigl\langle \Delta^{(\ell)},\Delta^{(\ell')}\bigr\rangle
=
2\left(2-\frac{2}{m}\right)\left(-\frac{2}{m}\right)
+(m-2)\left(\frac{2}{m}\right)^2
=
-\frac{4}{m}.
\]
合并两种情形，得到统一公式
\[
\bigl\langle \Delta^{(\ell)},\Delta^{(\ell')}\bigr\rangle
=
4\left(\delta_{\ell,\ell'}-\frac{1}{m}\right).
\]

于是全部向量具有相同范数、两两内积恒定，故在其仿射张成空间中构成正则单纯形。又
\[
\sum_{\ell\in\ZZ/m\ZZ}\Delta^{(\ell)}=0,
\]
故它们至多张成 \((m-1)\) 维空间。其 Gram 矩阵正是
\[
4\left(I-\frac{1}{m}J\right),
\]
而该矩阵的秩为 \(m-1\)。因此其线性张成空间维数恰为 \(m-1\)，并且只能等于同维的零均值超平面 \(\mathcal H_m\)。证毕。
\end{proof}

\begin{corollary}[单原子 centered residue 扰动在零均值子空间上的各向同性白噪声律]\label{cor:xi-time-part65c-atomic-centered-residue-isotropic}
\leanverified{paper\_xi\_time\_part65c\_atomic\_centered\_residue\_isotropic}
沿用定理 \ref{thm:xi-time-part65c-atomic-centered-residue-simplex} 的记号，定义均匀平均协方差算子
\[
\Sigma_m:=
\frac{1}{m}\sum_{\ell\in\ZZ/m\ZZ}
\Delta^{(\ell)}\bigl(\Delta^{(\ell)}\bigr)^{\mathsf T}.
\]
则有严格闭式
\[
\boxed{\;
\Sigma_m=\frac{4}{m}\left(I-\frac{1}{m}J\right),\;}
\]
其中 \(J\) 为全 \(1\) 矩阵。于是：
\begin{enumerate}
  \item \(\Sigma_m\) 在零均值超平面 \(\mathcal H_m\) 上恰等于
  \[
  \frac{4}{m}I_{\mathcal H_m},
  \]
  因而完全各向同性；
  \item 对任意 \(v\in\mathcal H_m\)，都有精确能量恒等式
  \[
  \boxed{\;
  \frac{1}{m}\sum_{\ell\in\ZZ/m\ZZ}
  \bigl|\langle v,\Delta^{(\ell)}\rangle\bigr|^2
  =
  \frac{4}{m}\|v\|_2^2.\;}
  \]
\end{enumerate}
换言之，single atom 在 centered residue 空间中的周期轨道平均不偏向任何方向；其二次统计贡献在 \(\mathcal H_m\) 上表现为一个完全各向同性的白噪声核。
\end{corollary}
\begin{proof}
由
\[
\Delta^{(\ell)}=2e_\ell-\frac{2}{m}\mathbf 1,
\]
可得
\[
\Sigma_m
=
\frac{1}{m}\sum_{\ell}
\left(2e_\ell-\frac{2}{m}\mathbf 1\right)
\left(2e_\ell-\frac{2}{m}\mathbf 1\right)^{\mathsf T}.
\]
逐项展开：
\[
\frac{1}{m}\sum_{\ell}4e_\ell e_\ell^{\mathsf T}=\frac{4}{m}I,
\]
\[
\frac{1}{m}\sum_{\ell}\left(
-\frac{4}{m}e_\ell\mathbf 1^{\mathsf T}
-\frac{4}{m}\mathbf 1 e_\ell^{\mathsf T}
\right)
=
-\frac{8}{m^2}J,
\]
\[
\frac{1}{m}\sum_{\ell}\frac{4}{m^2}\mathbf 1\mathbf 1^{\mathsf T}
=
\frac{4}{m^2}J.
\]
合并得
\[
\Sigma_m=\frac{4}{m}I-\frac{4}{m^2}J
=
\frac{4}{m}\left(I-\frac{1}{m}J\right).
\]
这便得到第一式。

若 \(v\in\mathcal H_m\)，则 \(Jv=0\)。因此
\[
v^{\mathsf T}\Sigma_m v
=
\frac{4}{m}v^{\mathsf T}v
=
\frac{4}{m}\|v\|_2^2.
\]
另一方面，
\[
v^{\mathsf T}\Sigma_m v
=
\frac{1}{m}\sum_{\ell}
\bigl|\langle v,\Delta^{(\ell)}\rangle\bigr|^2.
\]
两式合并即得所示能量恒等式。证毕。
\end{proof}

\begin{theorem}[归一化 periodic residue law 偏离均匀分布的精确指数]\label{thm:xi-time-part65c-normalized-residue-law-exact-exponent}
沿用结论 \ref{con:xi-time-part57ba-core-full-residue-rankone-perturbation} 与推论 \ref{cor:xi-single-atom-centered-residue-l2-exponent-truncation} 的记号。固定模数 \(m\ge 2\)，定义归一化 periodic residue law
\[
\mu_{m,r}^{\mathrm{full}}(n)
:=
\frac{A_{m,r}^{\mathrm{full}}(n)}{a_n^{\mathrm{full}}(1)},
\qquad
\mu_{m,\bullet}^{\mathrm{full}}(n)
:=
\bigl(\mu_{m,r}^{\mathrm{full}}(n)\bigr)_{r\in\ZZ/m\ZZ},
\qquad
u_m:=\frac{1}{m}\mathbf 1_m.
\]
\leanverified{paper\_xi\_time\_part65c\_normalized\_residue\_law\_exact\_exponent}
再记
\[
\lambda:=\rho(M(1))=\varphi^2,
\qquad
\rho_m^{\mathrm{core}}
:=
\max_{1\le j\le m-1}\rho\!\bigl(M_{\mathrm{core}}(\omega_m^j)\bigr).
\]
则对 \(\CC^{\ZZ/m\ZZ}\) 上任意固定向量范数 \(\|\cdot\|\)，都有
\[
\boxed{\;
\limsup_{n\to\infty}
\bigl\|
\mu_{m,\bullet}^{\mathrm{full}}(n)-u_m
\bigr\|^{1/n}
=
\frac{\max\{1,\rho_m^{\mathrm{core}}\}}{\lambda}.\;}
\]
\end{theorem}
\begin{proof}
记 centered residue 向量
\[
\widetilde A_{m,\bullet}^{\mathrm{full}}(n)
:=
\Bigl(
A_{m,r}^{\mathrm{full}}(n)-\frac{a_n^{\mathrm{full}}(1)}{m}
\Bigr)_{r\in\ZZ/m\ZZ}.
\]
则
\[
\mu_{m,\bullet}^{\mathrm{full}}(n)-u_m
=
\frac{\widetilde A_{m,\bullet}^{\mathrm{full}}(n)}{a_n^{\mathrm{full}}(1)}.
\]
由推论 \ref{cor:xi-single-atom-centered-residue-l2-exponent-truncation}，
\[
\limsup_{n\to\infty}
\bigl\|\widetilde A_{m,\bullet}^{\mathrm{full}}(n)\bigr\|_2^{1/n}
=
\max\{1,\rho_m^{\mathrm{core}}\}.
\]
由于 \(\CC^{\ZZ/m\ZZ}\) 有限维，任意两个固定范数等价，故上式中的指数不依赖于范数选择，从而
\[
\limsup_{n\to\infty}
\bigl\|\widetilde A_{m,\bullet}^{\mathrm{full}}(n)\bigr\|^{1/n}
=
\max\{1,\rho_m^{\mathrm{core}}\}.
\]
另一方面，无扭曲 periodic 总数满足
\[
\lim_{n\to\infty}\bigl(a_n^{\mathrm{full}}(1)\bigr)^{1/n}
=
\rho(M(1))
=
\lambda,
\]
这是无扭曲 Perron 根的直接结果。将上述两式相除即可得到题述结论。证毕。
\end{proof}

\begin{theorem}[原子主导区中归一化 residue law 的单纯形首项轨道]\label{thm:xi-time-part65c-atom-dominated-normalized-residue-profile}
沿用定理 \ref{thm:xi-time-part65c-atomic-centered-residue-simplex} 与定理 \ref{thm:xi-time-part65c-normalized-residue-law-exact-exponent} 的记号。若
\[
\rho_m^{\mathrm{core}}<1,
\]
则对偶数 \(n=2q\)，记 \(r_q\in\ZZ/m\ZZ\) 为 \(q\) 的剩余类，有
\[
\boxed{\;
\frac{a_{2q}^{\mathrm{full}}(1)}{2}
\bigl(\mu_{m,\bullet}^{\mathrm{full}}(2q)-u_m\bigr)
=
e_{r_q}-u_m+o(1)
=
\frac12\Delta^{(r_q)}+o(1),\;}
\]
而对奇数 \(n\) 有
\[
\boxed{\;
a_n^{\mathrm{full}}(1)\bigl(\mu_{m,\bullet}^{\mathrm{full}}(n)-u_m\bigr)=o(1).\;}
\]
因此在 \(\rho_m^{\mathrm{core}}<1\) 的原子主导区中，归一化 residue law 的首项偏差并非任意零均值向量，而是沿定理 \ref{thm:xi-time-part65c-atomic-centered-residue-simplex} 的正则单纯形顶点轨道
\[
\{e_r-u_m:\ r\in\ZZ/m\ZZ\}
\]
周期移动。
\end{theorem}
\begin{proof}
由结论 \ref{con:xi-time-part57ba-core-full-residue-rankone-perturbation}，对偶数 \(n=2q\) 有
\[
A_{m,r}^{\mathrm{full}}(2q)
=
A_{m,r}^{\mathrm{core}}(2q)+2\,\mathbf 1_{\{r\equiv q\;(\mathrm{mod}\,m)\}},
\qquad
a_{2q}^{\mathrm{full}}(1)=a_{2q}^{\mathrm{core}}(1)+2.
\]
因此
\[
\mu_{m,\bullet}^{\mathrm{full}}(2q)-u_m
=
\frac{1}{a_{2q}^{\mathrm{full}}(1)}
\Bigl(
\widetilde A_{m,\bullet}^{\mathrm{core}}(2q)+2(e_{r_q}-u_m)
\Bigr).
\]
另一方面，由 \(\rho_m^{\mathrm{core}}<1\) 与推论 \ref{cor:xi-single-atom-centered-residue-l2-exponent-truncation} 的第一条，
\[
\bigl\|\widetilde A_{m,\bullet}^{\mathrm{core}}(n)\bigr\|_2=o(1),
\]
故由范数等价亦有
\[
\bigl\|\widetilde A_{m,\bullet}^{\mathrm{core}}(n)\bigr\|=o(1)
\]
在任意固定范数下成立。把上式乘以 \(a_{2q}^{\mathrm{full}}(1)/2\) 即得偶数情形。

当 \(n\) 为奇数时，单原子 residue 修正消失，故
\[
\mu_{m,\bullet}^{\mathrm{full}}(n)-u_m
=
\frac{\widetilde A_{m,\bullet}^{\mathrm{core}}(n)}{a_n^{\mathrm{full}}(1)}.
\]
再乘以 \(a_n^{\mathrm{full}}(1)\) 并利用同一 \(o(1)\) 估计，即得奇数情形。最后，由定理 \ref{thm:xi-time-part65c-atomic-centered-residue-simplex} 中
\(
\Delta^{(r)}=2(e_r-u_m)
\)
的定义，首项轨道正是同一定理给出的正则单纯形顶点轨道。证毕。
\end{proof}

\begin{theorem}[输出计数中心窗口的一阶 Edgeworth 修正与严格左偏]\label{thm:xi-time-part65c-output-edgeworth-first-correction}
\leanverified{paper\_xi\_time\_part65c\_output\_edgeworth\_first\_correction}
沿用命题 \ref{prop:real-input-40-output-cumulants-closed} 与定理 \ref{thm:xi-real-input-40-output-edgeworth-quasipowers} 的记号。记
\[
\alpha_0:=P'(0)=0.367376207875073606\ldots,
\qquad
\sigma^2:=P''(0)=0.0487412456183376392\ldots,
\]
\[
\kappa_3:=P^{(3)}(0)=-0.0158881374258564278\ldots<0,
\qquad
\kappa_4:=P^{(4)}(0)=0.00568478997567274949\ldots>0.
\]
若 \(X_n\) 为定理 \ref{thm:xi-real-input-40-output-edgeworth-quasipowers} 中的输出 \(1\) 总数，并定义标准化变量
\[
Y_n:=\frac{X_n-n\alpha_0}{\sigma\sqrt n},
\]
则在中心窗口内一致有
\[
\boxed{\;
\PP(Y_n\le x)
=
\Phi(x)
+\frac{\kappa_3}{6\sigma^3\sqrt n}(1-x^2)\phi(x)
+O(n^{-1}).\;}
\]
特别地，标准化偏度常数
\[
\boxed{\;
\gamma_1:=\frac{\kappa_3}{\sigma^3}<0.\;}
\]
因此输出计数在一阶 Edgeworth 层面严格左偏。
\end{theorem}
\begin{proof}
定理 \ref{thm:xi-real-input-40-output-edgeworth-quasipowers} 已给出
\[
\PP\!\left(
\frac{X_n-n\mu}{\sqrt{n\kappa_2}}
\le x
\right)
=
\Phi(x)
+\frac{\kappa_3}{6\kappa_2^{3/2}\sqrt n}(1-x^2)\phi(x)
+O(n^{-1})
\]
在中心窗口内一致成立。将其中 \(\mu,\kappa_2\) 分别改写为 \(\alpha_0,\sigma^2\) 即得第一式。

再由定理 \ref{thm:xi-real-input-40-output-edgeworth-nongaussian} 与表 \ref{tab:real_input_40_output_potential_cumulants_closed}，
\[
\kappa_3<0,
\]
从而 \(\gamma_1=\kappa_3/\sigma^3<0\)。证毕。
\end{proof}

\begin{corollary}[Edgeworth 首项与局部大偏差三次项的同源符号锁定]\label{cor:xi-time-part65c-edgeworth-local-ldp-sign-lock}
沿用定理 \ref{thm:xi-time-part65c-output-edgeworth-first-correction} 与定理 \ref{thm:xi-output-potential-local-legendre-birkhoff-normal-form} 的记号。则
\[
\operatorname{sign}\!\left(\frac{\kappa_3}{6\sigma^3}\right)
=
\operatorname{sign}\!\left(\frac{\kappa_3}{6\sigma^6}\right)
=
\operatorname{sign}(\kappa_3).
\]
等价地，中心极限层的一阶 Edgeworth 偏斜与局部 Legendre 正规形中的三次偏斜常数具有同一符号源。对当前核，由 \(\kappa_3<0\) 可得
\[
I_{\mathrm{out}}(\alpha_0+\delta)-I_{\mathrm{out}}(\alpha_0-\delta)
=
-\frac{\kappa_3}{3\sigma^6}\delta^3+O(\delta^5)>0
\qquad (\delta>0\ \text{充分小}),
\]
故中心窗口严格左偏，而局部大偏差层的右偏离严格更昂贵。
\end{corollary}
\begin{proof}
由定理 \ref{thm:xi-time-part65c-output-edgeworth-first-correction}，一阶 Edgeworth 系数为 \(\kappa_3/(6\sigma^3)\)。由定理 \ref{thm:xi-output-potential-local-legendre-birkhoff-normal-form}，
\[
I_{\mathrm{out}}(\alpha_0+\delta)
=
\frac{\delta^2}{2\sigma^2}
-
\left(\frac{\kappa_3}{6\sigma^6}\right)\delta^3
+O(\delta^4).
\]
因此控制两种首阶奇次不对称的常数分别为 \(\kappa_3/(6\sigma^3)\) 与 \(\kappa_3/(6\sigma^6)\)，它们显然具有相同符号。再对 \(\delta\) 与 \(-\delta\) 相减，即得
\[
I_{\mathrm{out}}(\alpha_0+\delta)-I_{\mathrm{out}}(\alpha_0-\delta)
=
-\frac{\kappa_3}{3\sigma^6}\delta^3+O(\delta^5).
\]
由 \(\kappa_3<0\) 立得题述不等式。证毕。
\end{proof}

\noindent
以上诸结果把冻结相 pressure/escort 链、单原子 residue 手术与输出位势的首阶非高斯修正进一步压缩到同一层：一方面，escort--Rényi 熵率不仅在 \((a,s)\)-平面上沿 \(as=a_{\mathrm c}\) 发生严格冻结，而且在更强半直线 \(a>a_{\mathrm c,s}\) 上把有限尺度主项提升为 \(\log K_m\)，并把最大纤维扇区对预言机容量的贡献压缩为指数下界 \(g_*+\min\{\alpha_*,\beta\log 2\}\) 及其相应侧信息临界斜率；另一方面，单原子的一维局域手术不仅在 centered residue 空间中生成正则单纯形，而且把归一化 periodic residue law 的指数精确锁定为 \(\max\{1,\rho_m^{\mathrm{core}}\}/\lambda\)，并在 \(\rho_m^{\mathrm{core}}<1\) 时强迫首项沿该单纯形顶点轨道移动；最后，输出计数在中心窗口中的左偏 Edgeworth 修正与局部大偏差的右侧代价增强又由同一三阶累积量 \(\kappa_3\) 严格统一。由此，冻结热力学、原子局域手术与首阶非高斯非对称便在“冻结边界 + 有限尺度主项 + 侧信息斜率 + 单纯形轨道 + 偏斜常数”的共同语言下闭合为同一条可审计主线。

\endinput


\paragraph{宏观共振模、周期化绝对连续障碍与熵缺口尺度失配}
在 Fibonacci/Lucas 共振子列的正极限、Pisot--Bernoulli 周期化障碍、最大纤维偏置、谱零点稀疏律以及碰撞/KL 恒等式已经建立之后，还可进一步压缩出下列六条彼此衔接的直接后果。它们分别刻画宏观共振模的对数丰度、周期化测度的非 Rajchman 与绝对连续性排斥、最大纤维相对均值的无界放大、相对均匀基线的 \(\sqrt{M}\)-尺度偏离刚性、零点相对密度对碰撞尺度的根本失效，以及 Shannon 熵缺口在自然 \(1/M\) 标度上的发散。

\begin{theorem}[宏观共振模的对数丰度]\label{thm:xi-time-part66-macro-resonance-log-abundance}
沿用定理 \ref{thm:fold-critical-resonance-constant-integer-ladder} 与推论
\ref{cor:fold-resonance-ladder-fib-luc-limits} 的记号，记
\[
c_{\mathrm{Fib}}:=\lim_{n\to\infty}C_\varphi(F_n)>0.
\]
对任意常数
\[
0<\eta<c_{\mathrm{Fib}},
\]
存在 \(N_\eta\) 与 \(c_\eta>0\)，使得对所有 \(N\ge N_\eta\) 都有
\[
\boxed{\;
\#\{1\le u\le N:\ C_\varphi(u)\ge \eta\}
\ge
c_\eta\log N.\;}
\]
因此，非零共振峰并非孤立稀见事件，而是在对数尺度上持续出现的宏观族。
\end{theorem}
\begin{proof}
由推论 \ref{cor:fold-resonance-ladder-fib-luc-limits}，
\[
C_\varphi(F_n)\longrightarrow c_{\mathrm{Fib}}>0.
\]
故对任意 \(0<\eta<c_{\mathrm{Fib}}\)，存在 \(n_0\) 使得对所有 \(n\ge n_0\) 都有
\[
C_\varphi(F_n)\ge \eta.
\]
于是当 \(N\) 充分大时，
\[
\#\{1\le u\le N:\ C_\varphi(u)\ge \eta\}
\ge
\#\{n\ge n_0:\ F_n\le N\}.
\]
再由 Binet 公式，
\[
\#\{n:\ F_n\le N\}
=
\frac{\log N}{\log\varphi}+O(1).
\]
因此存在常数 \(c_\eta>0\) 与 \(N_\eta\)，使得对所有 \(N\ge N_\eta\) 都有
\[
\#\{1\le u\le N:\ C_\varphi(u)\ge \eta\}\ge c_\eta\log N.
\]
证毕。
\end{proof}

\begin{corollary}[\texorpdfstring{$\mu_{\varphi^{-1}}^{(2)}$}{mu} 的模 \texorpdfstring{$2$}{2} 周期化既非 Rajchman，亦无任何 \texorpdfstring{$L^p$}{Lp} 密度]\label{cor:xi-time-part66-periodized-bernoulli-nonrajchman-no-lp}
令
\[
q:=\varphi^{-1},
\]
并令 \(\mu_q^{(2)}\) 如定理 \ref{thm:fold-pisot-bernoulli-convolution-representation}。记
\[
\nu:=\bigl(\RR\to \mathbb T_2=\RR/(2\ZZ)\bigr)_\#\mu_q^{(2)}
\]
为其模 \(2\) 周期化测度。则：
\begin{enumerate}
  \item \(\nu\) 的 Fourier 系数不趋于 \(0\)，故 \(\nu\) 不是 Rajchman 测度；
  \item \(\nu\) 不可能具有任何 \(L^p(\mathbb T_2)\) 密度，其中 \(1\le p\le \infty\)。
\end{enumerate}
\end{corollary}
\begin{proof}
由推论 \ref{cor:fold-resonance-ladder-fib-luc-limits}，记
\[
c_{\mathrm{Fib}}:=\lim_{n\to\infty}C_\varphi(F_n)>0.
\]
再由定理 \ref{thm:xi-pisot-bernoulli-fibonacci-direction-nonrajchman}，
\[
\bigl|\widehat\nu(F_n)\bigr|
=
\bigl|\widehat{\mu_q^{(2)}}(\pi F_n)\bigr|
=
C_\varphi(F_n)
\longrightarrow
c_{\mathrm{Fib}}>0,
\]
故 \(\widehat\nu(u)\not\to 0\)，从而 \(\nu\) 不是 Rajchman 测度。

若 \(\nu\) 具有某个 \(L^1(\mathbb T_2)\) 密度，则由 Riemann--Lebesgue 引理，其 Fourier 系数必趋于 \(0\)，与上式矛盾。另一方面，\(\mathbb T_2\) 的 Haar 测度有限，故对任意 \(1\le p\le\infty\) 都有
\[
L^p(\mathbb T_2)\subseteq L^1(\mathbb T_2).
\]
因此 \(\nu\) 不可能具有任何 \(L^p\) 密度。证毕。
\end{proof}
\leanverified{paper\_xi\_time\_part66\_periodized\_bernoulli\_nonrajchman\_no\_lp}

\begin{theorem}[保 Fibonacci 共振的 Fourier 乘子不可能把周期化测度送入 \texorpdfstring{$L^2$}{L2}]\label{thm:xi-time-part66-fibonacci-visible-multiplier-no-l2}
沿用推论 \ref{cor:xi-time-part66-periodized-bernoulli-nonrajchman-no-lp} 的记号。设
\[
\sigma:\ZZ\to \CC
\]
为任意符号，并把 \(T_\sigma\nu\) 视为 Fourier 系数满足
\[
\widehat{T_\sigma\nu}(u)=\sigma(u)\widehat\nu(u)
\qquad (u\in\ZZ)
\]
的周期分布。若存在常数 \(c>0\) 与 \(n_0\)，使得
\[
|\sigma(F_n)|\ge c
\qquad (n\ge n_0),
\]
则 \(T_\sigma\nu\) 不可能由任何 \(L^2(\mathbb T_2)\) 密度表示。
\end{theorem}
\begin{proof}
由推论 \ref{cor:xi-time-part66-periodized-bernoulli-nonrajchman-no-lp} 的证明，
\[
|\widehat\nu(F_n)|=C_\varphi(F_n)\longrightarrow c_{\mathrm{Fib}}>0.
\]
故存在 \(n_1\) 使得对一切 \(n\ge n_1\) 都有
\[
|\widehat\nu(F_n)|\ge \frac{c_{\mathrm{Fib}}}{2}.
\]
于是对任意 \(N\ge \max\{n_0,n_1\}\)，都有
\[
\sum_{u\in\ZZ}\bigl|\widehat{T_\sigma\nu}(u)\bigr|^2
\ge
\sum_{n\ge N}\bigl|\sigma(F_n)\widehat\nu(F_n)\bigr|^2
\ge
\frac{c^2c_{\mathrm{Fib}}^2}{4}\sum_{n\ge N}1
=+\infty.
\]
若 \(T_\sigma\nu\) 由某个 \(f\in L^2(\mathbb T_2)\) 表示，则 Parseval 恒等式给出
\[
\sum_{u\in\ZZ}\bigl|\widehat{T_\sigma\nu}(u)\bigr|^2
=
\sum_{u\in\ZZ}|\widehat f(u)|^2
<
\infty,
\]
矛盾。故 \(T_\sigma\nu\) 不可能具有 \(L^2(\mathbb T_2)\) 密度。证毕。
\end{proof}

\begin{theorem}[最大纤维相对均值的放大因子无界]\label{thm:xi-time-part66-maxfiber-over-mean-diverges}
\leanverified{paper\_xi\_time\_part66\_maxfiber\_over\_mean\_diverges}
沿用定理 \ref{thm:xi-time-part9m-maxfiber-uniformity-gap-chain} 与推论
\ref{cor:fold-collision-scale-diverges} 的记号。记
\[
M:=F_{m+2},
\qquad
p_m(r):=\frac{d_m(r)}{2^m},
\qquad
M_m:=\max_r d_m(r),
\qquad
\overline d_m:=\frac{2^m}{M}.
\]
则对每个 \(m\) 都有
\[
\boxed{\;
\frac{M_m}{\overline d_m}\ge M\,\mathsf{Col}_m.\;}
\]
从而
\[
\boxed{\;
\frac{M_m}{\overline d_m}\longrightarrow +\infty.\;}
\]
换言之，最大纤维相对均值的偏置本身在尺度极限下发散。
\end{theorem}
\begin{proof}
设
\[
p_{\max}:=\max_r p_m(r)=\frac{M_m}{2^m}.
\]
由 \(\sum_r p_m(r)=1\) 可知
\[
\mathsf{Col}_m=\sum_r p_m(r)^2\le p_{\max}\sum_r p_m(r)=p_{\max}.
\]
两边乘以 \(M\)，得到
\[
M\,\mathsf{Col}_m\le M p_{\max}=M\frac{M_m}{2^m}=\frac{M_m}{\overline d_m}.
\]
这就证明了第一条不等式。

再由推论 \ref{cor:fold-collision-scale-diverges}，
\[
M\,\mathsf{Col}_m=F_{m+2}\mathsf{Col}_m\longrightarrow +\infty.
\]
结合上述下界，立即得到
\[
\frac{M_m}{\overline d_m}\longrightarrow +\infty.
\]
证毕。
\end{proof}

\begin{theorem}[零点相对密度永远不是碰撞尺度的充分统计量]\label{thm:xi-time-part66-zero-density-never-sufficient}
记
\[
M:=F_{m+2},
\qquad
\eta_m:=\frac{|\mathcal Z_m|}{M},
\]
其中 \(\mathcal Z_m\subset \ZZ/(M\ZZ)\) 为谱零点集合。设
\[
\Phi:[0,1]\to [0,\infty)
\]
在 \(0\) 的某邻域内有界。则存在 \(m_0\)，使得对所有 \(m\ge m_0\) 都有
\[
\boxed{\;
M\,\mathsf{Col}_m>\Phi(\eta_m).\;}
\]
因此，任何仅依赖零点相对密度并在 \(0\) 邻域有界的零点几何指标，都不可能控制归一化碰撞尺度。
\end{theorem}
\begin{proof}
先证 \(\eta_m\to 0\)。若 \(3\nmid (m+2)\)，则由定理
\ref{thm:killo-fold-zero-spectrum-affine-congruence}，
\[
\mathcal Z_m=\varnothing,
\qquad
\eta_m=0.
\]
若 \(3\mid (m+2)\)，则 \(M\) 为偶数，定理 \ref{thm:fold-zero-density-sparse} 给出
\[
\eta_m=\frac{|\mathcal Z_m|}{M}
\le
\tau(m+2)\,\frac{F_{\lfloor (m+2)/2\rfloor}}{F_{m+2}}
=
\tau(m+2)\,O(\varphi^{-m/2})\to 0.
\]
故在全体 \(m\) 上都有 \(\eta_m\to 0\)。

由于 \(\Phi\) 在 \(0\) 的某邻域内有界，存在 \(\delta>0\) 与 \(K<\infty\)，使得
\[
0\le t\le \delta\quad\Longrightarrow\quad \Phi(t)\le K.
\]
由 \(\eta_m\to 0\)，存在 \(m_1\) 使得对所有 \(m\ge m_1\) 都有 \(\eta_m\le \delta\)，故
\[
\Phi(\eta_m)\le K.
\]
另一方面，由推论 \ref{cor:fold-collision-scale-diverges}，
\[
M\,\mathsf{Col}_m\longrightarrow +\infty.
\]
故存在 \(m_2\) 使得对所有 \(m\ge m_2\) 都有
\[
M\,\mathsf{Col}_m>K.
\]
取 \(m_0:=\max\{m_1,m_2\}\)，便得
\[
M\,\mathsf{Col}_m>\Phi(\eta_m)
\qquad (m\ge m_0).
\]
证毕。
\end{proof}

\begin{theorem}[相对均匀基线的 \texorpdfstring{$\sqrt{M}$}{sqrtM} 尺度偏离不可消除]\label{thm:xi-time-part66-uniform-baseline-sqrtM-rigidity}
沿用结论 \ref{con:xi-fold-collision-spectral-trace-parseval} 与结论
\ref{con:xi-fold-chi2-constant-residual-energy} 的记号。记
\[
M:=F_{m+2},
\qquad
p_m(r):=\frac{d_m(r)}{2^m},
\qquad
u(r):=\frac{1}{M}
\qquad
\bigl(r\in \ZZ/(M\ZZ)\bigr).
\]
则有
\[
\boxed{\;
\liminf_{m\to\infty}\sqrt{M}\,\|p_m-u\|_{\ell^2}
\ge
\sqrt{2}\,C_\varphi.\;}
\]
以及
\[
\boxed{\;
\liminf_{m\to\infty}\sqrt{M}\,\|p_m-u\|_{\mathrm{TV}}
\ge
\frac{C_\varphi}{\sqrt{2}}.\;}
\]
因此，折叠推前分布相对均匀基线的偏离并非 \(o(M^{-1/2})\) 级扰动，而在中心极限定标下保留严格正的黄金常数下界。
\end{theorem}
\begin{proof}
由结论 \ref{con:xi-fold-collision-spectral-trace-parseval}，
\[
\|p_m-u\|_{\ell^2}^2
:=
\sum_{r}\bigl(p_m(r)-u(r)\bigr)^2
=
\mathsf{Col}_m-\frac1M.
\]
另一方面，结论 \ref{con:xi-fold-chi2-constant-residual-energy} 给出
\[
\liminf_{m\to\infty}M\Bigl(\mathsf{Col}_m-\frac1M\Bigr)\ge 2C_\varphi^2.
\]
两式合并即得
\[
\liminf_{m\to\infty}M\,\|p_m-u\|_{\ell^2}^2\ge 2C_\varphi^2,
\]
从而
\[
\liminf_{m\to\infty}\sqrt{M}\,\|p_m-u\|_{\ell^2}\ge \sqrt2\,C_\varphi.
\]
再由
\[
\|p_m-u\|_{\mathrm{TV}}
=
\frac12\|p_m-u\|_{\ell^1}
\ge
\frac12\|p_m-u\|_{\ell^2},
\]
便得到
\[
\liminf_{m\to\infty}\sqrt{M}\,\|p_m-u\|_{\mathrm{TV}}
\ge
\frac{1}{2}\cdot \sqrt2\,C_\varphi
=
\frac{C_\varphi}{\sqrt2}.
\]
证毕。
\end{proof}

\begin{theorem}[Shannon 熵缺口在 \texorpdfstring{$1/M$}{1/M} 标度上发散]\label{thm:xi-time-part66-shannon-gap-times-M-diverges}
\leanverified{paper\_xi\_time\_part66\_shannon\_gap\_times\_m\_diverges}
沿用结论 \ref{con:xi-fold-collision-spectral-trace-parseval} 与推论
\ref{cor:fold-collision-scale-diverges} 的记号。记
\[
M:=F_{m+2},
\qquad
p_m(r):=\frac{d_m(r)}{2^m},
\qquad
u(r):=\frac{1}{M}.
\]
则有下界
\[
\boxed{\;
D_{\mathrm{KL}}(p_m\|u)\ge \frac12\Bigl(\mathsf{Col}_m-\frac1M\Bigr).\;}
\]
特别地，
\[
\boxed{\;
M\bigl(\log M-H(p_m)\bigr)
\ge
\frac12\bigl(M\,\mathsf{Col}_m-1\bigr)
\longrightarrow +\infty.\;}
\]
因此，相对均匀分布的 Shannon 熵缺口乘以自然尺度 \(M\) 后并不趋于有限常数，而是发散到无穷。
\end{theorem}
\begin{proof}
由结论 \ref{con:xi-fold-collision-spectral-trace-parseval}，
\[
\sum_r\bigl(p_m(r)-u(r)\bigr)^2
=
\mathsf{Col}_m-\frac1M.
\]
另一方面，Pinsker 不等式给出
\[
D_{\mathrm{KL}}(p_m\|u)\ge 2\,\|p_m-u\|_{\mathrm{TV}}^2.
\]
而
\[
\|p_m-u\|_{\mathrm{TV}}
=
\frac12\|p_m-u\|_{\ell^1}
\ge
\frac12\|p_m-u\|_{\ell^2},
\]
故
\[
D_{\mathrm{KL}}(p_m\|u)
\ge
2\cdot \frac14\|p_m-u\|_{\ell^2}^2
=
\frac12\sum_r\bigl(p_m(r)-u(r)\bigr)^2
=
\frac12\Bigl(\mathsf{Col}_m-\frac1M\Bigr).
\]
再由
\[
D_{\mathrm{KL}}(p_m\|u)=\log M-H(p_m),
\]
两边乘以 \(M\)，得到
\[
M\bigl(\log M-H(p_m)\bigr)
\ge
\frac12\bigl(M\,\mathsf{Col}_m-1\bigr).
\]
最后由推论 \ref{cor:fold-collision-scale-diverges}，\(M\,\mathsf{Col}_m\to+\infty\)，故右端趋于 \(+\infty\)。证毕。
\end{proof}

\noindent
这六条结果把共振闭包常数、周期化 Fourier 采样、谱零点集合、碰撞尺度与 Shannon 熵缺口几条主线真正闭合起来：一方面，Fibonacci 共振并非只提供个别正极限，而在阈值意义下形成对数丰度的宏观模族，并迫使周期化测度同时失去 Rajchman 性与全部绝对连续可能；另一方面，碰撞尺度的发散不仅会被最大纤维相对均值直接放大到无界量级，还表明相对均匀基线的 \(\ell^2\) 与总变差偏离在 \(\sqrt{M}\) 定标下保留严格正下界，并且任何局部有界的零点密度指标都不足以控制归一化碰撞强度；最终，这一非均匀性又在信息论层面外显为 \(M(\log M-H(p_m))\to+\infty\) 的熵缺口发散。由此，有限窗口折叠的非均匀性便不再能被解释为少数零点位置的几何副产物，而必须被视为一族对数丰度的宏观共振模在振幅、峰值与熵三重层面上的共同驱动。

\paragraph{角色投影、平滑塌缩与外置算术核的函子分裂}
在模剩余 Fourier 反演、window-\(6\) 群商障碍、末位充分化、输出位势的解析腔室刚性、二步原子的 primitive/finite-part 位移、零温四态地基层分裂以及局部化数系的 Pontryagin 对偶链已经建立之后，还可把离散波纹与连续平滑之间的关系进一步压缩为下列七条彼此闭合的结论。它们分别刻画非平凡角色扇区、非自由群商障碍、末位二点终端塌缩、单解析相内的波滑分裂、两步原子的 Abel 有限部位移、零温平滑包络的有限状态化，以及连续相位层之后仍不可消去的 profinite 素数核。

\begin{theorem}[角色投影的扭曲扇区分裂]\label{thm:xi-time-part67-character-projection-twisted-sectors}
\leanverified{paper\_xi\_time\_part67\_character\_projection\_twisted\_sectors}
固定整数 \(q\ge 2\)。设 \(B_{q,a}(n)\) 为定义在 \(a\in \ZZ/(q\ZZ)\) 上的计数函数，并令
\[
\widehat B_{q,j}(n):=\sum_{a\in \ZZ/(q\ZZ)}B_{q,a}(n)\,\omega_q^{ja},
\qquad
\omega_q:=e^{2\pi i/q}.
\]
则有 Fourier 反演
\[
B_{q,a}(n)=\frac1q\sum_{j=0}^{q-1}\widehat B_{q,j}(n)\,\omega_q^{-ja},
\qquad
\frac1q\sum_{b\in \ZZ/(q\ZZ)}B_{q,b}(n)=\frac1q\,\widehat B_{q,0}(n),
\]
因而中心化波纹恰为
\[
B_{q,a}(n)-\frac1q\,\widehat B_{q,0}(n)
=
\frac1q\sum_{j=1}^{q-1}\widehat B_{q,j}(n)\,\omega_q^{-ja}.
\]
对真实输入 \(40\) 状态核取 \(B_{q,a}(n)=A_{q,a}(n)\) 或 \(B_{q,a}(n)=P_{q,a}(n)\)，则
\[
\widehat A_{q,j}(n)=a_n(\omega_q^j),
\qquad
\widehat P_{q,j}(n)=p_n(\omega_q^j),
\]
从而平滑均值恰是平凡角色扇区，而全部离散波纹严格由非平凡单位根扭曲 \(M(\omega_q^j)\) 承载。进一步，若记
\[
\rho_q:=\max_{1\le j\le q-1}\rho\!\bigl(M(\omega_q^j)\bigr),
\qquad
\theta_q:=\frac{\log\rho_q}{\log\lambda},
\qquad
\lambda=\varphi^2,
\]
则这些非平凡扇区在 periodic 与 primitive 两层均以 \(O(\rho_q^n)\) 衰减；并且存在明确模数 \(q\) 使 \(\theta_q>\frac12\)，且该性质沿倍数塔向上遗传。
\end{theorem}
\begin{proof}
有限循环群 \(\ZZ/(q\ZZ)\) 上的离散 Fourier 反演直接给出前两式；中心化公式只需把 \(j=0\) 项单独提出。

对真实输入 \(40\) 状态核，定理 \ref{thm:xi-dirichlet-residue-fourier-parseval-exact} 已在 periodic 与 primitive 两层分别给出
\[
a_n(\omega_q^j)=\sum_{a\in \ZZ/(q\ZZ)}A_{q,a}(n)\,\omega_q^{ja},
\qquad
p_n(\omega_q^j)=\sum_{a\in \ZZ/(q\ZZ)}P_{q,a}(n)\,\omega_q^{ja},
\]
故扭曲 trace 正是 residue counts 的非平凡角色分量。命题 \ref{prop:real-input-40-output-dirichlet-nonrh} 进一步给出：这些扭曲分量在 periodic/primitive 两层都以 \(O(\rho_q^n)\) 衰减；并且存在明确模数使 \(\theta_q>\frac12\)，其中 \(q=2\) 已足够。最后，由结论 \ref{con:xi-time-part9l-dirichlet-twist-divisibility-monotonicity}，一旦某个模数满足 \(\theta_q>\frac12\)，则其全部倍数也都满足同一严格不等式。证毕。
\end{proof}

\begin{theorem}[window-\texorpdfstring{$6$}{6} 平滑包络的非自由群商障碍]\label{thm:xi-time-part67-window6-nonfree-group-quotient-obstruction}
设
\[
\Omega_6:=\{0,1\}^6,
\qquad
F_6:=\Fold_6^{\mathrm{bin}}:\Omega_6\to X_6.
\]
若试图把 \(F_6\) 的纤维压缩解释为某个固定有限群的轨道商，或解释为某个阿贝尔线性陪集划分，则该方案必失败。更具体地，\(F_6\) 的纤维直方图
\[
2:8,\qquad 3:4,\qquad 4:9
\]
迫使任何与之同轨道的群作用都必须带三层稳定子分层；而在保持 fold-gauge 纤维独立置换语义时，完备异常签名仍需
\[
\Aut(\Fold_6^{\mathrm{bin}})^{\mathrm{ab}}\cong (\ZZ/2\ZZ)^{21}.
\]
因此，审计窗中的平滑包络所消去的并非单一统一对称冗余，而只能是随纤维变化的稳定子残差；其自然承载对象是带变稳定子的群胚型商，而非单一群商。
\end{theorem}
\begin{proof}
定理 \ref{thm:xi-window6-group-quotient-isotropy-strata} 已证明：任何把 \(F_6\) 纤维分解实现为群轨道分解的方案，都不可能是自由作用，也不可能退化为等势轨道；并且轨道大小只能取 \(2,3,4\)，从而稳定子恰分成三层。于是，固定自由群商与阿贝尔线性陪集解释同时被排除。

另一方面，定理 \ref{thm:xi-window6-oracle-anomaly-coordinate-bifurcation} 证明：若继续保持 fold-gauge 的纤维独立置换语义，则完备异常信息的最小承载对象仍为
\(
(\ZZ/2\ZZ)^{21}
\)；
任何把异常坐标压到 \(21\) 以下的方案都必须引入跨纤维关系，从而破坏单一纤维直积语义。故这里被“平滑”掉的并不是某个固定群的统一冗余，而是随纤维变化的稳定子残差；用群胚而非单一群来描述，正是这一变稳定子机制的自然形式化。证毕。
\end{proof}

\begin{theorem}[末位投影的渐近充分性与终端二点塌缩]\label{thm:xi-time-part67-lastbit-terminal-csiszar-collapse}
设
\[
\ell_m(w):=w_m,
\qquad
\mu_m(w):=\frac{d_m^{\mathrm{bin}}(w)}{2^m},
\qquad
u_m(w):=\frac{1}{|X_m|}
\qquad (w\in X_m).
\]
则 \(\ell_m\) 对分布对 \((\mu_m,u_m)\) 是渐近充分统计量：对任意凸函数 \(f:(0,\infty)\to\RR\) 满足 \(f(1)=0\)，有
\[
D_f(\mu_m\Vert u_m)
=
D_f\!\bigl((\ell_m)_\#\mu_m\Vert (\ell_m)_\#u_m\bigr)
+O_f\!\Bigl(\Bigl(\frac{\varphi}{2}\Bigr)^m\Bigr).
\]
并且若记
\[
p_*:=\frac{1}{\varphi\sqrt5},
\qquad
q_*:=\varphi^{-2},
\]
则
\[
\lim_{m\to\infty}D_f(\mu_m\Vert u_m)
=
D_f\!\bigl(\mathrm{Bernoulli}(p_*)\Vert \mathrm{Bernoulli}(q_*)\bigr)
=
\frac1\varphi\,f\!\left(\frac{\varphi^2}{\sqrt5}\right)
+
\frac1{\varphi^2}\,f\!\left(\frac{\varphi}{\sqrt5}\right).
\]
换言之，相对于稳定层均匀基线的全部渐近 Csisz\'ar 几何，都在末位二点实验上终端塌缩。
\end{theorem}
\begin{proof}
定理 \ref{thm:xi-fold-lastbit-asymptotic-sufficiency} 已在全变差与 Le Cam 意义下给出末位投影 \(\ell_m\) 的渐近充分性；下面只需把这一充分化提升到 \(f\)-散度层面。

由定理 \ref{thm:xi-time-part64b-fold-information-geometry-two-point-collapse}，
\[
\frac{\mu_m(w)}{u_m(w)}
=
\frac{\varphi^{2-\ell_m(w)}}{\sqrt5}
+O\!\Bigl(\Bigl(\frac{\varphi}{2}\Bigr)^m\Bigr)
\qquad (w\in X_m)
\]
在 \(w\) 上一致成立。于是存在紧区间 \(K\subset(0,\infty)\)，使得对充分大的 \(m\)，上式右端与左端都落在 \(K\) 内。由于凸函数在紧子区间上局部 Lipschitz，故对每个固定 \(f\) 有
\[
D_f(\mu_m\Vert u_m)
=
\sum_{b\in\{0,1\}}u_m(\ell_m=b)\,
f\!\left(\frac{\varphi^{2-b}}{\sqrt5}\right)
+O_f\!\Bigl(\Bigl(\frac{\varphi}{2}\Bigr)^m\Bigr).
\]

另一方面，设
\[
p_m:=((\ell_m)_\#\mu_m)(\{1\}),
\qquad
q_m:=((\ell_m)_\#u_m)(\{1\})=\frac{F_m}{F_{m+2}}.
\]
对上式在每个末位块上取 \(u_m\)-平均，得到
\[
\frac{p_m}{q_m}
=
\frac{\varphi}{\sqrt5}
+O\!\Bigl(\Bigl(\frac{\varphi}{2}\Bigr)^m\Bigr),
\qquad
\frac{1-p_m}{1-q_m}
=
\frac{\varphi^2}{\sqrt5}
+O\!\Bigl(\Bigl(\frac{\varphi}{2}\Bigr)^m\Bigr),
\]
从而
\[
D_f\!\bigl((\ell_m)_\#\mu_m\Vert (\ell_m)_\#u_m\bigr)
=
q_m\,f\!\left(\frac{p_m}{q_m}\right)
+(1-q_m)\,f\!\left(\frac{1-p_m}{1-q_m}\right)
\]
与前一式只差 \(O_f((\varphi/2)^m)\)。这便证明了渐近充分性。

最后，由
\[
q_m=\frac{F_m}{F_{m+2}}\longrightarrow \frac1{\varphi^2},
\qquad
1-q_m=\frac{F_{m+1}}{F_{m+2}}\longrightarrow \frac1{\varphi},
\]
以及上面的两个比值极限，立即得到
\[
\lim_{m\to\infty}D_f\!\bigl((\ell_m)_\#\mu_m\Vert (\ell_m)_\#u_m\bigr)
=
\frac1\varphi\,f\!\left(\frac{\varphi^2}{\sqrt5}\right)
+\frac1{\varphi^2}\,f\!\left(\frac{\varphi}{\sqrt5}\right).
\]
再与刚证得的误差估计合并，即得所示极限公式。证毕。
\end{proof}

\begin{theorem}[单一解析腔室内的波滑分裂]\label{thm:xi-time-part67-single-analytic-chamber-wave-smooth-splitting}
在真实输入 \(40\) 状态核的输出位势参数化
\[
u=e^\theta
\]
下，审计窗口
\[
|\theta|<\theta_0,
\qquad
\theta_0=0.0533477827\ldots,
\]
满足
\[
u\in (e^{-\theta_0},e^{\theta_0})
\Subset
(0,u_R),
\qquad
u_R\approx 3.59262181344.
\]
因而该窗口既不跨越正实退化点，也完全停留在
\[
\Lambda(u)<\lambda_+(u)^{1/2}
\]
的统一亚临界区间内。对于任何可归约为解析 \(\arg\max\) 读出的折叠谱观测量，若其在此窗口内出现离散波纹与连续平滑的分裂，则该分裂必发生于同一解析相内部，而不能归因于谱分支碰撞或正实相变。
\end{theorem}
\begin{proof}
推论 \ref{cor:xi-collision-u1-audit-window-uniform-subcriticality} 已证明：
\[
|\theta|<\theta_0
\Longrightarrow
u=e^\theta\in(0.932837\ldots,1.054796\ldots)\Subset(0,u_R),
\]
从而整个 \(u=1\) 邻域审计窗都处在统一亚临界腔室内，不会接触平方根门槛 \(\Lambda(u)=\lambda_+(u)^{1/2}\) 的临界端点。

另一方面，命题 \ref{prop:xi-time-part9m3-analytic-chamber-rigidity-sampling-readout} 说明：在任何不穿越退化点的实解析区间上，凡由解析 \(\arg\max\) 读出定义的纤维计数、容量与矩谱，都只能在离散交叉集之外分段常值。于是，在本窗口内若仍观察到“波纹--平滑”二分，则其来源只能是同一解析腔室内部不同读出函子的分裂，而不能被解释为谱支交换或正实相变。证毕。
\end{proof}

\begin{theorem}[两步原子剥离与 Abel 有限部位移]\label{thm:xi-time-part67-two-step-atomic-peeling-finitepart-shift}
对真实输入 \(40\) 状态核，有精确分解
\[
\Delta(z,u)=(1-uz^2)\,F(z,u),
\qquad
P(z,u)=uz^2+P_{\mathrm{core}}(z,u).
\]
在未加权情形 \(u=1\) 下，若 \(z_\star=\varphi^{-2}\)，并记 \(\mathfrak M,\mathfrak M_{\mathrm{core}}\) 分别为 full/core 的 Abel 有限部常数，则
\[
\log \mathfrak M
=
\log \mathfrak M_{\mathrm{core}}+\varphi^{-4},
\qquad
\varphi^{-4}=\frac{7-3\sqrt5}{2}.
\]
因此，由离散波纹统计转入平滑统计并非中性平均，而是一次精确的两步 primitive Euler 原子剥离。
\end{theorem}
\begin{proof}
定理 \ref{thm:xi-time-part65-atomic-primitive-abel-shift} 已直接给出
\[
P(z,u)=uz^2+P_{\mathrm{core}}(z,u),
\qquad
\log\mathfrak M=\log\mathfrak M_{\mathrm{core}}+\varphi^{-4}.
\]
其中 \(\varphi^{-4}=(7-3\sqrt5)/2\) 是同一定理中的闭式常数。于是 full/core 之差在 primitive 层只作用于长度 \(2\) 的单一原子，而在 Abel 有限部层只留下一个完全闭式的代数位移。证毕。
\end{proof}

\begin{theorem}[零温平滑包络的有限状态化]\label{thm:xi-time-part67-zero-temp-envelope-finite-stateization}
对真实输入 \(40\) 状态核的输出位势 \(P(\theta)=\log\lambda(e^\theta)\)，有
\[
P'(\theta)\longrightarrow \frac12
\qquad (\theta\to+\infty),
\]
并且满足平方根端点律
\[
\frac12-P'(\theta)=O(e^{-\theta/2}).
\]
另一方面，在零温缩放
\[
z=\sqrt{w/u}
\]
下，
\[
\lim_{u\to+\infty}\Delta\!\left(\sqrt{w/u},u\right)
=
(1-w)\det(I-wB_\infty),
\]
其中
\[
B_\infty=
\begin{pmatrix}
0&1&1&0\\
0&0&1&0\\
0&0&3&1\\
1&0&0&3
\end{pmatrix}.
\]
因此，零温极限中被选出的平滑包络并非整个 \(40\) 状态核的粗粒化投影，而是由显式四状态地基 SFT 所控制的有限状态对象；原子层仅留下因子 \((1-w)\)，而主导极点由地基层而非原子层决定。
\end{theorem}
\begin{proof}
定理 \ref{thm:xi-output-potential-zero-temp-square-root-critical-law} 给出
\[
\frac12-P'(\theta)=\frac{d}{2c}e^{-\theta/2}+O(e^{-\theta}),
\]
故零温端点的饱和斜率确为 \(1/2\)。

再由定理 \ref{thm:xi-zero-temp-atomic-ground-exact-splitting}，
\[
\lim_{u\to+\infty}\Delta\!\left(\sqrt{w/u},u\right)
=(1-w)\det(I-wB_\infty),
\]
这说明零温解析对象严格分裂为一个 primitive 原子因子与一个四状态地基 SFT 的动力学行列式。最后，推论 \ref{cor:xi-zero-temp-leading-pole-ground-dominance} 进一步表明正实最内层零点来自 \(\det(I-wB_\infty)\) 而非 \((1-w)\)，故主导平滑尺度由四状态地基层决定。证毕。
\end{proof}

\begin{theorem}[连续相位层与外置算术核]\label{thm:xi-time-part67-continuous-phase-external-arithmetic-kernel}
设 \(\varnothing\neq S\subset\mathbb P\) 为有限素数集，并记
\[
G_S:=\ZZ[S^{-1}].
\]
则其 Pontryagin 对偶满足短正合列
\[
0\longrightarrow \prod_{p\in S}\ZZ_p
\longrightarrow \widehat{G_S}
\longrightarrow \TT
\longrightarrow 0,
\]
并且该 profinite 核 \(\prod_{p\in S}\ZZ_p\) 可由相位采样函数 \(\Sigma^{\mathrm{ph}}_{G_S}\) 唯一恢复。因而圆相位因子 \(\TT\) 只是连通可见层，真正的算术内容则以外置 prime-register 核的形式保留在 profinite 方向上。

更进一步，不存在任何有限秩加法账本能够把正整数乘法幺半群忠实严格同态化；因此，上述 profinite 核不能被有限多个加法寄存器吸收。
\end{theorem}
\begin{proof}
定理 \ref{thm:xi-localized-phase-sampling-closed-form} 已证明：\(\Sigma^{\mathrm{ph}}_{G_S}\) 唯一恢复素数集 \(S\)，并因而唯一恢复对偶短正合列中的 profinite 核
\[
\prod_{p\in S}\ZZ_p.
\]
定理 \ref{thm:xi-localized-basic-extension-strictly-nonsplit} 进一步说明该短正合列在 \(S\neq\varnothing\) 时严格不分裂，因此 \(\TT\) 不能替代整个对偶对象；它只给出连通商，而非全部算术信息。

最后，定理 \ref{thm:xi-time-part64b-local-additivity-nonlocal-multiplicativity-polarization} 已排除一切把
\(
(\NN_{\times},\times)
\)
严格忠实同态化到有限秩加法账本的方案。故这里的 prime-register profinite 核确实是不可由有限加法寄存器消化的外置算术部分。证毕。
\end{proof}

\noindent
上述七条结果表明，离散间隔统计中的波纹并不应被理解为连续对象自身的粗糙化。在本文框架内，它们恰是非平凡角色扇区、window-\(6\) 的变稳定子残差、两步 primitive 原子与 prime-register profinite 核在离散读出层的共同显影；相应地，平滑极限则由平凡角色投影、末位二点塌缩、Abel 有限部反项剥离以及零温地基层选择共同构成。因而，“波”与“滑”并非同一对象的两种视觉描画，而是同一算术动力学对象在不同函子层上的两种范畴像。

\input{subsubsec__xi-time-protocol-conclusions-part68-spectral-ldp-finitepart-entropy-rigidity}
\input{subsubsec__xi-time-protocol-conclusions-part68a-primitive-parity-window-lambert-refinement}
\input{subsubsec__xi-time-protocol-conclusions-part68aa-positive-chamber-algebraic-anchor-rigidity}

\endinput


\endinput


\endinput

\paragraph{dyadic 屏幕 exactization、边界码退化、矩层析二歧与读出精度分离}
在 dyadic 体--界全息链的精确屏幕、边界像 LDPC 阴影、有限矩盒完备性与对角边界译码阈值都已固定之后，还可继续抽取出下列四条派生结果。它们分别把部分屏幕的 exactization 代价压缩为加权拟阵极小化，把 dyadic 边界像码的相对距离退化写成闭式，把固定阶边界矩与分辨率自适应矩盒之间的分界提升为严格二歧，并把单标量 Joukowsky--Gödel 读出与多通道边界译码之间的精度差距量化为指数分离。

\begin{theorem}[部分屏幕 exactization 的加权拟阵极小化与拓扑缺口恒等式]\label{thm:xi-time-part65a-weighted-exactization-matroid-top-gap}
\leanverified{paper\_xi\_time\_part65a\_weighted\_exactization\_matroid\_top\_gap}
沿用定理 \ref{thm:xi-time-part9fa-exact-screen-minimal-size-matroid-basis} 与定理 \ref{thm:spg-screen-rank-matroid-supermodularity} 的记号，并令
\[
M_{m,n}:=\bigl(\overrightarrow{\mathcal F_m^{n-1}},r\bigr)
\]
为由可见秩函数 \(r(\cdot)\) 定义的边界行拟阵，记其总秩为
\[
\rho_{m,n}:=2^{mn}.
\]
给定任意非负权
\[
c:\overrightarrow{\mathcal F_m^{n-1}}\to \RR_{\ge 0},
\]
并对任意部分屏幕 \(S\subseteq \overrightarrow{\mathcal F_m^{n-1}}\) 定义 exactization 的最小追加代价
\[
\mathcal C_c(S):=
\min\Bigl\{
\sum_{\vec f\in T\setminus S}c(\vec f):
\ T\supseteq S,\ \ker f_T=0
\Bigr\}.
\]
则有
\[
\mathcal C_c(S)
=
\min\Bigl\{
\sum_{\vec f\in B}c(\vec f):
\ B\subseteq \overrightarrow{\mathcal F_m^{n-1}}\setminus S,\
B\text{ 是 }M_{m,n}/S\text{ 的一组基}
\Bigr\}.
\]
特别地，当 \(c\equiv 1\) 时，
\[
\mathcal C_1(S)
=
\rho_{m,n}-r(S)
=
2^{mn}-r(S)
=
\delta(S)
=
\rank_{\ZZ}\widetilde H_{n-1}(A_{\neg S};\ZZ).
\]
\end{theorem}
\begin{proof}
由定理 \ref{thm:xi-time-part9fa-exact-screen-minimal-size-matroid-basis}，精确屏幕恰由边界行拟阵的基控制，且其总秩为
\(
\rho_{m,n}=2^{mn}
\)。
对任意 \(B\subseteq \overrightarrow{\mathcal F_m^{n-1}}\setminus S\)，收缩拟阵 \(M_{m,n}/S\) 的秩函数满足
\[
r_{M_{m,n}/S}(B)=r(S\cup B)-r(S).
\]
因此，\(B\) 为 \(M_{m,n}/S\) 的基当且仅当
\[
r(S\cup B)=\rho_{m,n},
\qquad
|B|=\rho_{m,n}-r(S),
\]
也即当且仅当 \(S\cup B\) 是一个达到满秩而不含冗余追加面的 exactization。

现取任意 \(T\supseteq S\) 满足 \(\ker f_T=0\)。由定理
\ref{thm:xi-time-part9fa-exact-screen-minimal-size-matroid-basis} 的首部分，\(\ker f_T=0\) 等价于 \(r(T)=\rho_{m,n}\)，故 \(T\setminus S\) 在收缩拟阵 \(M_{m,n}/S\) 中张成满秩子集。于是可取某组基
\[
B\subseteq T\setminus S.
\]
由于权函数 \(c\) 非负，便有
\[
\sum_{\vec f\in B}c(\vec f)\le \sum_{\vec f\in T\setminus S}c(\vec f).
\]
这说明任意 exactization 的追加代价都不小于某组收缩拟阵基的权和；反过来，任一基 \(B\) 都给出 \(T=S\cup B\) 的 exactization。故第一条公式成立。

当 \(c\equiv 1\) 时，收缩拟阵的每一组基都具有基数
\(
\rho_{m,n}-r(S)
\)，于是
\[
\mathcal C_1(S)=\rho_{m,n}-r(S)=2^{mn}-r(S).
\]
再由定理 \ref{thm:spg-screen-rank-matroid-supermodularity}，
\[
\delta(S)=\rho_{m,n}-r(S),
\]
而定理 \ref{thm:xi-time-part9fa-screen-kernel-top-homology} 给出
\[
\delta(S)=\rank_{\ZZ}\widetilde H_{n-1}(A_{\neg S};\ZZ).
\]
合并即得结论。证毕。
\end{proof}

\begin{corollary}[dyadic 边界像码的正码率--零相对距离退化与固定唯一纠错半径]\label{cor:xi-time-part65a-dyadic-boundary-code-not-asymptotically-good}
固定 \(n\ge 1\)，并沿用定理 \ref{thm:xi-time-part9fa-dyadic-boundary-image-ldpc} 的记号。则 dyadic 边界像码族
\[
\bigl(\mathcal C_{m,n}\bigr)_{m\ge 1}
\]
的相对距离满足
\[
\Delta_{m,n}
:=
\frac{d_{\min}(\mathcal C_{m,n})}{N_{m,n}}
=
\frac{2}{(2^m+1)\,2^{m(n-1)}}
\xrightarrow[m\to\infty]{}
0.
\]
另一方面，其码率满足
\[
\operatorname{Rate}_{m,n}
:=
\frac{\dim_{\FF_2}\mathcal C_{m,n}}{N_{m,n}}
=
\frac{2^m}{n(2^m+1)}
\xrightarrow[m\to\infty]{}
\frac1n.
\]
因此，对任意固定 \(n\)，码族 \((\mathcal C_{m,n})_{m\ge 1}\) 不是渐近优码族。并且任何最优最大似然唯一译码的可保证纠错半径都满足
\[
t_{m,n}^{\operatorname{ML}}
\le
\left\lfloor \frac{d_{\min}(\mathcal C_{m,n})-1}{2}\right\rfloor
=
n-1.
\]
特别地，在最坏情形的对抗错误模型下，可保证唯一纠错的边界面数与分辨率 \(m\) 无关。
\end{corollary}
\begin{proof}
由定理 \ref{thm:xi-time-part9fa-dyadic-boundary-image-ldpc}，
\[
N_{m,n}=n(2^m+1)2^{m(n-1)},
\qquad
\dim_{\FF_2}\mathcal C_{m,n}=2^{mn},
\qquad
d_{\min}(\mathcal C_{m,n})=2n.
\]
直接代入即可得到
\[
\Delta_{m,n}
=
\frac{2n}{n(2^m+1)2^{m(n-1)}}
=
\frac{2}{(2^m+1)2^{m(n-1)}}
\to 0
\]
以及
\[
\operatorname{Rate}_{m,n}
=
\frac{2^{mn}}{n(2^m+1)2^{m(n-1)}}
=
\frac{2^m}{n(2^m+1)}
\to \frac1n.
\]
故该码族具有正码率极限而相对距离趋于零，从而不是渐近优码族。

对于任意二元线性码，统一唯一纠错半径至多为
\[
\left\lfloor \frac{d_{\min}-1}{2}\right\rfloor.
\]
将 \(d_{\min}(\mathcal C_{m,n})=2n\) 代入便得
\[
t_{m,n}^{\operatorname{ML}}\le n-1.
\]
证毕。
\end{proof}

\begin{theorem}[固定阶边界矩不可分与分辨率自适应矩盒完备的严格二歧]\label{thm:xi-time-part65a-fixed-order-boundary-moment-vs-adaptive-box-dichotomy}
固定 \(n\ge 2\)。对每个分辨率 \(m\ge 1\)，设
\[
\mathcal A_m\subseteq \NN^n
\]
为一族多重指标。若存在常数 \(R<\infty\) 使得
\[
\mathcal A_m\subseteq \{\alpha\in\NN^n:\ |\alpha|\le R\}
\qquad(\forall m\ge 1),
\]
则边界通量矩字典
\[
U\longmapsto \bigl(M_\alpha(U)\bigr)_{\alpha\in\mathcal A_m}
\]
不可能在全部 \(m\)-级 dyadic polycube 上同时单射。

相反，若取分辨率自适应矩盒
\[
\mathcal B_m:=\{0,\dots,2^m-1\}^n,
\]
则体积分矩字典
\[
U\longmapsto \bigl(\mu_\alpha(U)\bigr)_{\alpha\in\mathcal B_m}
\]
在全部 \(m\)-级 dyadic polycube 上为单射；并且边界 Gödel 整数在固定分辨率类内唯一决定 \(U\)。
\end{theorem}
\begin{proof}
若存在统一总次数上界 \(R\)，则对 \(m=R+1\) 应用定理
\ref{thm:spg-prouhet-thue-morse-obstruction-dyadic-polyclube-flux-moments}，可得两个互异的 \(m\)-级 dyadic polycube
\[
U\neq V
\]
满足
\[
M_\alpha(U)=M_\alpha(V)
\qquad
(\forall \alpha\in\NN^n,\ |\alpha|\le R).
\]
由于 \(\mathcal A_m\subseteq \{|\alpha|\le R\}\)，遂有
\[
\bigl(M_\alpha(U)\bigr)_{\alpha\in\mathcal A_m}
=
\bigl(M_\alpha(V)\bigr)_{\alpha\in\mathcal A_m},
\]
故该字典在 \(m=R+1\) 时已非单射，从而不可能在全部分辨率上同时单射。

另一方面，定理 \ref{thm:xi-time-part9fa-dyadic-finite-moment-completeness} 已证明：对每个固定 \(m\)，矩盒
\[
\mathcal B_m=\{0,\dots,2^m-1\}^n
\]
给出的体积分矩字典是单射。再由推论
\ref{cor:spg-boundary-godel-finite-moment-completeness}，边界 Gödel 整数在固定分辨率类内唯一决定 dyadic polycube。本定理遂得证。证毕。
\end{proof}

\begin{theorem}[Joukowsky--Gödel 单标量读出与多通道边界译码的指数精度分离]\label{thm:xi-time-part65a-scalar-vs-multichannel-exponential-precision-gap}
\leanverified{paper\_xi\_time\_part65a\_scalar\_vs\_multichannel\_exponential\_precision\_gap}
取互异素数 \(p_1,\dots,p_k\)，并沿用定理 \ref{thm:spg-single-parameter-log-readout-vs-diagonal-multichannel-bifurcation} 的单标量读出
\[
u(b):=\sum_{i=1}^k b_i\log p_i,
\qquad
b\in\{0,1\}^k.
\]
若要求仅由 \(u(b)\) 的单标量观测，在统一绝对误差阈值 \(\eta_k\) 下对全部 \(2^k\) 个状态保证无碰撞唯一译码，则必有
\[
\eta_k
\le
\frac{1}{2(2^k-1)}
\sum_{i=1}^k\log p_i.
\]

另一方面，多通道对角边界译码允许的统一相对误差阈值为
\[
\varepsilon_\ast
:=
\frac{p_{\min}-1}{p_{\min}+1},
\qquad
p_{\min}:=\min_{1\le i\le k}p_i,
\]
并且该阈值由定理 \ref{thm:spg-single-parameter-log-readout-vs-diagonal-multichannel-bifurcation} 给出充分性、由命题 \ref{prop:spg-relative-error-threshold-sharpness} 给出最坏情形下的尖锐必要性。

因此，任意试图用单一标量 \(u(b)\) 取代多通道边界译码的方案，都必然支付指数级绝对精度代价。若进一步取 \(p_i\) 为前 \(k\) 个素数，则
\[
\eta_k=O\!\left(\frac{k\log k}{2^k}\right),
\qquad
\varepsilon_\ast=\frac13.
\]
\end{theorem}
\begin{proof}
由定理 \ref{thm:spg-single-parameter-log-readout-vs-diagonal-multichannel-bifurcation}，存在 \(b\neq b'\) 使
\[
|u(b)-u(b')|
\le
\frac{1}{2^k-1}\sum_{i=1}^k\log p_i.
\]
若要求在绝对误差不超过 \(\eta_k\) 的条件下对全部状态无碰撞唯一译码，则以各点 \(u(b)\) 为中心、半径 \(\eta_k\) 的闭区间必须两两不交。于是必要地有
\[
2\eta_k
\le
\min_{b\neq b'}|u(b)-u(b')|
\le
\frac{1}{2^k-1}\sum_{i=1}^k\log p_i,
\]
从而得到
\[
\eta_k
\le
\frac{1}{2(2^k-1)}
\sum_{i=1}^k\log p_i.
\]

多通道情形的充分性由定理
\ref{thm:spg-single-parameter-log-readout-vs-diagonal-multichannel-bifurcation} 直接给出，而阈值
\(
\varepsilon_\ast=\frac{p_{\min}-1}{p_{\min}+1}
\)
在最坏情形下的必要性由命题
\ref{prop:spg-relative-error-threshold-sharpness} 给出，故该相对误差门槛是 sharp 的。

若 \(p_i\) 取前 \(k\) 个素数，则 \(p_{\min}=2\)，故
\[
\varepsilon_\ast=\frac{2-1}{2+1}=\frac13.
\]
同时由素数定理的标准推论，
\[
\sum_{i=1}^k\log p_i=\vartheta(p_k)\sim p_k\sim k\log k,
\]
代回上式即得
\[
\eta_k=O\!\left(\frac{k\log k}{2^k}\right).
\]
证毕。
\end{proof}

\noindent
这四条派生结果表明，dyadic 边界协议中的可恢复性门槛可分解为三类严格量化的结构参数：部分屏幕的秩亏由收缩拟阵基选择与顶维同调秩共同控制，边界像码的渐近限制由相对距离退化而非码率塌缩主导，矩层析的完备性必须随分辨率同步扩张，而边界读出的鲁棒性则依赖于多通道坐标而不能被无损压缩为单一实参数。

\endinput

\paragraph{Gödel--Tate 地址像的严格自相似闭合与链内算子格的 Boolean 算术随机化}
在单向量 \(2^L\)-进全息像集的同胚、自相似、维数与 Haar 二歧律已经确定之后，Gödel--Tate 地址链还可进一步压缩为一条显式的 Cantor--Bernoulli 闭合；另一方面，链内算子格的平方自由 Gödel 算术化也可进一步收束为一组完全显式的 Boolean--概率闭式。下列四条结果分别统一这两条主线的拓扑、测度、格论与随机结构。

\begin{theorem}[Gödel--Tate 地址像的严格自相似、移位共轭与 clopen cylinder 分解]\label{thm:xi-time-part65d-godeltate-selfsimilar-shift-cylinder}
设 \(E\) 为有限原子事件集，\(\mathrm{code}:E\hookrightarrow \NN\) 为单射，并记
\[
M:=\max_{e\in E}\mathrm{code}(e),
\qquad
B:=2^L>M,
\qquad
T:=T_2(E_{\min}).
\]
取任意非零向量 \(v\in T\)，并记
\[
D:=\mathrm{code}(E)\subseteq \{0,\dots,B-1\},
\qquad
T_v:=\ZZ_2v\subseteq T.
\]
对任意无限事件流 \(\mathbf e=(e_t)_{t\ge 1}\in E^{\NN}\)，定义
\[
\Xi(\mathbf e):=X(\mathbf e):=\sum_{t\ge 1}\mathrm{code}(e_t)\,B^{t-1}v,
\qquad
\mathcal X_E:=\Xi(E^{\NN})\subseteq T_v.
\]
则成立：
\begin{enumerate}
  \item \(\mathcal X_E\) 是紧、全不连通集，并满足严格不交分解
  \[
  \mathcal X_E=\bigsqcup_{a\in D}\bigl(av+B\mathcal X_E\bigr).
  \]
  \item 记左移算子
  \[
  \sigma:E^{\NN}\to E^{\NN},
  \qquad
  \sigma(e_1,e_2,\dots)=(e_2,e_3,\dots).
  \]
  则 \(\Xi:E^{\NN}\to \mathcal X_E\) 为拓扑同胚；并且若定义
  \[
  S_B(av+Bx):=x
  \qquad
  (a\in D,\ x\in \mathcal X_E),
  \]
  则 \(S_B\) 良定义且满足共轭关系
  \[
  S_B\circ \Xi=\Xi\circ \sigma.
  \]
  等价地，对任意 \(\mathbf e\in E^{\NN}\)，有
  \[
  \Xi(\mathbf e)=\mathrm{code}(e_1)v+B\,\Xi(\sigma\mathbf e).
  \]
  \item 记 \(N:=|E|=|D|\)。对每个 \(n\ge 1\) 与每个 digit 前缀
  \[
  (a_1,\dots,a_n)\in D^n,
  \]
  定义
  \[
  \mathcal X_E[a_1,\dots,a_n]
  :=
  \left(\sum_{t=1}^{n}a_tB^{t-1}v\right)+B^n\mathcal X_E.
  \]
  则这 \(N^n\) 个集合两两不交，且都是 \(\mathcal X_E\) 中的 clopen cylinder，并满足
  \[
  \mathcal X_E=\bigsqcup_{(a_1,\dots,a_n)\in D^n}\mathcal X_E[a_1,\dots,a_n].
  \]
\end{enumerate}
\end{theorem}
\begin{proof}
由定理 \ref{thm:xi-time-part62d-hologram-image-cantor-homeomorphism}，地址映射
\[
\Xi:E^{\NN}\longrightarrow \mathcal X_E
\]
已知是拓扑同胚，并且满足严格不交分解
\[
\mathcal X_E=\bigsqcup_{a\in D}(av+B\mathcal X_E).
\]
因此第 \(1\) 条与第 \(2\) 条中的同胚性立即成立；而对任意
\(
\mathbf e=(e_1,e_2,\dots)
\)，按定义拆开首位即可得到
\[
\Xi(\mathbf e)=\mathrm{code}(e_1)v+B\,\Xi(\sigma\mathbf e),
\]
故 \(S_B\circ \Xi=\Xi\circ \sigma\)，并且 \(S_B\) 由上述严格不交分解决定而良定义。

对第 \(3\) 条，由同一定理，对任意事件前缀 \(\mathbf a=(e_1,\dots,e_n)\in E^n\) 及其 cylinder
\[
[\mathbf a]:=\{\mathbf f\in E^{\NN}:\ \mathrm{pref}_n(\mathbf f)=\mathbf a\},
\]
都有
\[
\Xi([\mathbf a])
=
\left(\sum_{t=1}^{n}\mathrm{code}(e_t)\,B^{t-1}v\right)+B^n\mathcal X_E.
\]
由于 \(\mathrm{code}\) 单射，\(D^n\) 中每个 digit 前缀与 \(E^n\) 中唯一事件前缀一一对应，故上式正是所定义的 \(\mathcal X_E[a_1,\dots,a_n]\)。而 \(\{[\mathbf a]:\mathbf a\in E^n\}\) 在 \(E^{\NN}\) 中本来就是由 \(N^n\) 个两两不交 clopen 集组成的分割，经同胚 \(\Xi\) 送到 \(\mathcal X_E\) 后，便得到所述 \(N^n\) 个两两不交 clopen cylinder 的不交并分解。证毕。
\end{proof}

\begin{theorem}[Gödel--Tate 地址像的精确 \texorpdfstring{$2$}{2}-进维数与 Haar 测度二分律]\label{thm:xi-time-part65d-godeltate-dimension-haar-dichotomy}
沿用定理 \ref{thm:xi-time-part65d-godeltate-selfsimilar-shift-cylinder} 的记号，并记
\[
N:=|E|=|D|.
\]
则 \(\mathcal X_E\subseteq T_v\cong \ZZ_2\) 的 Hausdorff 维数与上、下 Minkowski 维数存在且相等，并满足
\[
\dim_{\mathrm H}(\mathcal X_E)
=
\dim_{\mathrm M}(\mathcal X_E)
=
\frac{\log N}{\log B}
=
\frac{\log N}{L\log 2}.
\]
此外，\(\mathcal X_E\) 关于地址线 \(T_v\) 的归一化 Haar 测度满足二分律：
\begin{enumerate}
  \item 若 \(N<B\)，则
  \[
  \mu_{T_v}(\mathcal X_E)=0,
  \]
  并且 \(\mathcal X_E\) 在 \(T_v\) 中内部为空；特别地，当 \(2\le N<B\) 时，\(\mathcal X_E\) 为 nowhere dense 的 Cantor 型集；
  \item 若 \(N=B\)，则等价地 \(D=\{0,1,\dots,B-1\}\)，并且
  \[
  \mathcal X_E=T_v,
  \qquad
  \mu_{T_v}(\mathcal X_E)=1.
  \]
\end{enumerate}
\end{theorem}
\begin{proof}
由结论 \ref{con:xi-time-part62b-rankone-hologram-selfsimilar-measure-dimension}，对任意有限 digit 集 \(D\subseteq\{0,\dots,B-1\}\)，其单向量 \(2^L\)-进全息像集在地址线口径下都满足
\[
\dim_{\mathrm H}=\dim_{\mathrm M}=\frac{\log|D|}{\log B},
\]
并且 Haar 测度恰为
\[
\mu_{T_v}=
\begin{cases}
1,& |D|=B,\\
0,& |D|<B.
\end{cases}
\]
将 \(|D|=N\) 代入，立即得到维数公式及测度二分律。

若 \(N<B\)，同一结论还给出 \(\mathcal X_E\) 在 \(T_v\) 中内部为空。另一方面，由定理 \ref{thm:xi-time-part65d-godeltate-selfsimilar-shift-cylinder}，\(\mathcal X_E\) 为紧集，因而在紧 Hausdorff 群 \(T_v\) 中闭。闭且内部为空即意味着 nowhere dense。

若 \(N=B\)，则因 \(D\subseteq\{0,\dots,B-1\}\) 且 \(|D|=B\)，只能有
\[
D=\{0,1,\dots,B-1\}.
\]
此时同一结论给出 \(\mathcal X_E=T_v\)。证毕。
\end{proof}

\begin{theorem}[链内算子格的 Boolean 算术化、Möbius 函数与特征多项式]\label{thm:xi-time-part65d-chain-interior-boolean-arithmeticization}
对链
\[
C_n=\{0<1<\cdots<n-1\},
\]
记
\[
\mathcal I_n:=\operatorname{Int}(C_n)
\]
为其全部内算子所成偏序格。沿用定理 \ref{thm:killo-chain-interior-godel-gcd-lcm} 的记号，对每个 \(I\in \mathcal I_n\) 记
\[
\kappa(I):=\prod_{i\in \mathrm{Fix}(I)}q_i,
\qquad
Q_n:=\prod_{i=0}^{n-1}q_i.
\]
则成立：
\begin{enumerate}
  \item \((\mathcal I_n,\preceq)\) 是秩为 \(n-1\) 的 Boolean 格；其秩函数可取为
  \[
  \operatorname{rk}(I):=|\mathrm{Fix}(I)|-1.
  \]
  \item 映射 \(\kappa\) 给出格同构
  \[
  \mathcal I_n
  \cong
  \{d\mid Q_n:\ q_0\mid d\},
  \]
  其中右侧按整除序排序；并且对任意 \(I,J\in \mathcal I_n\)，有
  \[
  \kappa(I\wedge J)=\gcd(\kappa(I),\kappa(J)),
  \qquad
  \kappa(I\vee J)=\mathrm{lcm}(\kappa(I),\kappa(J)).
  \]
  \item 对任意 \(I\preceq J\)，其 Möbius 函数满足
  \[
  \mu_{\mathcal I_n}(I,J)
  =
  (-1)^{\operatorname{rk}(J)-\operatorname{rk}(I)}.
  \]
  \item 其特征多项式满足
  \[
  \chi_{\mathcal I_n}(t)=(t-1)^{n-1}.
  \]
\end{enumerate}
\end{theorem}
\begin{proof}
由定理 \ref{thm:killo-chain-interior-godel-gcd-lcm}，映射
\[
I\longmapsto \mathrm{Fix}(I)
\]
把 \(\mathcal I_n\) 规范识别为全部满足 \(0\in F\subseteq\{0,\dots,n-1\}\) 的子集格，并把 meet/join 分别送到交与并。去掉强制出现的元素 \(0\) 后，便得到
\[
\mathcal I_n\cong \mathcal P(\{1,\dots,n-1\}),
\]
故 \(\mathcal I_n\) 是秩为 \(n-1\) 的 Boolean 格，且秩函数正是
\(
|\mathrm{Fix}(I)|-1
\)。

同一条定理还给出
\[
\kappa(I)=\prod_{i\in \mathrm{Fix}(I)}q_i,
\]
因此 \(\kappa(\mathcal I_n)\) 恰为 \(Q_n\) 的全部含 \(q_0\) 因子的平方自由约数；由于 \(Q_n\) 本身平方自由，这正是集合
\[
\{d\mid Q_n:\ q_0\mid d\}.
\]
同时 \(\kappa\) 将 meet/join 精确送到 \(\gcd/\mathrm{lcm}\)，故第 \(2\) 条成立。

对 Möbius 函数与特征多项式，定理
\ref{thm:xi-chain-interior-boolean-mobius-characteristic-maxchains} 与推论
\ref{cor:xi-chain-interior-dirichlet-mobius-characteristic-closure} 已给出
\[
\mu_{\mathcal I_n}(I,J)=(-1)^{|\mathrm{Fix}(J)\setminus \mathrm{Fix}(I)|}
\]
以及
\[
\chi_{\mathcal I_n}(t)=(t-1)^{n-1}.
\]
而由秩函数定义，
\[
|\mathrm{Fix}(J)\setminus \mathrm{Fix}(I)|
=
\bigl(|\mathrm{Fix}(J)|-1\bigr)-\bigl(|\mathrm{Fix}(I)|-1\bigr)
=
\operatorname{rk}(J)-\operatorname{rk}(I),
\]
遂得第 \(3\) 条。证毕。
\end{proof}

\leanverified{paper\_xi\_time\_part65d\_chain\_interior\_interval\_squarefree\_divisor\_isomorphism}
\begin{theorem}[链内算子区间与平方自由整除区间的精确算术同构]\label{thm:xi-time-part65d-chain-interior-interval-squarefree-divisor-isomorphism}
沿用定理 \ref{thm:xi-time-part65d-chain-interior-boolean-arithmeticization} 的记号。对任意
\[
I\preceq J\in \mathcal I_n,
\]
定义
\[
m_{I,J}:=\frac{\kappa(J)}{\kappa(I)}.
\]
则 \(m_{I,J}\) 是平方自由整数，并且区间
\[
[I,J]:=\{K\in \mathcal I_n:\ I\preceq K\preceq J\}
\]
与整除格
\[
\{d\in\NN:\ d\mid m_{I,J}\}
\]
之间存在规范格同构
\[
\Phi_{I,J}:[I,J]\longrightarrow \{d\in\NN:\ d\mid m_{I,J}\},
\qquad
\Phi_{I,J}(K):=\frac{\kappa(K)}{\kappa(I)}.
\]
并且对任意 \(K_1,K_2\in [I,J]\)，有
\[
\Phi_{I,J}(K_1\wedge K_2)
=
\gcd\!\bigl(\Phi_{I,J}(K_1),\Phi_{I,J}(K_2)\bigr),
\]
\[
\Phi_{I,J}(K_1\vee K_2)
=
\mathrm{lcm}\!\bigl(\Phi_{I,J}(K_1),\Phi_{I,J}(K_2)\bigr).
\]
\end{theorem}
\begin{proof}
由定理 \ref{thm:killo-chain-interior-godel-gcd-lcm}，
\[
I\longmapsto \mathrm{Fix}(I)
\]
把 \(\mathcal I_n\) 识别为满足 \(0\in F\subseteq\{0,\dots,n-1\}\) 的子集格，并把偏序送到包含关系。因而
\[
I\preceq K\preceq J
\quad\Longleftrightarrow\quad
\mathrm{Fix}(I)\subseteq \mathrm{Fix}(K)\subseteq \mathrm{Fix}(J).
\]
于是
\[
m_{I,J}
=
\frac{\kappa(J)}{\kappa(I)}
=
\prod_{i\in \mathrm{Fix}(J)\setminus \mathrm{Fix}(I)}q_i
\]
确为平方自由整数。

对任意 \(K\in[I,J]\)，同样有
\[
\Phi_{I,J}(K)
=
\frac{\kappa(K)}{\kappa(I)}
=
\prod_{i\in \mathrm{Fix}(K)\setminus \mathrm{Fix}(I)}q_i.
\]
当 \(K\) 在区间 \([I,J]\) 内变化时，\(\mathrm{Fix}(K)\setminus \mathrm{Fix}(I)\) 恰遍历
\(\mathrm{Fix}(J)\setminus \mathrm{Fix}(I)\) 的全部子集，故 \(\Phi_{I,J}(K)\) 恰遍历 \(m_{I,J}\) 的全部正因子。于是 \(\Phi_{I,J}\) 为双射。

再由定理 \ref{thm:killo-chain-interior-godel-gcd-lcm}，
\[
\kappa(K_1\wedge K_2)=\gcd(\kappa(K_1),\kappa(K_2)),
\qquad
\kappa(K_1\vee K_2)=\mathrm{lcm}(\kappa(K_1),\kappa(K_2)).
\]
而 \(\Phi_{I,J}(K_\nu)\) 的素因子都来自 \(\mathrm{Fix}(J)\setminus \mathrm{Fix}(I)\)，故与 \(\kappa(I)\) 互素，从而
\[
\gcd(\kappa(K_1),\kappa(K_2))
=
\kappa(I)\,
\gcd\!\bigl(\Phi_{I,J}(K_1),\Phi_{I,J}(K_2)\bigr),
\]
\[
\mathrm{lcm}(\kappa(K_1),\kappa(K_2))
=
\kappa(I)\,
\mathrm{lcm}\!\bigl(\Phi_{I,J}(K_1),\Phi_{I,J}(K_2)\bigr).
\]
两边同除以 \(\kappa(I)\) 即得所述运算相容性。证毕。
\end{proof}

\begin{corollary}[链内算子区间的 Möbius 函数与特征多项式的整数退化]\label{cor:xi-time-part65d-chain-interval-mobius-characteristic-arithmetic}
沿用定理 \ref{thm:xi-time-part65d-chain-interior-interval-squarefree-divisor-isomorphism} 的记号。则对任意 \(I\preceq J\)，有
\[
\mu_{\mathcal I_n}(I,J)=\mu(m_{I,J}),
\]
其中右端 \(\mu\) 表示经典整数 Möbius 函数。若 \(\omega(m)\) 表示整数 \(m\) 的不同素因子个数，则区间 \([I,J]\) 的特征多项式满足
\[
\chi_{[I,J]}(t)=(t-1)^{\omega(m_{I,J})}.
\]
\end{corollary}
\begin{proof}
把
\[
m_{I,J}=\prod_{r=1}^{s}p_r
\]
写成互异素数乘积，则 \(\{d\in\NN:\ d\mid m_{I,J}\}\) 恰是秩为 \(s=\omega(m_{I,J})\) 的 Boolean 格。由定理
\ref{thm:xi-time-part65d-chain-interior-interval-squarefree-divisor-isomorphism}，区间 \([I,J]\) 与该 Boolean 格同构，故其 Möbius 函数为
\[
(-1)^s=\mu(m_{I,J}),
\]
而特征多项式即
\[
(t-1)^s=(t-1)^{\omega(m_{I,J})}.
\]
证毕。
\end{proof}

\begin{theorem}[链内算子算术化的 Dirichlet 生成函数、矩母函数与 Hoeffding 型偏差界]\label{thm:xi-time-part65d-chain-interior-dirichlet-mgf-hoeffding}
沿用定理 \ref{thm:xi-time-part65d-chain-interior-boolean-arithmeticization} 的记号，并在 \(\mathcal I_n\) 上取均匀分布。定义随机变量
\[
K_n(I):=\kappa(I)=\prod_{i\in \mathrm{Fix}(I)}q_i.
\]
则成立：
\begin{enumerate}
  \item 对任意 \(s\in \CC\)，其 Dirichlet 生成函数满足
  \[
  \sum_{I\in \mathcal I_n}K_n(I)^{-s}
  =
  q_0^{-s}\prod_{i=1}^{n-1}(1+q_i^{-s}).
  \]
  \item 对任意 \(t\in \RR\)，\(\log K_n\) 的矩母函数满足
  \[
  \EE\!\left[e^{t\log K_n}\right]
  =
  q_0^t\prod_{i=1}^{n-1}\frac{1+q_i^t}{2}.
  \]
  \item 其均值与方差为
  \[
  \EE[\log K_n]
  =
  \log q_0+\frac12\sum_{i=1}^{n-1}\log q_i,
  \qquad
  \mathrm{Var}(\log K_n)
  =
  \frac14\sum_{i=1}^{n-1}(\log q_i)^2.
  \]
  \item 对任意 \(u>0\)，有 Hoeffding 型偏差界
  \[
  \PP\!\left(\left|\log K_n-\EE[\log K_n]\right|\ge u\right)
  \le
  2\exp\!\left(
  -\frac{2u^2}{\sum_{i=1}^{n-1}(\log q_i)^2}
  \right).
  \]
\end{enumerate}
\end{theorem}
\begin{proof}
由定理 \ref{thm:xi-chain-interior-random-godel-euler-logvariance}，存在独立同分布随机变量
\[
\xi_1,\dots,\xi_{n-1}\stackrel{\mathrm{i.i.d.}}{\sim}\mathrm{Bernoulli}\!\left(\frac12\right)
\]
使得
\[
K_n=q_0\prod_{i=1}^{n-1}q_i^{\xi_i},
\qquad
\log K_n=\log q_0+\sum_{i=1}^{n-1}\xi_i\log q_i.
\]

第 \(1\) 条也可由定理 \ref{thm:xi-chain-interior-godel-euler-mobius-factorization} 直接得到；等价地，按 \(\xi_i\in\{0,1\}\) 求和可写为
\[
\sum_{I\in \mathcal I_n}K_n(I)^{-s}
=
q_0^{-s}\sum_{\xi_i\in\{0,1\}}
\prod_{i=1}^{n-1}q_i^{-s\xi_i}
=
q_0^{-s}\prod_{i=1}^{n-1}(1+q_i^{-s}).
\]

对第 \(2\) 条，由独立性，
\[
\EE\!\left[e^{t\log K_n}\right]
=
\EE\!\left[K_n^t\right]
=
q_0^t\prod_{i=1}^{n-1}\EE\!\left[q_i^{t\xi_i}\right]
=
q_0^t\prod_{i=1}^{n-1}\frac{1+q_i^t}{2}.
\]

第 \(3\) 条则由
\[
\EE[\xi_i]=\frac12,
\qquad
\mathrm{Var}(\xi_i)=\frac14
\]
与独立性直接得到：
\[
\EE[\log K_n]
=
\log q_0+\sum_{i=1}^{n-1}\EE[\xi_i]\log q_i
=
\log q_0+\frac12\sum_{i=1}^{n-1}\log q_i,
\]
\[
\mathrm{Var}(\log K_n)
=
\sum_{i=1}^{n-1}\mathrm{Var}(\xi_i)(\log q_i)^2
=
\frac14\sum_{i=1}^{n-1}(\log q_i)^2.
\]

最后定义中心化独立变量
\[
X_i:=\left(\xi_i-\frac12\right)\log q_i
\in
\left[-\frac12\log q_i,\frac12\log q_i\right].
\]
则
\[
\log K_n-\EE[\log K_n]=\sum_{i=1}^{n-1}X_i,
\]
且每个 \(X_i\) 的取值区间长度为 \(\log q_i\)。由 Hoeffding 不等式，
\[
\PP\!\left(\left|\sum_{i=1}^{n-1}X_i\right|\ge u\right)
\le
2\exp\!\left(
-\frac{2u^2}{\sum_{i=1}^{n-1}(\log q_i)^2}
\right),
\]
即得第 \(4\) 条。证毕。
\end{proof}

\noindent
由此，Gödel--Tate 地址像在 Tate 模块中闭合为一个严格的 \(B\)-进 Cantor--Bernoulli 模型，而链内算子格则在平方自由算术化下闭合为一个 Boolean 整除格与对数可加的 Bernoulli 随机模型。前者把地址几何与移位动力学统一到同一条 clopen 自相似链上，后者则把纯序结构直接转写为可完全积化、可求矩并可作非渐近偏差估计的算术对象。

\endinput

\paragraph{Toeplitz 主子式的离散凹锥、Blaschke 原点模维数证书、Rouch\'e 边界裕度与 window-\(6\) parity 的几何压缩}
在 Toeplitz--Verblunsky 乘积分解、离切片 Blaschke 指标、Jensen 缺陷的视界纯度判据以及 window-\(6\) boundary uplift 的二层/三层分岔都已明确之后，还可继续抽取出下列五条后果。它们分别把 Toeplitz 主子式链在对数坐标下精确识别为离散凹锥，把单标量 \(|B_{\mathcal Z}(0)|\) 升格为残差维数的 sharp 证书，把 Hurwitz 型零自由逼近进一步有限化为一列圆周上的 Rouch\'e 裕度不等式，并把 window-\(6\) 的 \((\ZZ_2)^3\) boundary parity 先压缩为局域几何只保留的对角 \(\ZZ_2\) 方向，再识别为不能在“无额外协议门、无额外寄存器”的口径下跨 \(m=7,8\) 自然延拓的最小锚点现象。

\begin{theorem}[Toeplitz 主子式锥在对数坐标下恰为离散凹锥]\label{thm:xi-time-part65e-toeplitz-local-logconcavity-realizability}
\leanverified{paper\_xi\_time\_part65e\_toeplitz\_local\_logconcavity\_realizability}
设 \(N\ge 0\)，并给定正数列
\[
\Delta_{-1}:=1,
\qquad
\Delta_0>0,\ \Delta_1>0,\dots,\Delta_N>0.
\]
定义
\[
u_k:=\log \Delta_k
\qquad(-1\le k\le N).
\]
则以下两命题等价：
\begin{enumerate}
  \item 存在某个 Carath\'eodory 函数的截断 Toeplitz 主子式链
  \[
  \mathbf T_k=(c_{j-\ell})_{0\le j,\ell\le k}
  \qquad(0\le k\le N)
  \]
  使得
  \[
  \det \mathbf T_k=\Delta_k
  \qquad(0\le k\le N);
  \]
  \item 对所有 \(0\le k\le N-1\)，都有
  \[
  u_{k+1}-2u_k+u_{k-1}\le 0.
  \]
\end{enumerate}
并且在该情形下，对每个 \(0\le k\le N-1\) 都有严格恒等式
\[
\boxed{\;
u_{k+1}-2u_k+u_{k-1}
=
\log(1-|\alpha_k|^2)\le 0,
\;}
\]
从而相应 Verblunsky 系数的模长由主子式序列唯一恢复为
\[
|\alpha_k|^2
=
1-\exp\!\bigl(u_{k+1}-2u_k+u_{k-1}\bigr)
\qquad(0\le k\le N-1).
\]
换言之，Toeplitz--PSD 主子式链在忘却相位后恰等价于离散凹锥条件。
\end{theorem}
\begin{proof}
当 \(N=0\) 时结论平凡：取常值 Carath\'eodory 函数 \(\mathscr C\equiv \Delta_0\) 即可。以下设 \(N\ge 1\)。

若存在所述 Toeplitz 主子式链，则由定理 \ref{thm:xi-toeplitz-principal-minor-discrete-curvature-recovery}，对 \(k=0\) 有
\[
1-|\alpha_0|^2
=
\frac{\Delta_1}{\Delta_0^2}
=
\exp(u_1-2u_0+u_{-1}),
\]
而对 \(1\le k\le N-1\) 有
\[
1-|\alpha_k|^2
=
\frac{\Delta_{k+1}\Delta_{k-1}}{\Delta_k^2}
=
\exp(u_{k+1}-2u_k+u_{k-1}).
\]
两侧取对数即得
\[
u_{k+1}-2u_k+u_{k-1}
=
\log(1-|\alpha_k|^2)\le 0,
\]
并同时得到模长恢复公式。

反过来，若已知
\[
u_{k+1}-2u_k+u_{k-1}\le 0
\qquad(0\le k\le N-1),
\]
则可定义
\[
\rho_k:=\exp(u_{k+1}-2u_k+u_{k-1})\in(0,1],
\qquad
r_k:=\rho_k,
\qquad
|\alpha_k|:=\sqrt{1-\rho_k}\in[0,1)
\qquad(0\le k\le N-1),
\]
并任选相位；例如取全部 \(\alpha_k\in[0,1)\subset\RR\)。再令
\[
\alpha_j:=0
\qquad(j\ge N).
\]
由 Schur--Carath\'eodory 对应，存在某个 Carath\'eodory 函数 \(\mathscr C\) 具有这列 Verblunsky 系数。记其 Toeplitz 主子式为
\[
\widetilde\Delta_k:=\det \widetilde{\mathbf T}_k
\qquad(k\ge 0).
\]
由定理 \ref{thm:xi-toeplitz-principal-minor-discrete-curvature-recovery}，有
\[
\widetilde\Delta_0=\Delta_0,
\qquad
\widetilde\Delta_1
=
\Delta_0^2(1-|\alpha_0|^2)
=
\Delta_0^2\rho_0
=
\Delta_1.
\]
对每个 \(1\le k\le N-1\)，同一定理还给出递推
\[
\widetilde\Delta_{k+1}
=(1-|\alpha_k|^2)\frac{\widetilde\Delta_k^2}{\widetilde\Delta_{k-1}}
=r_k\frac{\widetilde\Delta_k^2}{\widetilde\Delta_{k-1}}.
\]
若已知 \(\widetilde\Delta_{k-1}=\Delta_{k-1}\) 与 \(\widetilde\Delta_k=\Delta_k\)，则
\[
\widetilde\Delta_{k+1}
=r_k\frac{\Delta_k^2}{\Delta_{k-1}}
=\Delta_{k+1}.
\]
归纳即得 \(\widetilde\Delta_k=\Delta_k\) 对所有 \(0\le k\le N\) 成立，故所需实现性得证。最后，模长唯一恢复公式已由上面的曲率恒等式直接给出。证毕。
\end{proof}

\begin{theorem}[Blaschke 原点模对离切片残差维数的 sharp 上界]\label{thm:xi-time-part65e-blaschke-origin-modulus-dimension-bound}
\leanverified{paper\_xi\_time\_part65e\_blaschke\_origin\_modulus\_dimension\_bound}
沿用定义 \ref{def:xi-offslice-ind-weight} 的记号，并设 \(\mathcal Z_{+}\neq\varnothing\) 为有限离切片零点多重集，\(B_{\mathcal Z}\) 为其对应的有限 Blaschke 乘积。定义
\[
\delta_{\min}(\mathcal Z)
:=
\min_{s\in\mathcal Z_{+}}
\Bigl(\Re(s)-\frac12\Bigr)
>0.
\]
则有 sharp 估计
\[
\mathrm{Ind}(\mathcal Z)
\le
\frac{\mathrm W(\mathcal Z)}{2L\,\delta_{\min}(\mathcal Z)}
=
\frac{-\log|B_{\mathcal Z}(0)|}{2L\,\delta_{\min}(\mathcal Z)}.
\]
等价地，
\[
\dim K_{B_{\mathcal Z}}
\le
\frac{-\log|B_{\mathcal Z}(0)|}{2L\,\delta_{\min}(\mathcal Z)}.
\]
并且等号成立当且仅当全部离切片零点具有相同的实部偏差，即
\[
\Re(s)-\frac12\equiv \delta_{\min}(\mathcal Z)
\qquad(s\in\mathcal Z_{+}).
\]
\end{theorem}
\begin{proof}
由定理 \ref{thm:xi-offslice-realpart-sum-law}，
\[
\mathrm W(\mathcal Z)
=
2L\sum_{s\in\mathcal Z_{+}}
\Bigl(\Re(s)-\frac12\Bigr).
\]
由于每一项都不小于 \(\delta_{\min}(\mathcal Z)\)，便有
\[
\mathrm W(\mathcal Z)\ge 2L\,\mathrm{Ind}(\mathcal Z)\,\delta_{\min}(\mathcal Z),
\]
即
\[
\mathrm{Ind}(\mathcal Z)
\le
\frac{\mathrm W(\mathcal Z)}{2L\,\delta_{\min}(\mathcal Z)}.
\]
再由定理 \ref{thm:xi-offslice-blaschke-origin-modulus}，
\[
\mathrm W(\mathcal Z)=-\log|B_{\mathcal Z}(0)|,
\]
从而得到第一条闭式。另一方面，命题 \ref{prop:xi-offslice-model-space-dimension} 给出
\[
\dim K_{B_{\mathcal Z}}=\mathrm{Ind}(\mathcal Z),
\]
遂得第二条闭式。

等号情形同样直接：上式中的和达到下界，当且仅当每个求和项都等于 \(\delta_{\min}(\mathcal Z)\)，亦即全部 \(\Re(s)-\tfrac12\) 完全相同。证毕。
\end{proof}

\begin{theorem}[Blaschke 原点模在 strip 约束下的残差维数下界]\label{thm:xi-time-part65e-strip-blaschke-origin-modulus-index-lower-bound}
\leanverified{paper\_xi\_time\_part65e\_strip\_blaschke\_origin\_modulus\_index\_lower\_bound}
沿用定义 \ref{def:xi-offslice-ind-weight} 的记号，并设 \(\mathcal Z_{+}\neq\varnothing\) 为有限离切片零点多重集，\(B_{\mathcal Z}\) 为其对应的有限 Blaschke 乘积。
若存在常数 \(\eta>0\) 使
\[
0<\Re(s)-\frac12\le \eta
\qquad(s\in\mathcal Z_{+}),
\]
则有 sharp 下界
\[
\boxed{\;
\mathrm{Ind}(\mathcal Z)
\ge
\frac{\mathrm W(\mathcal Z)}{2L\,\eta}
=
\frac{-\log|B_{\mathcal Z}(0)|}{2L\,\eta}.\;}
\]
等价地，
\[
\boxed{\;
\dim K_{B_{\mathcal Z}}
\ge
\frac{-\log|B_{\mathcal Z}(0)|}{2L\,\eta}.\;}
\]
并且等号成立当且仅当全部离切片零点具有相同的实部偏差
\[
\Re(s)-\frac12\equiv \eta
\qquad(s\in\mathcal Z_{+}).
\]
\end{theorem}
\begin{proof}
由定理 \ref{thm:xi-offslice-realpart-sum-law}，
\[
\mathrm W(\mathcal Z)
=
2L\sum_{s\in\mathcal Z_{+}}
\Bigl(\Re(s)-\frac12\Bigr).
\]
由于每一项都不超过 \(\eta\)，故
\[
\mathrm W(\mathcal Z)\le 2L\,\eta\,\mathrm{Ind}(\mathcal Z),
\]
从而
\[
\mathrm{Ind}(\mathcal Z)\ge \frac{\mathrm W(\mathcal Z)}{2L\,\eta}.
\]
再由定理 \ref{thm:xi-offslice-blaschke-origin-modulus}，
\(
\mathrm W(\mathcal Z)=-\log|B_{\mathcal Z}(0)|
\)，
即得第一条闭式。
另一方面，由命题 \ref{prop:xi-offslice-model-space-dimension}，
\[
\dim K_{B_{\mathcal Z}}=\mathrm{Ind}(\mathcal Z),
\]
于是第二条闭式随即成立。

等号成立当且仅当上式求和中的每一项都达到其上界 \(\eta\)，亦即全部 \(\Re(s)-\tfrac12\) 完全相同且等于 \(\eta\)。证毕。
\end{proof}

\begin{theorem}[Rouch\'e 边界裕度所给出的 RH 有限证书]\label{thm:xi-time-part65e-rouche-boundary-margin-rh-certificate}
沿用 \eqref{eq:xi-ramanujan-horizon-Fzeta} 的圆盘完成化函数 \(F_{\zeta}\)。设 \(\{F_N\}_{N\ge 1}\) 为一列解析函数 \(F_N:\mathbb{D}\to\CC\)，并满足定理 \ref{thm:xi-hurwitz-horizon-reduction} 的自对偶完成化一致性。若存在一列半径
\[
\rho_N\uparrow 1
\qquad(N\to\infty)
\]
使得对每个 \(N\) 都有
\[
F_N\text{ 在 }\overline{D(0,\rho_N)}\text{ 内无零点}
\]
且
\[
\sup_{|w|=\rho_N}|F_{\zeta}(w)-F_N(w)|
<
\inf_{|w|=\rho_N}|F_N(w)|,
\]
则 RH 成立。
\end{theorem}
\begin{proof}
固定 \(N\)。由于 \(F_N\) 在 \(\overline{D(0,\rho_N)}\) 内无零点，故
\[
\inf_{|w|=\rho_N}|F_N(w)|>0.
\]
于是由 Rouch\'e 定理，\(F_{\zeta}\) 与 \(F_N\) 在圆盘 \(D(0,\rho_N)\) 内具有相同的零点总数（按重数计）。而 \(F_N\) 在该圆盘内无零点，故 \(F_{\zeta}\) 在 \(D(0,\rho_N)\) 内亦无零点。

由于 \(\rho_N\uparrow 1\)，对任意 \(0<\rho<1\) 都可取 \(N\) 充分大使 \(\rho<\rho_N\)。于是 \(F_{\zeta}\) 在 \(D(0,\rho)\) 内无零点。因为这对每个 \(\rho<1\) 都成立，便知 \(F_{\zeta}\) 在整个 \(\mathbb D\) 内无零点。最后由定理 \ref{thm:app-rh-iff-disk-zero-free}，RH 成立。证毕。
\end{proof}

\begin{theorem}[window-\(6\) 的几何可实现 boundary parity 子群恰为对角 \texorpdfstring{$\ZZ_2$}{Z2}，且精确损失两比特]\label{thm:xi-time-part65e-window6-geometric-diagonal-z2-parity}
\leanverified{paper\_xi\_time\_part65e\_window6\_geometric\_diagonal\_z2\_parity}
记
\[
P_6:=C_{\mathrm{bdry}}\cong (\ZZ/2\ZZ)^3
\]
为命题 \ref{prop:xi-time-part60ab3-window6-boundary-center-residual} 的 boundary parity 中心子群。若只允许保持稳定类型标签且由 \(\Aut(Q_6)\) 实现的强局域 geometric uplift，则其几何可实现的 boundary parity 翻转子群恰为
\[
P_6^{\mathrm{geo}}=\langle\zeta\rangle\cong \ZZ/2\ZZ,
\]
其中 \(\zeta\) 为定理 \ref{thm:xi-time-part53da-boundary-z6-torsor} 中同时翻转三条两点纤维的全局 sheet-flip。更精确地，在坐标识别
\[
P_6\cong (\ZZ/2\ZZ)^3
\]
下有
\[
P_6^{\mathrm{geo}}=\langle(1,1,1)\rangle.
\]
因此
\[
P_6/P_6^{\mathrm{geo}}\cong (\ZZ/2\ZZ)^2.
\]
换言之，局域几何可实现性仅保留三位 boundary parity 的同步对角方向，并精确剥离两个独立的二值自由度。
\end{theorem}
\begin{proof}
由命题 \ref{prop:xi-time-part60ab3-window6-boundary-center-residual}，三条 boundary 二点纤维各自产生一个彼此独立的中心 sheet-flip，故
\[
P_6\cong (\ZZ/2\ZZ)^3.
\]
另一方面，由定理 \ref{thm:xi-time-part53da-boundary-z6-torsor}，\(\widetilde X_6^{\mathrm{bdry}}\) 上存在一个全局对合
\[
\zeta:\widetilde X_6^{\mathrm{bdry}}\to \widetilde X_6^{\mathrm{bdry}},
\]
它在每条两点纤维上都施加唯一非平凡交换，因此在上述三维 parity 坐标中正对应于同步翻转 \((1,1,1)\)。

再由结论 \ref{con:xi-terminal-window6-local-uplift-so10-unique}，在保持稳定类型标签 \(F_6\) 不变且由 \(\Aut(Q_6)\) 实现的强局域 uplift 口径下，除平凡位移外仅允许 \(\delta=\pm 34\)。而在 window-\(6\) 的 boundary uplift 中，同一纤维内两点之差正恒等于 \(34\)（定理 \ref{thm:xi-time-part53da-boundary-z6-torsor} 的设定），故这唯一允许的几何非平凡分支恰好就是 \(\zeta\)。因此
\[
P_6^{\mathrm{geo}}=\langle \zeta\rangle=\langle(1,1,1)\rangle\cong \ZZ/2\ZZ.
\]
最后，\((\ZZ/2\ZZ)^3\) 除以其一维对角子群后得到
\[
(\ZZ/2\ZZ)^3/\langle(1,1,1)\rangle\cong (\ZZ/2\ZZ)^2,
\]
从而 cokernel 精确具有两位二进自由度。证毕。
\end{proof}

\begin{theorem}[window-\(6\) boundary parity 的无协议门跨尺度 \(2\)-挠中心延拓不存在]\label{thm:xi-time-part65e-window6-no-gateless-z2-central-extension}
\leanverified{paper\_xi\_time\_part65e\_window6\_no\_gateless\_z2\_central\_extension}
不存在某个固定紧连通群 \(G\) 与一族仅依赖于分辨率的规则
\[
\mathfrak F_m:X_m^{\mathrm{bdry}}\longrightarrow Z(G)[2]
\qquad(m\in\{6,7,8\}),
\]
其中
\[
Z(G)[2]:=\{z\in Z(G):z^2=e\},
\]
使得同时满足下列三条条件：
\begin{enumerate}
  \item 在 \(m=6\) 时，\(\mathfrak F_6\) 恢复推论 \ref{cor:terminal-delta-parity-rule} 中三条 boundary 纤维所给出的独立 \((\ZZ_2)^3\)-sheet parity；
  \item 该规则对 dyadic uplift 的 boundary 纤维置换保持自然；
  \item 构造过程中不引入任何额外协议门或额外读出寄存器。
\end{enumerate}
\end{theorem}
\begin{proof}
反设存在这样一族规则。条件 (1) 表明：在 \(m=6\) 的 boundary 扇区上，每条二点纤维都被赋予一个非平凡的二值中心选择；这正是推论 \ref{cor:terminal-delta-parity-rule} 中由
\[
\Delta_6^{\mathrm{bdry}}=\{0,34\}
\]
所诱导的 sheet parity。若条件 (2) 与 (3) 也同时成立，则这一二值对象必须沿 dyadic uplift 在更高分辨率的 boundary 纤维上继续以“置换自然且无需外加侧信息”的方式存在。

然而，命题 \ref{prop:terminal-sign-covariant-label-boundary-bifurcation} 已经给出严格判据：这样的符号协变二值对象存在当且仅当纤维基数等于 \(2\)。同一命题并明确指出，在 boundary uplift 口径下有
\[
m=6:\ \#(\mathrm{Fold}^{\mathrm{bin}}_6)^{-1}(w)=2,
\]
\[
m=7,8:\ \#(\mathrm{Fold}^{\mathrm{bin}}_m)^{-1}(w)=3,
\qquad
\Delta_7^{\mathrm{bdry}}=\{0,55,89\},
\qquad
\Delta_8^{\mathrm{bdry}}=\{0,89,144\}.
\]
因此 \(m=7,8\) 的三点 boundary 纤维上根本不存在所需的置换自然二值延拓，这与条件 (2)--(3) 矛盾。故这样的统一规则不存在。

最后，备注 \ref{rem:window6-sheet-parity-nonpersistent} 已指出：若欲把 window-\(6\) 的 \(\ZZ_2\) sheet parity 提升为跨尺度残余选择律，就必须额外施加可审计协议门以重新锁定二层覆盖，或引入额外读出寄存器。并且定理 \ref{thm:xi-boundary-uplift-mixed-radix-register-lower-bound} 进一步量化了这一下限：若要求完整继承 window-\(6\) 的三位 sheet parity，则必有
\[
r_{\mathrm{ext}}(7)\ge 11,
\qquad
r_{\mathrm{ext}}(8)\ge 16.
\]
证毕。
\end{proof}

\noindent
由此，Toeplitz 主子式、离切片 Blaschke 数据、圆盘完成化边界比较与 window-\(6\) boundary uplift 这五条原本分属不同模块的链条，都被进一步压缩为有限局域证书：Toeplitz 可实现性不再需要搜索整个正定锥，而可在对数坐标下完全转写为离散凹锥；Blaschke 原点模不再只是总偏离负载，而在给定最小偏离阈值时直接控制残差维数；RH 的一个充分条件不再需要整列紧集一致逼近，而只需一列逼近边界的圆周裕度不等式；至于 window-\(6\) 的 \((\ZZ_2)^3\) boundary parity，则在局域几何 uplift 口径下仅保留同步对角 \(\ZZ_2\) 方向，而其跨尺度延拓仍必须显式支付新的协议成本。

\endinput

\paragraph{部分屏幕观测的后验熵、唯一回放阈值、压缩树补全与单面边际律}
在部分边界观测的相对同调核、屏幕图连通分量公式、最小审计补全成本与可见秩函数的拟阵结构已经确立之后，这条主线还可继续压缩为下列五条直接后果。它们分别把 \(\FF_2\)-后验分布写成完全显式的公平比特模型，把唯一回放屏幕的极小阈值识别为全局边界多重图的生成树阈值，把任意部分屏幕的极小精确化补全面集完全刻画为压缩多重图的生成树，把带权精确化编译为最小生成树问题，并把单张观测面的边际信息增益提升为严格的 \(0\text{--}1\) 律。

\begin{theorem}[部分屏幕观测的后验分布完全等价于独立公平比特]\label{thm:xi-time-part65f-screen-posterior-fair-bits}
\leanverified{paper\_xi\_time\_part65f\_screen\_posterior\_fair\_bits}
固定 \(m,n\ge 1\)，并沿用定理 \ref{thm:spg-partial-boundary-screen-kernel-relative-homology}、定理 \ref{thm:spg-screen-kernel-connected-components} 与推论 \ref{cor:spg-partial-boundary-fiber-cardinality-f2} 的记号。将系数域取为 \(\FF_2\)，令
\[
f_S:C_n(\mathcal Q_m;\FF_2)\longrightarrow \FF_2^{S},
\qquad
Y:=f_S(X),
\]
其中 \(X\) 在 \(C_n(\mathcal Q_m;\FF_2)\) 上服从均匀分布。记
\[
c_S:=c(\Gamma_S).
\]
则对每个 \(y\in \operatorname{im}(f_S)\)，纤维 \(f_S^{-1}(y)\) 是维数 \(c_S-1\) 的仿射 \(\FF_2\)-子空间；特别地，\(f_S^{-1}(y)\) 与 \(\FF_2^{\,c_S-1}\) 仿射同构。

若以下 Shannon 熵与互信息均以 \(\log_2\) 为底，则有
\[
H(X\mid Y)=c_S-1,
\]
\[
H(Y)=2^{mn}-c_S+1,
\]
\[
I(X;Y)=2^{mn}-c_S+1.
\]
\end{theorem}
\begin{proof}
顶维链群 \(C_n(\mathcal Q_m;\FF_2)\) 的维数等于顶维胞腔数 \(2^{mn}\)，故
\[
H(X)=\log_2|C_n(\mathcal Q_m;\FF_2)|=2^{mn}.
\]
另一方面，定理 \ref{thm:spg-screen-kernel-connected-components} 的“分量常值”论证在系数域 \(\FF_2\) 上逐字适用，从而得到
\[
\ker f_S\cong \FF_2^{\,c_S-1}.
\]
由于 \(f_S\) 为 \(\FF_2\)-线性映射，每个非空纤维 \(f_S^{-1}(y)\) 都是某个 \(\ker f_S\) 的仿射平移，因此与 \(\FF_2^{\,c_S-1}\) 仿射同构。

又由推论 \ref{cor:spg-partial-boundary-fiber-cardinality-f2}，
\[
|f_S^{-1}(y)|=2^{\dim_{\FF_2}\ker f_S}=2^{c_S-1}
\qquad
(y\in\operatorname{im}(f_S)),
\]
右端与 \(y\) 无关。故在均匀先验下，条件分布 \(X\mid Y=y\) 在纤维上均匀，因而
\[
H(X\mid Y=y)=\log_2|f_S^{-1}(y)|=c_S-1
\]
对所有可达 \(y\) 都成立，从而
\[
H(X\mid Y)=c_S-1.
\]

由链式法则，
\[
H(Y)=H(X)-H(X\mid Y)=2^{mn}-c_S+1.
\]
最后，\(Y=f_S(X)\) 是 \(X\) 的确定函数，因此
\[
I(X;Y)=H(Y)-H(Y\mid X)=H(Y)=2^{mn}-c_S+1.
\]
证毕。
\end{proof}

\begin{theorem}[唯一回放屏幕的极小规模与生成树刻画]\label{thm:xi-time-part65f-exact-screen-spanning-tree-threshold}
沿用定理 \ref{thm:spg-screen-kernel-connected-components} 的边界多重图 \(\Gamma\) 记号：其顶点集为全部 \(n\)-胞腔外加外部顶点 \(\infty\)，并把每个可观测取向 \((n-1)\)-面视为 \(\Gamma\) 的一条边；内部面连接其两侧顶维胞腔，外边界面连接相应胞腔与 \(\infty\)。对任意屏幕 \(S\subset \overrightarrow{\mathcal F_m^{n-1}}\)，记其所诱导的生成子多重图为 \(\Gamma_S\)。

则成立：
\begin{enumerate}
  \item \(f_S\) 单射当且仅当 \(\Gamma_S\) 连通；
  \item 任意使 \(f_S\) 单射的屏幕都满足
        \[
        |S|\ge 2^{mn};
        \]
  \item 等号成立当且仅当 \(\Gamma_S\) 是 \(\Gamma\) 的一棵生成树；
  \item 因而极小唯一回放屏幕的总数恰为 \(\tau(\Gamma)\)，其中 \(\tau(\Gamma)\) 表示 \(\Gamma\) 的生成树数。
\end{enumerate}
\end{theorem}
\begin{proof}
第 \(1\) 条正是推论 \ref{cor:spg-screen-injectivity-and-audit-cost-components} 的内容。

由于 \(\mathcal Q_m\) 的顶维胞腔数为 \(2^{mn}\)，故
\[
|V(\Gamma)|=2^{mn}+1.
\]
任意使 \(f_S\) 单射的屏幕都对应一个连通 spanning sub-multigraph \(\Gamma_S\)。而任意连通图至少有 \(|V|-1\) 条边，因此
\[
|S|=|E(\Gamma_S)|\ge |V(\Gamma)|-1=2^{mn}.
\]
这就得到第 \(2\) 条。

若 \(|S|=2^{mn}\) 且 \(f_S\) 单射，则 \(\Gamma_S\) 连通并且
\[
|E(\Gamma_S)|=|V(\Gamma)|-1,
\]
故 \(\Gamma_S\) 必为生成树。反之，任一生成树都连通且边数恰为 \(|V(\Gamma)|-1=2^{mn}\)，于是对应屏幕必为极小唯一回放屏幕。于是第 \(3\) 条成立。

第 \(4\) 条随即成立：极小唯一回放屏幕与 \(\Gamma\) 的生成树一一对应，其总数即为 \(\tau(\Gamma)\)。证毕。
\end{proof}
\leanverified{paper\_xi\_time\_part65f\_exact\_screen\_spanning\_tree\_threshold}

\begin{theorem}[任意部分屏幕的极小精确化补全面集由压缩图生成树完全刻画]\label{thm:xi-time-part65f-screen-quotient-tree-completion}
沿用定理 \ref{thm:spg-screen-kernel-connected-components} 的记号。给定任意部分屏幕 \(S\subset \overrightarrow{\mathcal F_m^{n-1}}\)，设
\[
\Gamma_S=\bigsqcup_{i=1}^{c}\mathcal C_i,
\qquad
c:=c(\Gamma_S).
\]
将每个连通分量 \(\mathcal C_i\) 缩并为一个顶点，并把 \(\Gamma\setminus E(\Gamma_S)\) 中连接不同分量的边保留下来、删去缩并后产生的回路，得到压缩多重图
\[
\Gamma/\Gamma_S.
\]
以下把每条跨分量补全边与其在 \(\Gamma/\Gamma_S\) 中对应的边仍记为同一符号。
则成立：
\begin{enumerate}
  \item 同一个整数 \(c-1\) 同时等于最小圆维审计补全成本与最少追加观测面数，即
        \[
        \min_{(R,r_S)}\cdim(R)
        =
        c-1
        =
        \min\Bigl\{
        |T|:\ T\subseteq E(\Gamma)\setminus E(\Gamma_S),\ \Gamma_{S\cup T}\text{ 连通}
        \Bigr\};
        \]
  \item 对任意 \(T\subseteq E(\Gamma)\setminus E(\Gamma_S)\) 且 \(|T|=c-1\)，屏幕 \(S\cup T\) 达到唯一回放当且仅当 \(T\) 在压缩图 \(\Gamma/\Gamma_S\) 中的像是一棵生成树；
  \item 因而全部极小精确化补全面集的总数恰为
        \[
        \tau(\Gamma/\Gamma_S).
        \]
\end{enumerate}
\end{theorem}
\begin{proof}
第一条中的左端等式正是推论 \ref{cor:spg-screen-injectivity-and-audit-cost-components}。

现证右端等式。由于 \(\Gamma_S\) 恰有 \(c\) 个连通分量，若要通过继续添加观测面使 \(\Gamma_{S\cup T}\) 连通，则 \(T\) 在压缩图上必须把这 \(c\) 个缩并顶点连成一个连通图。任意 \(c\) 点连通图至少含 \(c-1\) 条边，因此
\[
|T|\ge c-1.
\]
另一方面，压缩图任取一棵生成树，其对应的原图边集 \(T\) 恰含 \(c-1\) 条边，并使 \(\Gamma_{S\cup T}\) 连通，从而上界可达。故右端最小值也等于 \(c-1\)。

再证第 \(2\) 条。若 \(|T|=c-1\) 且 \(\Gamma_{S\cup T}\) 连通，则 \(T\) 在压缩图上的像给出一个 \(c\) 点连通图，且边数恰为 \(c-1\)，故必为生成树。反之，若该像是一棵生成树，则对应的原图边集把 \(c\) 个缩并顶点连通，从而 \(\Gamma_{S\cup T}\) 连通，推论 \ref{cor:spg-screen-injectivity-and-audit-cost-components} 遂给出 \(f_{S\cup T}\) 单射。

第 \(3\) 条随之成立：压缩图的每棵生成树与一个极小精确化补全面集唯一对应，故总数恰为 \(\tau(\Gamma/\Gamma_S)\)。证毕。
\end{proof}
\leanverified{paper\_xi\_time\_part65f\_screen\_quotient\_tree\_completion}

\begin{theorem}[带权精确化可完全编译为压缩图的最小生成树问题]\label{thm:xi-time-part65f-weighted-screen-completion-mst}
\leanverified{paper\_xi\_time\_part65f\_weighted\_screen\_completion\_mst}
在定理 \ref{thm:xi-time-part65f-screen-quotient-tree-completion} 的设定下，对每条补全边赋予非负代价
\[
\mathrm{cost}:E(\Gamma)\setminus E(\Gamma_S)\longrightarrow \RR_{\ge 0}.
\]
对任意补全面集 \(T\subseteq E(\Gamma)\setminus E(\Gamma_S)\)，定义总代价
\[
\mathrm{Cost}(T):=\sum_{e\in T}\mathrm{cost}(e).
\]
则
\[
\min\Bigl\{\mathrm{Cost}(T):\ \Gamma_{S\cup T}\text{ 连通}\Bigr\}
=
\min\Bigl\{\mathrm{Cost}(T):\ T\text{ 在 }\Gamma/\Gamma_S\text{ 中是一棵生成树}\Bigr\}.
\]
特别地：
\begin{enumerate}
  \item 任一压缩图最小生成树对应的补全面集都给出成本最优的精确化方案；
  \item 反之，任一按包含极小的成本最优补全面集，其压缩像必为 \(\Gamma/\Gamma_S\) 的一棵最小生成树。
\end{enumerate}
\end{theorem}
\begin{proof}
任取使 \(\Gamma_{S\cup T}\) 连通的补全面集 \(T\)。其在压缩图上的像是一个连通生成子多重图。取其任意生成树 \(T_{\mathrm{tr}}\subseteq T\)。由于代价非负，
\[
\mathrm{Cost}(T_{\mathrm{tr}})\le \mathrm{Cost}(T).
\]
因此，在全部可行补全面集中最小化总代价，等价于在压缩图的全部生成树中最小化总代价。第一个等式得证。

于是任一压缩图最小生成树都直接给出成本最优补全，这证明了第 \(1\) 条。

若 \(T\) 是按包含极小的成本最优补全面集，而其压缩像含有回路，则删去该回路上一条边后图仍连通；由于边权非负，总代价不会增加。若总代价严格减小，则与成本最优性矛盾；若总代价不变，则与按包含极小矛盾。故其压缩像无回路。既然它仍连通并覆盖全部缩并顶点，便是一棵生成树；再由成本最优性可知它必是最小生成树。第 \(2\) 条成立。证毕。
\end{proof}

\begin{theorem}[单张观测面的边际信息增益具有严格的 \(0\text{--}1\) 律]\label{thm:xi-time-part65f-single-face-zero-one-marginal-law}
\leanverified{paper\_xi\_time\_part65f\_single\_face\_zero\_one\_marginal\_law}
沿用定理 \ref{thm:spg-screen-rank-matroid-supermodularity} 的记号，令
\[
r(S):=\rank_{\QQ}\bigl(\operatorname{im}(f_S\otimes\QQ)\bigr),
\qquad
\delta(S):=\rank_{\QQ}\bigl(\ker(f_S\otimes\QQ)\bigr).
\]
再把任一未观测取向 \((n-1)\)-面与其在 \(\Gamma\) 中对应的边同记为 \(e\in E(\Gamma)\setminus E(\Gamma_S)\)。若 \(e\) 的两个端点分别落在 \(\Gamma_S\) 的连通分量 \(\mathcal C_i,\mathcal C_j\) 中，则有精确公式
\[
\delta(S)-\delta(S\cup\{e\})
=
\begin{cases}
1,& \mathcal C_i\neq \mathcal C_j,\\[4pt]
0,& \mathcal C_i=\mathcal C_j,
\end{cases}
\]
等价地，
\[
r(S\cup\{e\})-r(S)
=
\begin{cases}
1,& \mathcal C_i\neq \mathcal C_j,\\[4pt]
0,& \mathcal C_i=\mathcal C_j.
\end{cases}
\]
\end{theorem}
\begin{proof}
由定理 \ref{thm:spg-screen-kernel-connected-components} 与推论 \ref{cor:spg-screen-injectivity-and-audit-cost-components}，
\[
\delta(S)=c(\Gamma_S)-1.
\]
若 \(e\) 的两个端点位于同一连通分量中，则向 \(\Gamma_S\) 添加 \(e\) 不改变连通分量数，因此
\[
c(\Gamma_{S\cup\{e\}})=c(\Gamma_S),
\qquad
\delta(S\cup\{e\})=\delta(S).
\]

若 \(e\) 连接两个不同分量，则添加 \(e\) 恰好把这两个分量合并，从而
\[
c(\Gamma_{S\cup\{e\}})=c(\Gamma_S)-1,
\]
进而
\[
\delta(S\cup\{e\})=c(\Gamma_{S\cup\{e\}})-1=\delta(S)-1.
\]
这就得到第一条 \(0\text{--}1\) 公式。

再由定理 \ref{thm:spg-screen-rank-matroid-supermodularity} 的秩亏恒等式
\[
\delta(S)=\rank_{\QQ}C_n(\mathcal Q_m;\QQ)-r(S)
=
2^{mn}-r(S),
\]
立刻推出
\[
r(S\cup\{e\})-r(S)=\delta(S)-\delta(S\cup\{e\}),
\]
从而得到第二条公式。证毕。
\end{proof}

\noindent
综上，部分边界观测这条链已经在概率、图论、优化与拟阵四个层面同时闭合：纤维常数不再只是一个计数事实，而是完整的后验公平比特模型；唯一回放屏幕的极小阈值不再只是 \(2^{mn}\) 的抽象秩值，而是全局边界多重图的生成树阈值；任意部分屏幕的极小精确化不再只是“还差 \(c(\Gamma_S)-1\) 张面”，而是由压缩图的生成树与最小生成树完全组织；与此同时，单张新观测面所带来的信息秩增量也不再只是超模意义下的平均改善，而是具有严格可判定的 \(0\text{--}1\) 局域律。

\endinput

\paragraph{局部化奇异环的 \(S\)-根单值化、逆极限相位账本、Fold 相位游走与 Dirichlet 轴交换阻断}
在有限 prime-register solenoid 的 Kronecker 读出、有理玫瑰/Lissajous 的几何提升、局部化相位采样的素数恢复与 Fold 谱的乘积分解已经确立之后，还可把 singular ring 侧的根提升、分母一致性、干涉能量的概率读出以及 Dirichlet 频率上的换轴障碍继续压缩为下列一组闭式后果。其核心是：\(\Sigma_S=\widehat{\ZZ[S^{-1}]}\) 不仅是圆周商上的 prime-register 扩张，而且恰是全部 \(S\)-光滑有理幂的普适单值化对象；一旦把这一普适对象与 Fold 的乘积谱和 Dirichlet 频率的线性无关性并置，分母账本、概率干涉与换轴阻断三条链便可统一落到同一相位审计框架之内。

\begin{theorem}[奇异环上的 \(S\)-根提升与单位圆有理幂的普适单值化]\label{thm:xi-time-part65g-solenoid-s-root-universal-single-valuedization}
沿用定理 \ref{thm:cdim-star-z1s-dual-extension} 的记号。设 \(S\subset\mathbb P\) 为有限素数集，并记
\[
G_S:=\ZZ[S^{-1}],
\qquad
\Sigma_S:=\widehat{G_S},
\qquad
\chi_q(x):=x(q)
\quad (q\in G_S,\ x\in \Sigma_S).
\]
令
\[
\pi:=\chi_1:\Sigma_S\to \TT.
\]
则 \(\pi\) 为满同态。对每个 \(S\)-光滑正整数 \(d\)，定义
\[
\rho_d:=\chi_{1/d}:\Sigma_S\to \TT.
\]
则对任意 \(x\in \Sigma_S\) 都有
\[
(\rho_d(x))^d=\pi(x),
\]
并且若 \(q=m/d\in G_S\) 为既约分式且 \(d\) 为 \(S\)-光滑，则
\[
\chi_q(x)=\rho_d(x)^m.
\]

更进一步，\(\Sigma_S\) 在如下意义下具有普适性：若 \(K\) 为紧交换群，存在满同态
\[
\pi_K:K\twoheadrightarrow \TT
\]
与一族连续同态
\[
\rho_d^K:K\to \TT
\]
（\(d\) 遍历全部 \(S\)-光滑正整数）满足
\[
(\rho_d^K)^d=\pi_K,
\qquad
(\rho_{d'}^K)^{d'/d}=\rho_d^K
\quad (d\mid d'),
\]
则存在唯一连续满同态
\[
F:K\twoheadrightarrow \Sigma_S
\]
使得
\[
\pi_K=\pi\circ F,
\qquad
\rho_d^K=\rho_d\circ F
\]
对全部 \(S\)-光滑 \(d\) 都成立。
\end{theorem}
\begin{proof}
由字符乘法法则，\(\chi_{q_1}\chi_{q_2}=\chi_{q_1+q_2}\)。因此对每个 \(S\)-光滑 \(d\) 都有
\[
(\rho_d)^d
=
\chi_{1/d}^d
=
\chi_{d\cdot (1/d)}
=
\chi_1
=
\pi,
\]
而对 \(q=m/d\in G_S\) 则有
\[
\chi_q
=
\chi_{m\cdot(1/d)}
=
\chi_{1/d}^{\,m}
=
\rho_d^{\,m}.
\]

由定理 \ref{thm:cdim-star-z1s-dual-extension}，\(\Sigma_S\) 连通。又 \(\chi_1\) 非平凡，故其像是 \(\TT\) 的非平凡闭连通子群，只能等于 \(\TT\)。于是 \(\pi\) 满射。

对普适性，令 \(\gamma_d\in K^\wedge\) 为 \(\rho_d^K\) 的对偶元素。则
\[
d\gamma_d=\gamma_1,
\qquad
\frac{d'}{d}\gamma_{d'}=\gamma_d
\quad (d\mid d').
\]
定义映射
\[
\iota:G_S\longrightarrow K^\wedge,
\qquad
\iota\!\left(\frac{m}{d}\right):=m\gamma_d.
\]
需证其良定义。若 \(m/d=m'/d'\)，取 \(S\)-光滑公倍数 \(\ell\) 使 \(d\mid \ell\)、\(d'\mid \ell\)。则
\[
m\gamma_d
=
m\frac{\ell}{d}\gamma_\ell
=
\frac{m\ell}{d}\gamma_\ell
=
\frac{m'\ell}{d'}\gamma_\ell
=
m'\frac{\ell}{d'}\gamma_\ell
=
m'\gamma_{d'},
\]
故 \(\iota\) 良定义。并且
\[
\iota\!\left(\frac{m}{d}+\frac{m'}{d'}\right)
=
\iota\!\left(\frac{m\ell/d+m'\ell/d'}{\ell}\right)
=
\left(\frac{m\ell}{d}+\frac{m'\ell}{d'}\right)\gamma_\ell
=
m\gamma_d+m'\gamma_{d'}
\]
说明 \(\iota\) 为群同态。

再证 \(\iota\) 单射。若 \(\iota(m/d)=0\)，则 \(m\gamma_d=0\)，等价于
\[
(\rho_d^K)^m=1.
\]
而 \(\rho_d^K(K)\) 是 \(\TT\) 的闭子群，且其 \(d\)-次幂像等于 \(\pi_K(K)=\TT\)。因此 \(\rho_d^K(K)\) 不可能是有限循环群，只能等于 \(\TT\)。于是 \(z^m=1\) 对所有 \(z\in\TT\) 都成立，故 \(m=0\)，从而 \(m/d=0\)。所以 \(\iota\) 单射。

Pontryagin 对偶化后得到连续满同态
\[
F:=\widehat{\iota}:K\twoheadrightarrow \Sigma_S.
\]
对任意 \(S\)-光滑 \(d\)，\(\rho_d\circ F\) 的对偶元素正是
\[
\iota(1/d)=\gamma_d,
\]
故 \(\rho_d\circ F=\rho_d^K\)。同理 \(\pi\circ F=\pi_K\)。唯一性则来自
\[
\widehat{\Sigma_S}\cong G_S
\]
由全部元素 \(m/d\) 生成，而 \(\chi_{m/d}=\rho_d^m\) 已完全由上述相容关系决定。证毕。
\end{proof}

\begin{corollary}[有理玫瑰与 Lissajous 在奇异环上的同态提升]\label{cor:xi-time-part65g-rose-lissajous-singular-ring-homomorphic-lift}
沿用定理 \ref{thm:xi-time-part65g-solenoid-s-root-universal-single-valuedization} 与定理 \ref{thm:xi-time-part62d-solenoid-kronecker-rose-lissajous-lift} 的记号，记
\[
\iota_S:=\kappa_S:\RR\longrightarrow \Sigma_S,
\qquad
\iota_S(t)(q):=e^{2\pi iqt}
\quad (q\in G_S).
\]
取 \(S\)-光滑正整数 \(d\) 与整数 \(n\) 满足 \(\gcd(n,d)=1\)，定义
\[
\mathcal R_{n,d}(x):=\frac12\Bigl(\rho_d(x)^{d+n}+\rho_d(x)^{d-n}\Bigr).
\]
则对任意 \(\theta\in\RR\) 都有
\[
\mathcal R_{n,d}\bigl(\iota_S(\theta/2\pi)\bigr)
=
e^{i\theta}\cos\!\Bigl(\frac{n}{d}\theta\Bigr).
\]
因此有理玫瑰轨道
\[
\theta\longmapsto e^{i\theta}\cos\!\Bigl(\frac{n}{d}\theta\Bigr)
\]
恰是奇异环字符读出的单值提升。

若 \(a=a_0/d\)、\(b=b_0/d\in G_S\) 且 \(a_0,b_0\in\ZZ\)，定义
\[
\mathcal L_{a,b,\delta}(x):=
\Bigl(
\Re\bigl(e^{i\delta}\rho_d(x)^{a_0}\bigr),\
\Re\bigl(\rho_d(x)^{b_0}\bigr)
\Bigr).
\]
则对任意 \(t\in\RR\) 都有
\[
\mathcal L_{a,b,\delta}\bigl(\iota_S(t/2\pi)\bigr)
=
\bigl(\cos(at+\delta),\cos(bt)\bigr).
\]
故一切分母仅由 \(S\) 中素数支撑的有理玫瑰与有理 Lissajous 轨道，都可在 \(\Sigma_S\) 上以单值同态方式实现。
\end{corollary}
\begin{proof}
由定理 \ref{thm:xi-time-part62d-solenoid-kronecker-rose-lissajous-lift}，对任意 \(q\in G_S\) 与 \(u\in\RR\) 都有
\[
\chi_q\bigl(\iota_S(u/2\pi)\bigr)=e^{iqu}.
\]
特别地，
\[
\rho_d\bigl(\iota_S(u/2\pi)\bigr)=\chi_{1/d}\bigl(\iota_S(u/2\pi)\bigr)=e^{iu/d}.
\]
于是取 \(u=\theta\) 得
\[
\mathcal R_{n,d}\bigl(\iota_S(\theta/2\pi)\bigr)
=
\frac12\Bigl(
e^{i(d+n)\theta/d}+e^{i(d-n)\theta/d}
\Bigr)
=
e^{i\theta}\cos\!\Bigl(\frac{n}{d}\theta\Bigr).
\]
这证明了玫瑰公式。

同理，取 \(u=t\) 得
\[
\rho_d\bigl(\iota_S(t/2\pi)\bigr)^{a_0}=e^{ia_0t/d}=e^{iat},
\qquad
\rho_d\bigl(\iota_S(t/2\pi)\bigr)^{b_0}=e^{ib_0t/d}=e^{ibt},
\]
故
\[
\mathcal L_{a,b,\delta}\bigl(\iota_S(t/2\pi)\bigr)
=
\bigl(\Re(e^{i\delta}e^{iat}),\Re(e^{ibt})\bigr)
=
\bigl(\cos(at+\delta),\cos(bt)\bigr).
\]
证毕。
\end{proof}

\begin{proposition}[局部化奇异环的逆极限相位账本刻画]\label{prop:xi-time-part65g-solenoid-inverse-limit-phase-ledger}
\leanverified{paper\_xi\_time\_part65g\_solenoid\_inverse\_limit\_phase\_ledger}
沿用定理 \ref{thm:xi-time-part65g-solenoid-s-root-universal-single-valuedization} 的记号。记
\[
\mathcal D_S:=\left\{
d\in\NN_{\ge 1}:\ \text{\(d\) 的全部素因子都属于 }S
\right\},
\]
并按整除关系赋予其有向结构。对 \(d\mid d'\) 定义连接映射
\[
\varphi_{d',d}:\TT\to\TT,
\qquad
\varphi_{d',d}(z):=z^{d'/d}.
\]
则映射
\[
\Sigma_S\longrightarrow \varprojlim_{d\in\mathcal D_S}\bigl(\TT,\varphi_{d',d}\bigr),
\qquad
x\longmapsto \bigl(\rho_d(x)\bigr)_{d\in\mathcal D_S}
\]
是拓扑同构。

等价地，一个相位账本 \((z_d)_{d\in\mathcal D_S}\in \TT^{\mathcal D_S}\) 来自某个 \(x\in\Sigma_S\) 当且仅当它满足一致性
\[
z_d=z_{d'}^{\,d'/d}
\qquad (d\mid d').
\]
在该情形下，相应字符满足
\[
x\!\left(\frac{m}{d}\right)=z_d^{\,m}
\qquad
\left(\frac{m}{d}\in G_S\right).
\]
\end{proposition}
\begin{proof}
有直极限分解
\[
G_S
=
\bigcup_{d\in\mathcal D_S}\frac{1}{d}\ZZ
=
\varinjlim_{d\in\mathcal D_S}\frac{1}{d}\ZZ,
\]
其中当 \(d\mid d'\) 时，\((1/d)\ZZ\hookrightarrow (1/d')\ZZ\) 为自然包含。Pontryagin 对偶把离散群直极限送为紧群逆极限，因此
\[
\Sigma_S
\cong
\varprojlim_{d\in\mathcal D_S}\widehat{(1/d)\ZZ}.
\]
把每个 \(\widehat{(1/d)\ZZ}\) 识别为 \(\TT\) 后，自然包含的对偶恰为
\[
z\longmapsto z^{d'/d},
\]
故得到所示逆系统。

显式地，若 \(x\in\Sigma_S\)，记
\[
z_d:=x(1/d)=\rho_d(x).
\]
当 \(d\mid d'\) 时，
\[
z_{d'}^{\,d'/d}
=
x(1/d')^{d'/d}
=
x\!\left(\frac{d'}{d}\cdot \frac{1}{d'}\right)
=
x(1/d)
=
z_d,
\]
故 \((z_d)\) 一致。

反之，给定一致族 \((z_d)_{d\in\mathcal D_S}\)，对任意 \(q=m/d\in G_S\) 定义
\[
x(q):=z_d^{\,m}.
\]
若 \(m/d=m'/d'\)，取 \(\ell\in\mathcal D_S\) 使 \(d\mid \ell\)、\(d'\mid \ell\)。则由一致性
\[
z_d^{\,m}
=
\left(z_\ell^{\,\ell/d}\right)^m
=
z_\ell^{\,m\ell/d}
=
z_\ell^{\,m'\ell/d'}
=
\left(z_\ell^{\,\ell/d'}\right)^{m'}
=
z_{d'}^{\,m'},
\]
故定义与表示无关。并且
\[
x(q+q')=x(q)x(q')
\]
由同分母计算立即可得，所以 \(x\in\operatorname{Hom}(G_S,\TT)=\Sigma_S\)。显然这两个构造互为逆映射。紧 Hausdorff 范畴中的连续双射即为同胚，故结论成立。证毕。
\end{proof}

\begin{theorem}[有限频率集的根闭包最小载体与 prime-register 共同核的素数恢复]\label{thm:xi-time-part65g-frequency-prime-support-kernel-recovery}
\leanverified{paper\_xi\_time\_part65g\_frequency\_prime\_support\_kernel\_recovery}
设 \(Q=\{q_1,\dots,q_r\}\subset\QQ^\times\) 为有限频率族，并记 \(S(Q)\subset\mathbb P\) 为其既约分母中出现的全部素数。再记
\[
\Sigma_{S(Q)}:=\widehat{\ZZ[S(Q)^{-1}]},
\qquad
K_{S(Q)}:=\prod_{p\in S(Q)}\ZZ_p.
\]
则成立：
\begin{enumerate}
  \item 若一个紧交换群 \(K\) 配备满同态 \(\pi_K:K\twoheadrightarrow \TT\) 及其对全部 \(S(Q)\)-光滑正整数的兼容根提升 \(\{\rho_d^K\}\)，并且对每个 \(q_j=m_j/d_j\)（既约）都以
  \[
  (\rho_{d_j}^K)^{m_j}
  \]
  作为对应频率 \(q_j\) 的读出，
  则存在唯一连续满同态
  \[
  F:K\twoheadrightarrow \Sigma_{S(Q)}
  \]
  使得全部 \(\rho_d^K\) 以及每个频率读出 \((\rho_{d_j}^K)^{m_j}\) 都由 \(\Sigma_{S(Q)}\) 上的标准字符经 \(F\) 拉回。换言之，一旦把 \(Q\) 的实现口径要求为对全部 \(S(Q)\)-光滑分母闭合的根一致读出，则 \(\Sigma_{S(Q)}\) 是普适最小载体。
  \item 对 \(\Sigma_{S(Q)}\) 上的角色族 \(\chi_{q_1},\dots,\chi_{q_r}\)，定义
  \[
  M_p:=\max_{1\le j\le r}\max\{0,-v_p(q_j)\}
  \qquad (p\in S(Q)).
  \]
  则有开子群包含
  \[
  \prod_{p\in S(Q)}p^{M_p}\ZZ_p
  \subseteq
  \bigcap_{j=1}^{r}\ker\bigl(\chi_{q_j}|_{K_{S(Q)}}\bigr),
  \]
  并且对每个素数 \(p\) 都有
  \[
  p\in S(Q)
  \iff
  M_p\ge 1
  \iff
  \bigcap_{j=1}^{r}\ker\bigl(\chi_{q_j}|_{K_{S(Q)}}\bigr)
  \text{ 含有非平凡开 }p\text{-进子群}.
  \]
  因而分母素数支撑可由共同 prime-register 核完全恢复。
\end{enumerate}
\end{theorem}
\begin{proof}
第 \(1\) 条是定理 \ref{thm:xi-time-part65g-solenoid-s-root-universal-single-valuedization} 在 \(S=S(Q)\) 时的直接应用；对每个 \(q_j=m_j/d_j\)，由该定理中的恒等式
\[
\chi_{q_j}=\rho_{d_j}^{\,m_j}
\]
可知 \(Q\) 的全部频率读出都被同一普适对象 \(\Sigma_{S(Q)}\) 吸收。

对第 \(2\) 条，把定理 \ref{thm:xi-time-part69-finite-character-prime-ledger-finite-depth} 应用于角色族
\[
\chi_{q_1},\dots,\chi_{q_r}\in \widehat{\Sigma_{S(Q)}}\cong \ZZ[S(Q)^{-1}].
\]
该定理给出开子群
\[
U(\chi_{q_1},\dots,\chi_{q_r})
=
\prod_{p\in S(Q)}p^{M_p}\ZZ_p
\]
并满足
\[
U(\chi_{q_1},\dots,\chi_{q_r})
\subseteq
\bigcap_{j=1}^{r}\ker\bigl(\chi_{q_j}|_{K_{S(Q)}}\bigr).
\]
另一方面，按 \(M_p\) 的定义，
\[
M_p\ge 1
\iff
\exists\,j\ \text{使}\ v_p(q_j)<0
\iff
p\ \text{出现在某个 }q_j\text{ 的既约分母中}
\iff
p\in S(Q).
\]
若 \(p\in S(Q)\)，则 \(p^{M_p}\ZZ_p\) 是 \(\ZZ_p\) 的非平凡开子群，因而共同核中含有非平凡开 \(p\)-进子群；反之若 \(p\notin S(Q)\)，则 \(K_{S(Q)}\) 本身并无 \(p\)-主方向，更不存在此类 \(p\)-进开子群。故所述等价成立。证毕。
\end{proof}

\begin{theorem}[Fold 多重度 Fourier 能量的三值相位游走表示]\label{thm:xi-time-part65g-fold-energy-threevalued-phase-walk}
\leanverified{paper\_xi\_time\_part65g\_fold\_energy\_threevalued\_phase\_walk}
沿用定理 \ref{thm:fold-spectrum-factorization} 的记号。对固定 \(m\ge 1\)，记
\[
M:=F_{m+2},
\qquad
\widehat d_m(k)=\prod_{j=1}^{m}\left(1+e^{-2\pi i kF_{j+1}/M}\right).
\]
令 \(\tau_1,\dots,\tau_m\) 独立同分布，并满足
\[
\PP(\tau_j=0)=\frac12,
\qquad
\PP(\tau_j=\pm 1)=\frac14.
\]
定义随机相位游走
\[
X_m:=
\sum_{j=1}^{m}\tau_j\frac{F_{j+1}}{F_{m+2}}
\in
\RR/\ZZ.
\]
则对任意整数 \(k\) 都有
\[
\EE\bigl[e^{2\pi i kX_m}\bigr]
=
\prod_{j=1}^{m}
\cos^2\!\Bigl(\pi k\frac{F_{j+1}}{F_{m+2}}\Bigr)
=
\frac{|\widehat d_m(k)|^2}{2^{2m}}.
\]
特别地，归一化 Fourier 能量剖面既是一个有限 Riesz 乘积，也是一个显式三值相位游走的特征函数。
\end{theorem}
\begin{proof}
由独立性，
\[
\EE\bigl[e^{2\pi i kX_m}\bigr]
=
\prod_{j=1}^{m}
\EE\!\left[
e^{2\pi i k\tau_jF_{j+1}/F_{m+2}}
\right].
\]
对任意实数 \(\theta\)，单步特征函数满足
\[
\EE[e^{i\theta\tau_j}]
=
\frac12+\frac14e^{i\theta}+\frac14e^{-i\theta}
=
\frac12(1+\cos\theta)
=
\cos^2(\theta/2).
\]
取
\[
\theta=2\pi k\frac{F_{j+1}}{F_{m+2}}
\]
便得到
\[
\EE\!\left[
e^{2\pi i k\tau_jF_{j+1}/F_{m+2}}
\right]
=
\cos^2\!\Bigl(\pi k\frac{F_{j+1}}{F_{m+2}}\Bigr).
\]
于是
\[
\EE\bigl[e^{2\pi i kX_m}\bigr]
=
\prod_{j=1}^{m}
\cos^2\!\Bigl(\pi k\frac{F_{j+1}}{F_{m+2}}\Bigr).
\]

另一方面，由定理 \ref{thm:fold-spectrum-factorization}，
\[
\frac{|\widehat d_m(k)|^2}{2^{2m}}
=
\prod_{j=1}^{m}
\frac{|1+e^{-2\pi i kF_{j+1}/F_{m+2}}|^2}{4}
=
\prod_{j=1}^{m}
\cos^2\!\Bigl(\pi k\frac{F_{j+1}}{F_{m+2}}\Bigr),
\]
因为
\[
|1+e^{-i\theta}|^2=4\cos^2(\theta/2).
\]
两式合并即得结论。证毕。
\end{proof}

\begin{theorem}[Dirichlet 频率上单一时间平移的实虚轴交换阻断]\label{thm:xi-time-part65g-dirichlet-axis-swap-obstruction}
设 \(A\subset\NN\) 为有限集，\(a_n\in\RR\) 为实系数，并定义 Dirichlet 多项式
\[
D(t):=\sum_{n\in A}a_n n^{-it}
=
\sum_{n\in A}a_n e^{-it\log n}.
\]
若存在 \(t_0\in\RR\) 使得对一切 \(t\in\RR\) 都有
\[
\bigl(\Re D(t+t_0),\Im D(t+t_0)\bigr)
=
\bigl(\Im D(t),\Re D(t)\bigr),
\]
则必有
\[
D\equiv 0.
\]
特别地，任何非零实系数 Dirichlet 多项式都不可能通过单一时间平移实现实轴与虚轴的全局交换。
\end{theorem}
\begin{proof}
记
\[
D(t)=x(t)+iy(t).
\]
题设等价于
\[
x(t+t_0)+iy(t+t_0)=y(t)+ix(t)=i\overline{D(t)}.
\]
由于所有 \(a_n\) 都是实数，
\[
\overline{D(t)}=D(-t),
\]
故得到函数恒等式
\[
D(t+t_0)=iD(-t)
\qquad (\forall t\in\RR).
\]
展开两边可得
\[
\sum_{n\in A}a_n e^{-it_0\log n}e^{-it\log n}
=
i\sum_{n\in A}a_n e^{it\log n}.
\]

指数函数族 \(\{e^{it\lambda}\}_{\lambda\in\Lambda}\) 在 \(\lambda\) 互异时于 \(\RR\) 上线性无关。对每个 \(n>1\)，左侧只出现频率 \(-\log n<0\)，右侧只出现频率 \(+\log n>0\)；因此这些频率的系数必须全部为零，从而
\[
a_n=0
\qquad (n>1).
\]
剩下唯一可能的频率是 \(\log 1=0\)。比较常数项得
\[
a_1=ia_1,
\]
故 \(a_1=0\)。于是全部系数皆为零，即 \(D\equiv 0\)。证毕。
\end{proof}

由此，局部化奇异环一侧的 \(S\)-光滑根提升、逆极限相位账本、prime-register 共同核恢复与 Fold 能量的概率化读出便形成了一条内部闭合的单值化链；与此同时，Dirichlet 频率侧试图以单一时间平移完成的全局换轴则被指数频率的线性无关性严格排除。于是“根单值化可行、相位账本可逆、干涉能量可概率化、单轴换位不可行”四个层面被压缩为同一条可审计的 singular-ring 主线。

\paragraph{压力缺口的乘法射线势、escort 单比特实验与 Hellinger--Chernoff 母核}
在整数压力轴的冻结斜率恢复、escort--Rényi 熵率的压力差公式、末位统计的渐近充分化、escort 熵常数的闭式化以及 Fisher--Rao 几何已经确立之后，还可继续压缩出下列六条彼此闭合的后果。它们分别把乘法细化中的压力损失写成一条精确射线势，把可除偏序上的斜率单调与整倍数冻结传播压缩为纯序结构命题，把有界温度 escort 实验族在 Blackwell 意义下指数塌缩为单比特 Bernoulli 家族，并把全阶 Rényi 常数层、Shannon 温度耗散以及二温度判别几何统一到同一个 logistic 母核之下。

\begin{theorem}[压力缺口是乘法熵亏损的精确射线势]\label{thm:xi-time-part65h-pressure-gap-ray-potential}
沿用定理 \ref{thm:xi-time-part65c-escort-renyi-two-parameter-freezing-line} 与定理 \ref{thm:xi-time-part57b-pressure-spectrum-slope-intercept-recovery} 的记号。设整数压力极限
\[
P_a:=\lim_{m\to\infty}\frac1m\log S_a(m)
\qquad (a\ge 1)
\]
都存在，并记
\[
\alpha_*:=\lim_{q\to\infty}\frac{P_q}{q}.
\]
对每个整数 \(a\ge 1\) 记 escort 分布
\[
\pi_{m,a}(x):=\frac{d_m(x)^a}{S_a(m)}.
\]
再对整数 \(a\ge 1\)、\(b\ge 2\) 定义
\[
\delta(a):=\frac{P_a}{a}-\alpha_*,
\qquad
H_{\mathrm{loss}}(a,b):=bP_a-P_{ab}.
\]
再记 escort--Rényi 熵率
\[
h_b^{\mathrm{esc}}(a):=
\lim_{m\to\infty}\frac1m H_b(\pi_{m,a}).
\]
其存在性由定理 \ref{thm:xi-time-part65c-escort-renyi-two-parameter-freezing-line} 保证。则对任意整数 \(a\ge 1\)、\(b\ge 2\) 都有严格恒等式
\[
\delta(a)-\delta(ab)=\frac{H_{\mathrm{loss}}(a,b)}{ab}.
\]
并且
\[
\delta(a)-\delta(ab)=\frac{b-1}{ab}\,h_b^{\mathrm{esc}}(a)\ge 0.
\]
更进一步，沿任意乘法射线都有绝对收敛展开
\[
\delta(a)=\sum_{k=0}^{\infty}\frac{H_{\mathrm{loss}}(ab^k,b)}{ab^{k+1}}.
\]
\leanverified{paper\_xi\_time\_part65h\_pressure\_gap\_ray\_potential}
\end{theorem}
\begin{proof}
第一式只是代数恒等式：
\[
\delta(a)-\delta(ab)
=
\left(\frac{P_a}{a}-\alpha_*\right)
-\left(\frac{P_{ab}}{ab}-\alpha_*\right)
=
\frac{bP_a-P_{ab}}{ab}
=
\frac{H_{\mathrm{loss}}(a,b)}{ab}.
\]

另一方面，由定理 \ref{thm:xi-time-part65c-escort-renyi-two-parameter-freezing-line}，
\[
h_b^{\mathrm{esc}}(a)
=
\frac{bP_a-P_{ab}}{b-1}
=
\frac{H_{\mathrm{loss}}(a,b)}{b-1}.
\]
代回即得
\[
\delta(a)-\delta(ab)=\frac{b-1}{ab}\,h_b^{\mathrm{esc}}(a)\ge 0,
\]
其中非负性来自 Rényi 熵率本身非负。

最后，对任意 \(N\ge 1\) 作有限望远镜求和：
\[
\delta(a)-\delta(ab^N)
=
\sum_{k=0}^{N-1}\bigl(\delta(ab^k)-\delta(ab^{k+1})\bigr)
=
\sum_{k=0}^{N-1}\frac{H_{\mathrm{loss}}(ab^k,b)}{ab^{k+1}}.
\]
由定理 \ref{thm:xi-time-part57b-pressure-spectrum-slope-intercept-recovery}，
\[
\frac{P_q}{q}\longrightarrow \alpha_*
\qquad (q\to\infty),
\]
故当 \(N\to\infty\) 时，
\[
\delta(ab^N)=\frac{P_{ab^N}}{ab^N}-\alpha_*\longrightarrow 0.
\]
于是得到
\[
\delta(a)=\sum_{k=0}^{\infty}\frac{H_{\mathrm{loss}}(ab^k,b)}{ab^{k+1}}.
\]
又由于各项都可写成
\[
\frac{H_{\mathrm{loss}}(ab^k,b)}{ab^{k+1}}
=
\frac{b-1}{ab^{k+1}}\,h_b^{\mathrm{esc}}(ab^k)\ge 0,
\]
该级数因而绝对收敛。证毕。
\end{proof}

\begin{corollary}[可除偏序上的斜率单调、几何逼近与整倍数冻结传播]\label{cor:xi-time-part65h-divisibility-slope-freezing-propagation}
沿用定理 \ref{thm:xi-time-part65h-pressure-gap-ray-potential} 的记号。则成立：
\begin{enumerate}
  \item 若 \(a\mid c\)，则
  \[
  \frac{P_c}{c}\le \frac{P_a}{a}.
  \]
  \item 对任意整数 \(a\ge 1\)、\(b\ge 2\)、\(N\ge 0\)，都有
  \[
  0\le \frac{P_{ab^N}}{ab^N}-\alpha_*
  \le \frac{\log\varphi}{ab^N}.
  \]
  \item 若对某个整数 \(a_0\ge 1\) 恰有
  \[
  P_{a_0}=a_0\alpha_*,
  \]
  则对每个 \(n\ge 1\) 都有
  \[
  P_{a_0n}=a_0n\,\alpha_*.
  \]
\end{enumerate}
\end{corollary}
\begin{proof}
若 \(c=a\)，结论平凡。若 \(c>a\) 且 \(a\mid c\)，写 \(c=ab\) 且 \(b\ge 2\)。由定理 \ref{thm:xi-time-part65h-pressure-gap-ray-potential}，
\[
\frac{P_a}{a}-\frac{P_c}{c}
=
\delta(a)-\delta(ab)
=
\frac{b-1}{ab}\,h_b^{\mathrm{esc}}(a)\ge 0,
\]
故第 \(1\) 条成立。

对第 \(2\) 条，令 \(n:=ab^N\)。由定义，
\[
M_m^n\le S_n(m)\le |X_m|\,M_m^n.
\]
两边取对数后除以 \(mn\)，再令 \(m\to\infty\)，得到
\[
\alpha_*\le \frac{P_n}{n}\le \alpha_*+\frac{1}{n}\lim_{m\to\infty}\frac1m\log|X_m|.
\]
而 \(|X_m|=F_{m+2}\)，故
\[
\lim_{m\to\infty}\frac1m\log|X_m|=\log\varphi.
\]
于是
\[
0\le \frac{P_n}{n}-\alpha_*\le \frac{\log\varphi}{n}.
\]
代回 \(n=ab^N\) 即得所示估计。

最后，若 \(P_{a_0}=a_0\alpha_*\)，则对任意 \(n\ge 1\)，由第 \(1\) 条与第 \(2\) 条分别得到
\[
\alpha_*\le \frac{P_{a_0n}}{a_0n}\le \frac{P_{a_0}}{a_0}=\alpha_*.
\]
故必有
\[
\frac{P_{a_0n}}{a_0n}=\alpha_*,
\]
即
\[
P_{a_0n}=a_0n\,\alpha_*.
\]
证毕。
\end{proof}

\begin{theorem}[有界温度 escort 实验族与单比特 Bernoulli 实验的 Blackwell 指数等价]\label{thm:xi-time-part65h-escort-blackwell-onebit-equivalence}
沿用定理 \ref{thm:killo-fold-bin-escort-onebit-fisher} 的记号。固定 \(Q<\infty\)，定义统计实验族
\[
\mathcal E_m^{(Q)}:=\{\pi_{m,q}:0\le q\le Q\},
\]
以及极限单比特实验族
\[
\mathcal B^{(Q)}:=\{\beta_q:0\le q\le Q\},
\qquad
\beta_q:=\mathrm{Bernoulli}(\theta(q)),
\qquad
\theta(q):=\frac{1}{1+\varphi^{q+1}}.
\]
再定义
\[
\Delta_Q(m):=
\inf_{K,L}\sup_{0\le q\le Q}
\max\Bigl\{
D_{\mathrm{TV}}(K_{\#}\pi_{m,q},\beta_q),\
D_{\mathrm{TV}}(L_{\#}\beta_q,\pi_{m,q})
\Bigr\},
\]
其中 \(K:X_m\rightsquigarrow\{0,1\}\) 与 \(L:\{0,1\}\rightsquigarrow X_m\) 取遍全部 Markov 核。则存在常数 \(C_Q>0\) 使得
\[
\Delta_Q(m)\le C_Q\Bigl(\frac{\varphi}{2}\Bigr)^m.
\]
因此，对整个有界温度区间而言，escort 统计实验在 Blackwell 意义下以指数速度塌缩为由末位比特承载的一维 logistic Bernoulli 家族。
\end{theorem}
\begin{proof}
记
\[
A_i:=\{w\in X_m:w_m=i\},
\qquad
U_{m,i}:=\mathrm{Unif}(A_i)
\qquad (i=0,1).
\]
定义前向核
\[
K_m(w):=w_m\in\{0,1\},
\]
以及反向核
\[
L_m(i,\cdot):=U_{m,i}.
\]
由定理 \ref{thm:killo-fold-bin-escort-onebit-fisher}，存在两态 Gibbs 近似
\[
\nu_{m,q}(w):=\frac{u_m(w)\varphi^{-qw_m}}{Z_{m,q}},
\qquad
Z_{m,q}:=\frac{F_{m+1}+\varphi^{-q}F_m}{F_{m+2}},
\]
并且
\[
\sup_{0\le q\le Q}\|\pi_{m,q}-\nu_{m,q}\|_{\mathrm{TV}}
=
O_Q\!\Bigl(\Bigl(\frac{\varphi}{2}\Bigr)^m\Bigr).
\]
由于 \(\nu_{m,q}\) 在每个块 \(A_i\) 上严格均匀，可把它写成
\[
\nu_{m,q}=(1-\theta_m(q))U_{m,0}+\theta_m(q)U_{m,1},
\]
其中
\[
\theta_m(q):=\nu_{m,q}(A_1)
=
\frac{\varphi^{-q}F_m}{F_{m+1}+\varphi^{-q}F_m}.
\]

先看前向方向。由 \(K_m\) 的定义，
\[
K_{m\#}\pi_{m,q}=\mathrm{Bernoulli}(\theta_m(q)).
\]
而
\[
\theta_m(q)=\frac{\varphi^{-q}}{F_{m+1}/F_m+\varphi^{-q}},
\qquad
\theta(q)=\frac{\varphi^{-q}}{\varphi+\varphi^{-q}}.
\]
又因为
\[
\frac{F_{m+1}}{F_m}=\varphi+O(\varphi^{-2m}),
\]
并且 \(q\in[0,Q]\) 落在紧区间上，所以
\[
\sup_{0\le q\le Q}|\theta_m(q)-\theta(q)|
=
O_Q(\varphi^{-2m})
=
O_Q\!\Bigl(\Bigl(\frac{\varphi}{2}\Bigr)^m\Bigr).
\]
因此
\[
\sup_{0\le q\le Q}
D_{\mathrm{TV}}(K_{m\#}\pi_{m,q},\beta_q)
=
O_Q\!\Bigl(\Bigl(\frac{\varphi}{2}\Bigr)^m\Bigr).
\]

再看反向方向。由构造，
\[
L_{m\#}\beta_q=(1-\theta(q))U_{m,0}+\theta(q)U_{m,1}.
\]
而与 \(\nu_{m,q}\) 相比，两者仅块权重不同，故
\[
\|L_{m\#}\beta_q-\nu_{m,q}\|_{\mathrm{TV}}
=
|\theta(q)-\theta_m(q)|
=
O_Q\!\Bigl(\Bigl(\frac{\varphi}{2}\Bigr)^m\Bigr).
\]
再与 \(\pi_{m,q}\) 到 \(\nu_{m,q}\) 的估计合并，得到
\[
\sup_{0\le q\le Q}
D_{\mathrm{TV}}(L_{m\#}\beta_q,\pi_{m,q})
=
O_Q\!\Bigl(\Bigl(\frac{\varphi}{2}\Bigr)^m\Bigr).
\]
对 \(\Delta_Q(m)\) 取上述这对核 \((K_m,L_m)\) 作为候选，即得所示上界。证毕。
\end{proof}

\begin{theorem}[bin-fold escort 的全阶 Rényi 熵在常数层具有显式温度剖面]\label{thm:xi-time-part65h-escort-renyi-o1-temperature-profile}
沿用定理 \ref{thm:killo-fold-bin-escort-onebit-fisher} 的记号。对固定 \(q\ge 0\) 与 \(\alpha>0\)，定义 escort Rényi 熵
\[
H_{\alpha}(\pi_{m,q})
:=
\begin{cases}
\dfrac{1}{1-\alpha}\log\sum_{w\in X_m}\pi_{m,q}(w)^\alpha,
& \alpha\neq 1,\\[2mm]
-\sum_{w\in X_m}\pi_{m,q}(w)\log\pi_{m,q}(w),
& \alpha=1.
\end{cases}
\]
则当 \(m\to\infty\) 时，若 \(\alpha\neq 1\)，有
\[
H_{\alpha}(\pi_{m,q})
=
m\log\varphi
-\frac12\log 5
+
\frac{\alpha\log(\varphi+\varphi^{-q})-\log(\varphi+\varphi^{-\alpha q})}{\alpha-1}
+O\!\Bigl(\Bigl(\frac{\varphi}{2}\Bigr)^m\Bigr),
\]
而当 \(\alpha=1\) 时，有
\[
H(\pi_{m,q})
=
m\log\varphi
-\frac12\log 5
+\log(\varphi+\varphi^{-q})
+\frac{q\log\varphi}{1+\varphi^{q+1}}
+O\!\Bigl(\Bigl(\frac{\varphi}{2}\Bigr)^m\Bigr).
\]
因此，escort 温度族不只在熵率层面塌缩到 \(\log\varphi\)；其全部 \(O(1)\) 剩余结构也由 \(q\) 的显式闭式函数完全控制。
\end{theorem}
\begin{proof}
先设 \(\alpha\neq 1\)。由定义，
\[
\sum_{w\in X_m}\pi_{m,q}(w)^\alpha
=
\frac{S_{\alpha q}^{\mathrm{bin}}(m)}{(S_q^{\mathrm{bin}}(m))^\alpha}.
\]
而命题 \ref{prop:fold-bin-power-sum-asymptotic} 给出，对每个固定 \(t\ge 0\)，
\[
S_t^{\mathrm{bin}}(m)
=
\Bigl(\frac{2}{\varphi}\Bigr)^{tm}
\bigl(F_{m+1}+\varphi^{-t}F_m\bigr)
\Bigl(1+O\!\Bigl(\Bigl(\frac{\varphi}{2}\Bigr)^m\Bigr)\Bigr).
\]
将 \(t=q\) 与 \(t=\alpha q\) 代入，即得
\[
\log\sum_{w\in X_m}\pi_{m,q}(w)^\alpha
=
\log\bigl(F_{m+1}+\varphi^{-\alpha q}F_m\bigr)
-\alpha\log\bigl(F_{m+1}+\varphi^{-q}F_m\bigr)
+O\!\Bigl(\Bigl(\frac{\varphi}{2}\Bigr)^m\Bigr).
\]
再由 Binet 公式，
\[
F_{m+1}+\varphi^{-u}F_m
=
\frac{\varphi^m}{\sqrt5}\bigl(\varphi+\varphi^{-u}\bigr)+O(\varphi^{-m}),
\]
故
\[
\log\bigl(F_{m+1}+\varphi^{-u}F_m\bigr)
=
m\log\varphi-\frac12\log5+\log(\varphi+\varphi^{-u})
+O(\varphi^{-2m}).
\]
代回并除以 \(1-\alpha\)，便得到所示 Rényi 公式。

再看 Shannon 情形。由定理 \ref{thm:xi-fold-escort-renyi-shannon-constants} 的证明所用恒等式，
\[
H(\pi_{m,q})
=
\log S_q^{\mathrm{bin}}(m)-q\,\kappa_m^{\mathrm{bin}}(q),
\]
其中
\[
\kappa_m^{\mathrm{bin}}(q):=
\sum_{w\in X_m}\pi_{m,q}(w)\log d_m^{\mathrm{bin}}(w).
\]
由推论 \ref{cor:xi-time-part9o-escort-fiberinfo-constantterm-inversion}，
\[
\kappa_m^{\mathrm{bin}}(q)
=
m\log\!\Bigl(\frac{2}{\varphi}\Bigr)
-\frac{\log\varphi}{1+\varphi^{q+1}}
+O\!\Bigl(\Bigl(\frac{\varphi}{2}\Bigr)^m\Bigr).
\]
另一方面，把 \(t=q\) 代入上述幂和展开并取对数，得到
\[
\log S_q^{\mathrm{bin}}(m)
=
qm\log\!\Bigl(\frac{2}{\varphi}\Bigr)
+m\log\varphi
-\frac12\log5
+\log(\varphi+\varphi^{-q})
+O\!\Bigl(\Bigl(\frac{\varphi}{2}\Bigr)^m\Bigr).
\]
两式合并即得
\[
H(\pi_{m,q})
=
m\log\varphi
-\frac12\log5
+\log(\varphi+\varphi^{-q})
+\frac{q\log\varphi}{1+\varphi^{q+1}}
+O\!\Bigl(\Bigl(\frac{\varphi}{2}\Bigr)^m\Bigr).
\]
证毕。
\end{proof}

\begin{theorem}[escort Shannon 常数满足精确 de Bruijn--Fisher 身份]\label{thm:xi-time-part65h-escort-shannon-debruijn-fisher}
沿用定理 \ref{thm:xi-time-part65h-escort-renyi-o1-temperature-profile} 的记号。定义 Shannon 常数层
\[
\mathfrak h(q):=
\lim_{m\to\infty}\bigl(H(\pi_{m,q})-m\log\varphi\bigr).
\]
则有闭式
\[
\mathfrak h(q)
=
-\frac12\log5+\log(\varphi+\varphi^{-q})
+\frac{q\log\varphi}{1+\varphi^{q+1}}.
\]
进一步，记
\[
\theta(q):=\frac{1}{1+\varphi^{q+1}},
\qquad
\mathcal I(q):=(\log\varphi)^2\theta(q)\bigl(1-\theta(q)\bigr).
\]
则成立精确微分恒等式
\[
\mathfrak h'(q)=-q\,\mathcal I(q),
\]
以及积分表示
\[
\mathfrak h(q)=\mathfrak h(0)-\int_0^q t\,\mathcal I(t)\,\dd t.
\]
其中 \(\mathcal I(q)\) 正是结论 \ref{con:xi-time-part9o-escort-fisher-rao-closed} 给出的极限 Fisher 信息密度。
\end{theorem}
\begin{proof}
第一条公式正是定理 \ref{thm:xi-time-part65h-escort-renyi-o1-temperature-profile} 在 \(\alpha=1\) 时的常数项。

令
\[
A(q):=\varphi+\varphi^{-q}.
\]
则
\[
\theta(q)=\frac{\varphi^{-q}}{A(q)}=\frac{1}{1+\varphi^{q+1}},
\]
并且
\[
\mathfrak h(q)
=
-\frac12\log5+\log A(q)+q\,\theta(q)\log\varphi.
\]
直接求导得
\[
\frac{A'(q)}{A(q)}
=
-\frac{\varphi^{-q}\log\varphi}{A(q)}
=
-\theta(q)\log\varphi,
\]
而
\[
\theta'(q)
=
-(\log\varphi)\theta(q)\bigl(1-\theta(q)\bigr).
\]
因此
\[
\mathfrak h'(q)
=
-\theta(q)\log\varphi+\theta(q)\log\varphi+q\,\theta'(q)\log\varphi
=
-q(\log\varphi)^2\theta(q)\bigl(1-\theta(q)\bigr)
=
-q\,\mathcal I(q).
\]
最后对 \(q\) 从 \(0\) 到任意给定值积分，即得
\[
\mathfrak h(q)=\mathfrak h(0)-\int_0^q t\,\mathcal I(t)\,\dd t.
\]
证毕。
\end{proof}

\begin{theorem}[escort 温度间的 Hellinger--Chernoff 几何由闭式母核完全支配]\label{thm:xi-time-part65h-escort-hellinger-chernoff-kernel}
沿用定理 \ref{thm:killo-fold-bin-escort-onebit-fisher} 的记号。对任意 \(p,q\ge 0\) 与 \(s\in[0,1]\)，定义混合亲和量
\[
\mathcal M_s^{(m)}(p,q):=
\sum_{w\in X_m}\pi_{m,p}(w)^s\pi_{m,q}(w)^{1-s}.
\]
则极限
\[
\mathcal M_s(p,q):=\lim_{m\to\infty}\mathcal M_s^{(m)}(p,q)
\]
存在，并满足精确闭式
\[
\mathcal M_s(p,q)
=
\frac{\varphi+\varphi^{-(sp+(1-s)q)}}
{(\varphi+\varphi^{-p})^s(\varphi+\varphi^{-q})^{1-s}}.
\]
特别地，Hellinger 亲和量
\[
\mathcal H(p,q):=\mathcal M_{1/2}(p,q)
\]
满足
\[
\mathcal H(p,q)
=
\frac{\varphi+\varphi^{-(p+q)/2}}
{\sqrt{(\varphi+\varphi^{-p})(\varphi+\varphi^{-q})}},
\]
而 Chernoff 母核可写为
\[
\mathcal C(p,q):=\inf_{0\le s\le 1}\mathcal M_s(p,q).
\]
因此，两温度 escort 之间的全部 Hellinger--Chernoff 判别几何都由单参数显式核
\[
s\longmapsto \mathcal M_s(p,q)
\]
完全控制。
\end{theorem}
\begin{proof}
对任意 \(s\in[0,1]\)，由 escort 分布的定义，
\[
\pi_{m,p}(w)^s\pi_{m,q}(w)^{1-s}
=
\frac{d_m^{\mathrm{bin}}(w)^{sp+(1-s)q}}
{\bigl(S_p^{\mathrm{bin}}(m)\bigr)^s\bigl(S_q^{\mathrm{bin}}(m)\bigr)^{1-s}}.
\]
故有严格恒等式
\[
\mathcal M_s^{(m)}(p,q)
=
\frac{S_{sp+(1-s)q}^{\mathrm{bin}}(m)}
{\bigl(S_p^{\mathrm{bin}}(m)\bigr)^s\bigl(S_q^{\mathrm{bin}}(m)\bigr)^{1-s}}.
\]
对分子与分母分别应用命题 \ref{prop:fold-bin-power-sum-asymptotic}，
\[
S_t^{\mathrm{bin}}(m)
=
\Bigl(\frac{2}{\varphi}\Bigr)^{tm}
\bigl(F_{m+1}+\varphi^{-t}F_m\bigr)
\Bigl(1+O\!\Bigl(\Bigl(\frac{\varphi}{2}\Bigr)^m\Bigr)\Bigr)
\qquad (t\ge 0).
\]
其中指数因子完全抵消。
于是
\[
\mathcal M_s^{(m)}(p,q)
=
\frac{F_{m+1}+\varphi^{-(sp+(1-s)q)}F_m}
{\bigl(F_{m+1}+\varphi^{-p}F_m\bigr)^s
\bigl(F_{m+1}+\varphi^{-q}F_m\bigr)^{1-s}}
\Bigl(1+O\!\Bigl(\Bigl(\frac{\varphi}{2}\Bigr)^m\Bigr)\Bigr).
\]
再由 Binet 公式
\[
F_{m+1}+\varphi^{-u}F_m
=
\frac{\varphi^m}{\sqrt5}\bigl(\varphi+\varphi^{-u}\bigr)+O(\varphi^{-m}),
\]
便得到
\[
\mathcal M_s^{(m)}(p,q)
=
\frac{\varphi+\varphi^{-(sp+(1-s)q)}}
{(\varphi+\varphi^{-p})^s(\varphi+\varphi^{-q})^{1-s}}
+O\!\Bigl(\Bigl(\frac{\varphi}{2}\Bigr)^m\Bigr).
\]
令 \(m\to\infty\) 即得所示极限核。

取 \(s=\tfrac12\) 立即得到 Hellinger 亲和量公式；而 Chernoff 母核不过是同一闭式核在区间 \([0,1]\) 上的下确界。证毕。
\end{proof}

由此，escort 主线中的整数压力差、末位单比特充分化、Rényi 常数层、Fisher 温度耗散与 Hellinger--Chernoff 判别核便被进一步压缩到同一条闭合链上：乘法偏序一侧的全部亏损由射线势 \(\delta\) 精确编码，而统计实验一侧的全部有限温度几何则完全退化为由 \(\theta(q)=1/(1+\varphi^{q+1})\) 所驱动的一维 logistic 母模型。


\endinput
